\documentclass[12pt]{article}
\usepackage[T1]{fontenc}
\usepackage{lmodern,amsmath,amssymb,amsthm,booktabs,longtable,microtype}
\usepackage[a4paper,margin=24mm]{geometry}
\usepackage{xcolor}
\definecolor{paperlink}{RGB}{0,0,114}
\usepackage[colorlinks=true,linkcolor=paperlink,citecolor=paperlink,urlcolor=paperlink]{hyperref}
\setlength{\emergencystretch}{2em}
\newtheorem{theorem}{Theorem}[section]
\newtheorem{proposition}[theorem]{Proposition}
\newtheorem{lemma}[theorem]{Lemma}
\newtheorem{corollary}[theorem]{Corollary}
\newtheorem{remark}[theorem]{Remark}
\newcommand{\OK}{\mathcal O_K}
\numberwithin{equation}{section}
\title{A uniform bound of five for Diophantine tuples\\in imaginary quadratic integer rings}
\author{Zrinka Franu\v{s}i\'{c} \and Tomislav Pejkovi\'{c}}
\date{}
\begin{document}
\maketitle
\begin{abstract}
We prove that every $D(1)$-tuple in the full ring of integers of an
imaginary quadratic field has at most five elements. The proof combines
determinant estimates, interpolation of square roots, and explicit
simultaneous approximation. The remaining cases are reduced to finite
computations. We give explicit search bounds and enumeration procedures,
together with tables of the resulting counts and cutoff constants.
\end{abstract}

\section{Introduction and conventions}

Let $K$ be an imaginary quadratic field with full ring of integers
$\mathcal O_K$. A finite set $\mathcal A\subset\mathcal O_K\setminus\{0\}$
is a $D(1)$-tuple if $xy+1$ is a square in $\mathcal O_K$ whenever
$x,y\in\mathcal A$ are distinct. Zero is allowed as such a square.
Equivalently, the square roots may be required to lie in $K$. A square
root in $K$ of an algebraic integer is again in $\mathcal O_K$.

Ad\v zaga obtained the uniform bound $42$ for these tuples~\cite{A}.
Cipu and Fujita proved the bound $24$ for both $D(1)$-tuples and
$D(-1)$-tuples in imaginary quadratic integer rings~\cite{CF}.
We obtain the following bound for $D(1)$-tuples.

\begin{theorem}\label{thm:main}
For every imaginary quadratic field $K$, every $D(1)$-tuple in
$\mathcal O_K$ has cardinality at most five.
\end{theorem}

Call an entry small if its modulus is less than three. We order a
hypothetical sextuple by modulus and exclude the following three cases.
\begin{center}
\begin{tabular}{ccl}
Number of small entries&Sections&Closing result\\\hline
$0$&\ref{sec:all-large}&Proposition~\ref{prop:no-all-large-six}\\
$1$&\ref{sec:pair-reduction}--\ref{sec:one-small-completion}&Proposition~\ref{prop:one-small}\\
at least $2$&\ref{sec:two-small}&Proposition~\ref{sm:main}
\end{tabular}
\end{center}
Determinant estimates and interpolation give lower bounds for the larger
entries. We then construct integral linear forms in square roots that
exclude overlapping ranges for the largest entry. Together these
estimates settle the case with no small entries and reduce the other
cases to bounded computations. For two small entries, a descent generates
all common extensions below the required bound.

The simultaneous approximation construction follows Bennett~\cite{Bennett96}
and Jadrijevi\'c and Ziegler~\cite{JZ}. We give the polynomial formulas and
the coefficient and remainder estimates needed here. The determinant
identity of Bennett is stated explicitly. The nonvanishing argument for
a $3\times3$ matrix of roots of $a_i b_j+1$ also appears in
Wildner Fraga's preprint~\cite[Lemma~3.1]{W}. We include that argument, without using
the quantitative Subspace Theorem considered there. The factorization
used for the pair $\{-1,1\}$ also appears in Cipu's lecture
slides~\cite[pp.~27--30]{Cipu25}. Franu\v{s}i\'c~\cite[Theorem~1]{F10} proved that
$\{1,3\}$ cannot be extended to a Diophantine quintuple in
$\mathbb Z[\sqrt{-d}]$ for positive squarefree $d\ne2$.
Franu\v{s}i\'c and Kreso~\cite[Corollary~1.2]{FK11} treated $d=2$.
For $d\equiv3\pmod4$, the ring $\mathbb Z[\sqrt{-d}]$ is properly contained
in the full ring of integers of $\mathbb Q(\sqrt{-d})$. We handle the pair in the full ring of integers using
the same approximation estimates as for the other pairs.

Dujella~\cite[Problem~1.11]{DujellaOpen} asks whether a $D(1)$-quintuple
exists in the ring of integers of an imaginary quadratic field.
Our bound does not settle this question.
For a general account of Diophantine tuples and their connections with
elliptic curves, see Dujella~\cite{Dujella24}.

Fix one complex embedding of $K$. We write $N(z)=z\overline z=|z|^2$.
For $z\in\mathcal O_K\setminus\{0\}$ this is a positive integer, so
$|z|\ge1$. When lower-case letters denote ordered tuple entries,
the corresponding capitals denote their moduli. Equal moduli are allowed throughout. Roots of the elements $xy+1$
supplied by the tuple are chosen in $\mathcal O_K$. Auxiliary square
roots used in analytic estimates may lie in $\mathbb C$.

For the exact enumerations, write $K=\mathbb Q(\sqrt{-d})$, with
$d>0$ squarefree, and
\[
 \omega=\begin{cases}(1+\sqrt{-d})/2,&d\equiv3\pmod4,\\
                     \sqrt{-d},&d\not\equiv3\pmod4.
       \end{cases}
\]
Then $\mathcal O_K=\mathbb Z[\omega]$ and
$\omega^2-\tau_K\omega+\nu_K=0$, where
$(\tau_K,\nu_K)=(1,(d+1)/4)$ or $(0,d)$, respectively.
The absolute field discriminant is $\Delta_K=4\nu_K-\tau_K^2$.
Tables using the coordinates $(u+v\sqrt{-d})/2$ are identified explicitly.
All other coordinate pairs denote $u+v\omega$. Both descriptions include
the full ring of integers.

Simultaneous negation and complex conjugation preserve the tuple
property. These, together with permutation, are the only equivalences
used in the finite reductions. Multiplication by an arbitrary unit
is not used as a tuple symmetry.

All numerical verification below reduces to integer or rational
arithmetic. Where an irrational modulus occurs, its square is an
explicit integer or rational number and the chosen enclosure is
certified by squaring. The finite-sum formulas, bounded enumeration
procedures, and tables of counts and cutoff constants are included in the article.

\section{Structural growth estimates}

A triple is \emph{regular} if
\[
 a^2+b^2+c^2-2ab-2ac-2bc=4,
\]
equivalently \(c=a+b\pm2\sqrt{ab+1}\). This condition is symmetric.

\begin{lemma}\label{lem:irregular-growth}
An irregular triple with \(3\le A=|a|\le B=|b|\le C=|c|\) satisfies
\(C>2AB\).
\end{lemma}
\begin{proof}
Choose \(r^2=ab+1,s^2=ac+1,t^2=bc+1\) in \(\mathcal O_K\), and put
\(u,v=a+b+c+2abc\pm2rst\). Their product is
\(a^2+b^2+c^2-2ab-2ac-2bc-4\), so irregularity makes both nonzero.
Since \(u,v\) are nonzero algebraic integers,
\(|u|+|v|\le|uv|+1\). Hence
\begin{equation}\label{eq:irregular-quadratic}
 0\le C^2-2\bigl(|1+2ab|-|a+b|\bigr)C
      +|(a-b)^2-4|+2|a+b|+1.
\end{equation}
Using \(|1+2ab|\ge2AB-1\), \(|a+b|\le A+B\), and
\(|(a-b)^2-4|\le(A+B)^2+4\), we obtain \(q(C)\ge0\),
where \(q(X)=(A+B+X+1)^2+4-4ABX\).
But
\[
 \frac{q(B)}{B^2}\le\left(\frac{10}{3}\right)^2+\frac49-12=-\frac49.
\]
Writing \(h=(A+B+1)/(AB)\le7/9\), we also have
\[
 \frac{q(2AB)}{4A^2B^2}
 =-1+h+\frac{h^2}{4}+\frac1{A^2B^2}\le-\frac{19}{324}.
\]
Convexity excludes \(B\le C\le2AB\).
\end{proof}

\begin{lemma}\label{lem:tail-irregular}
If an ordered quadruple has \(1\le A\le B\le C\le D\) and \(B\ge3\),
its last three entries form an irregular triple. In particular \(D>2BC\).
\end{lemma}
\begin{proof}
Suppose instead \(d=b+c+2r\), \(r^2=bc+1\). First let \(C\ge4\).
Choose \(\beta^2=b,\gamma^2=c\) so that
\(r=\beta\gamma\sqrt{1+1/(bc)}\), where \(\sqrt{1+z}\) is the
analytic branch with value \(1\) at \(z=0\).
Since \(BC\ge12\),
\[
 |r-\beta\gamma|<\frac{20}{39\sqrt{BC}}.
\]
Put \(p=\beta+\gamma\). The ordering \(D\ge C\) gives
\[
 |p|^2>C-\frac{40}{39\sqrt{BC}}>
 \frac{12}{13}C>\left(\frac{24}{25}\right)^2C.
\]
Indeed \(\sqrt{BC}>24/7\), so the relative loss is below
\(35/468<1/13\).
Choose \(\alpha^2=a\), and select the integral roots \(u^2=ab+1\),
\(v^2=ac+1\) using the same analytic branch. Then
\[
 |u-\alpha\beta|<\frac5{9\sqrt{AB}},\qquad
 |v-\alpha\gamma|<\frac7{13\sqrt{AC}}.
\]
For \(\delta=ad+1-ap^2=1+2a(r-\beta\gamma)\),
\[
 \left|\frac{\delta}{ap^2}\right|
 <\frac{1+35/117}{48/13}=\frac{19}{54}<\frac9{25}.
\]
Consequently the integral root
\(w=\alpha p\sqrt{1+\delta/(ap^2)}\) satisfies
\[
 |w-\alpha p|<
 \frac{125}{216}\left(\frac1{\sqrt{AC}}
       +\frac{40\sqrt A}{39\sqrt B\,C}\right).
\]
Write \(q_0=7/13+125/216\), \(r_0=625/1053\). It follows that
\[
 |u+v-w|<
 \frac5{9\sqrt{AB}}+\frac{q_0}{\sqrt{AC}}
                    +\frac{r_0\sqrt A}{\sqrt B\,C}<1.
\]
For completeness, the last expression is convex in \(\sqrt A\)
on \([1,\sqrt B]\). At the two endpoints, using \(B\ge3,C\ge4\),
its upper bounds are respectively
\[
 \frac{35}{108}+\frac{q_0}{2}+\frac{7r_0}{48}
   =\frac{12247}{12636}<1,\qquad
 \frac5{27}+\frac{7q_0}{24}+\frac{r_0}{4}
   =\frac{14813}{22464}<1.
\]
Integrality gives \(w=u+v\). Squaring and using
\(d=b+c+2r\) yields \(2(ar-uv)=1\), impossible because
\(ar-uv\) is an algebraic integer.

It remains to handle \(3\le B\le C<4\). The following finite computation
is included explicitly. Write
\(\omega^2=\tau_K\omega-\nu_K\) for the usual integral basis of
\(K=\mathbb Q(\sqrt{-q})\), with \(\Delta_K=4\nu_K-\tau_K^2\). Then
\begin{equation}\label{eq:integrated-disk}
\begin{aligned}
 N(x+y\omega)&=x^2+\tau_Kxy+\nu_Ky^2,\\
 4N(x+y\omega)&=(2x+\tau_Ky)^2+\Delta_Ky^2.
\end{aligned}
\end{equation}
If either \(b,c\) is nonreal, \(\Delta_K\le60\). If both are rational,
they must be \(-3,3\), and the square condition restricts the field to
\(\mathbb Q(\sqrt{-2})\). It suffices to enumerate the following squarefree values of $q$ with $\Delta_K\le60$.
\[
 1,2,3,5,6,7,10,11,13,14,15,19,23,31,35,39,43,47,51,55,59.
\]
For each field enumerate, by \eqref{eq:integrated-disk}, all elements of norm at most 15.
For distinct \(b,c\) with \(9\le N(b)\le N(c)\le15\), retain
\(bc+1=r^2\), form \(d=b+c\pm2r\), and retain
\(d\notin\{0,b,c\}\), \(N(d)\ge N(c)\). Finally enumerate
\(a\notin\{0,b,c,d\}\), \(N(a)\le N(b)\), testing \(ab+1,ac+1,ad+1\).
Every square test uses an exact finite table. Since
\(|d|<8+2\sqrt{17}<17\), all roots tested have norm at most 70. Enumerating those roots by \eqref{eq:integrated-disk} and squaring
them gives the complete square table, including zero.

Summing over these fields gives 90 admissible unordered pairs \(b,c\),
118 retained triples, and 16 candidates before the final square test. Simultaneous negation,
conjugation, and interchange of \(b,c\) reduce these to the six rows below.
Coordinates mean \(x+y\omega\). Here \(\omega=\sqrt{-q}\) unless
\(q\equiv3\pmod4\), when \(\omega=(1+\sqrt{-q})/2\).
\[
\begin{array}{c|rrrrr}
q&a&b&c&d&ad+1\\\hline
2&(-1,0)&(-3,0)&(3,0)&(0,-4)&(1,4)\\
2&(-1,0)&(2,-2)&(3,0)&(11,-4)&(-10,4)\\
3&(-1,2)&(-1,-3)&(4,-3)&(7,-14)&(22,0)\\
6&(-2,0)&(3,-1)&(3,1)&(14,0)&(-27,0)\\
11&(-1,0)&(3,1)&(4,-1)&(15,0)&(-14,0)\\
35&(-3,0)&(2,1)&(3,-1)&(13,0)&(-38,0)
\end{array}
\]
The first two final values have nonsquare norms 33 and 132. The third is
the positive rational nonsquare 22. For the last three, the ratios to
\(-q\) are \(9/2,14/11,38/35\), none a rational square. Thus no row
passes the last test. This proves irregularity. The growth assertion
follows from Lemma~\ref{lem:irregular-growth}.
\end{proof}

\subsection{The determinant estimates}
For entries \(a_1,\ldots,a_n\), choose \(u_i^2=a_i\) in \(\mathbb C\)
and signs of the integral roots so that
\[
 M_{ii}=a_i,\quad
 M_{ij}=u_i u_j\sqrt{1+1/(a_i a_j)}\quad(i\ne j).
\]
All products below have modulus at least 3, so these branches are defined.
Write $I$ for the identity matrix and $J$ for the matrix whose entries
are all one, in the required dimension. With \(v=(a_i)^T\),
\begin{equation}\label{eq:integrated-1}
N=\operatorname{diag}(u_i)M\operatorname{diag}(u_i)
   =vv^T+F,\quad F_{ii}=0,\quad
 F_{ij}=\frac1{1+\sqrt{1+1/(a_i a_j)}}.
\end{equation}
In particular \(\det N=(\prod_i a_i)\det M\), and
\begin{equation}\label{eq:integrated-2}
F_{ij}-\frac12=-\frac{F_{ij}^2}{2a_i a_j}.
\end{equation}
In even dimension \(n\), \(M\) is nonsingular. Indeed, in
\(\mathcal O_K/2\mathcal O_K\), entrywise squaring gives
\[
 (\det M)^2=\det(vv^T+I+J)=1.
\]
Here \(J^2=0\), \((I+J)^{-1}=I+J\), \(\det(I+J)=1\), and
\(v^T(I+J)v=\sum a_i^2+(\sum a_i)^2=0\). The rank-one determinant
identity is valid over this quotient ring even when it has zero divisors.

\begin{lemma}\label{lem:product-four}
An ordered quadruple satisfies
\[
 D>2ABC\quad\hbox{if }A\ge3,\qquad
 D>\frac32ABC\quad\hbox{if }A\ge1,\ B\ge3.
\]
\end{lemma}
\begin{proof}
In the first case Lemma~\ref{lem:tail-irregular} gives \(D>6C\) and
\(D>18\). For \(|z|\le1/9\),
\(\operatorname{Re}\sqrt{1+z}>\frac{15}{16}\), so \eqref{eq:integrated-1}--\eqref{eq:integrated-2} give
\[
 |F_{ij}|<16/31,\qquad |F_{ij}-1/2|<e:=128/8649.
\]
Put \(s=1/2+3e,\ h=3/2+3e\). The matrix \(F_0=(J_4-I_4)/2\)
has singular values \(3/2,1/2,1/2,1/2\), and
\(\|F-F_0\|_2<3e\). Therefore
\[
 |\det F|<hs^3,\qquad \|\operatorname{adj}F\|_2<hs^2.
\]
The latter follows from products of three singular values, including
singular matrices by continuity. The rank-one identity yields
\[
 |\det N|<hs^2\left(s+\sum A_i^2\right)
 <hs^2\left(1+\frac3{36}+\frac{s}{324}\right)D^2<\frac{D^2}{2}.
\]
The last inequality follows by substituting the definitions of \(s\) and \(h\). Since
\(1\le|\det M|=|\det N|/(ABCD)\), this proves the first assertion.

For the second assertion, the three off-diagonal entries of \(E=F-F_0\)
in the row corresponding to \(a\) have modulus less than \(1/18\),
since \(|F_{1j}|<5/9\) and
\(25/486<1/18\). For \(2\le i<j\le4\), we have \(|E_{ij}|<1/64\).
The symmetric nonnegative matrix of these entrywise upper bounds has
weighted row sums \(1/9,11/96\) for weights \((3/2,1,1,1)\).
Its spectral radius bounds $\|E\|_2$, so $\|E\|_2<11/96$.
Put $s=59/96$ and $h=155/96$. As above,
$|\det F|<hs^3$ and $\|\operatorname{adj}F\|_2<hs^2$.
Lemma~\ref{lem:tail-irregular} also gives $D>6C$ and $D>18$ in this case.
Consequently
\[
 |\det N|<hs^2\left(\frac{13}{12}+\frac{s}{324}\right)D^2
 =\frac{18212679025}{27518828544}D^2<\frac23D^2.
\]
The integral determinant bound $ABCD\le|\det N|$ proves the second assertion.
\end{proof}

\begin{lemma}\label{lem:product-five}
Five ordered entries of moduli \(1\le A\le B\le S\le C\le D\),
with \(B\ge3\), satisfy \(D>3ABSC\).
\end{lemma}
\begin{proof}
The four large entries give \(D>2BSC\ge18C\), \(D>54\).
In \eqref{eq:integrated-1} take \(n=5\), \(w=(a,b,s,c)^T\),
\(f=(F_{15},\ldots,F_{45})^T\), and let \(F_{\rm pre}\) be the
submatrix formed by the first four rows and columns of \(F\). The last diagonal entry of \(N\) is \(d^2\), so
\begin{equation}\label{eq:integrated-3}
\det N=d^2\det G,\qquad
 G=F_{\rm pre}-\frac{wf^T+fw^T}{d}-\frac{ff^T}{d^2}.
\end{equation}
The weighted bound just proved gives
\(\|F_{\rm pre}-(J_4-I_4)/2\|_2<1/8\), hence
\(\|F_{\rm pre}^{-1}\|_2<8/3\). The last-column products have
modulus above 54, so \eqref{eq:integrated-2} gives
\(|F_{i5}-1/2|<625/264627<1/400\).
For \(f_0=(1/2,\ldots,1/2)^T\),
\begin{equation}\label{eq:integrated-4}
\|f-f_0\|_2<1/200,\quad \|f\|_2<201/200,\quad
 \|F_{\rm pre}^{-1}f\|_2
 <\frac23+\frac83\left(\frac1{200}+\frac1{12}\right)
 =\frac{203}{225}<\frac{11}{12}.
\end{equation}
The inverse estimate follows by applying the resolvent identity to
\(F_0^{-1}f_0=(1/3,\ldots,1/3)^T\).

For a zero-diagonal symmetric matrix of order four, set
\(p=F_{12}F_{34},q=F_{13}F_{24},r=F_{14}F_{23}\). Then
\[
 \det F_{\rm pre}=p^2+q^2+r^2-2pq-2pr-2qr.
\]
Each of \(p,q,r\) is within \(7/192\) of \(1/4\). Expansion gives
\[
 |\det F_{\rm pre}|<
 \frac3{16}+\frac32\frac7{192}+9\left(\frac7{192}\right)^2
 <\frac{33}{128}.
\]
Put \(h=f^TF_{\rm pre}^{-1}w\), \(p_0=f^TF_{\rm pre}^{-1}f\),
\(q_0=w^TF_{\rm pre}^{-1}w\). From \eqref{eq:integrated-4},
\[
 |h|<11C/6,\qquad |p_0|<1,\qquad |q_0|<32C^2/3.
\]
The rank-two determinant identity in \eqref{eq:integrated-3} is exactly
\[
 \det G=\det F_{\rm pre}
 \left(1-\frac{2h}{d}+\frac{h^2-p_0q_0-p_0}{d^2}\right).
\]
It follows that
\[
 |\det G|<
 \frac{33}{128}\left(1+\frac{11}{54}+\frac{505}{11664}
                              +\frac1{2916}\right)<\frac13.
\]
Moreover
\[
 \|G-F_0\|_2<
 \frac18+\frac{67}{300}+\frac{(201/200)^2}{2916}<\frac12,
\]
so \(G\), and hence \(M\), is nonsingular. Integrality now gives
\[
 1\le|\det M|=\frac{D|\det G|}{ABSC}<\frac{D}{3ABSC}.
\]
\end{proof}

\subsection{Fourth-entry gaps}

\begin{lemma}\label{lem:five-frontier-gaps}
Let an ordered quadruple have \(1\le A\le B\le S\le T\), with
\(B\ge3\). If \(A\ge3\), then \(T>3ABS\).
If \(A<3\), the first two columns of each row in the following
table imply \(T>tABS\). The number \(\alpha\) is a rational lower
bound for \(A\), used in the proof.
\[
\begin{array}{c|c|c|c}
\hbox{condition on }N(a)&B\ge B_0&\alpha&t\\\hline
1&8&1&27/8\\
2&4&140/99&3\\
3&3&12/7&11/4\\
4&3&2&29/10\\
5\le N(a)\le8&3&20/9&3
\end{array}
\]
\end{lemma}
\begin{proof}
The $4\times4$ integral matrix constructed above is nonsingular by
the same congruence argument. We now sharpen the determinant estimate. Write
\[
 M=\operatorname{diag}(u_i)(J+E)\operatorname{diag}(u_i),\quad
 E_{ii}=0,\quad E_{ij}=\delta_{ij}
 =\sqrt{1+z_{ij}}-1,\quad z_{ij}=\frac1{a_i a_j}.
\]
With triangles and paths taken in the complete graph on the four indices,
the cubic part of \(\det(J+E)\) is
\begin{equation}\label{eq:integrated-7}
\mathcal C(\delta)=
 2\sum_{\rm triangles}\prod_{\rm edges}\delta
 -2\sum_{\rm paths}\prod_{\rm edges}\delta
 +2\sum_{\{i,j\}}\delta_{ij}^{\,2}\delta_{kl},
\end{equation}
where the middle sum runs over the twelve undirected Hamiltonian paths,
and \(\{k,l\}\) is complementary to \(\{i,j\}\) in the last sum.
Its quartic part is \(\det E\), the sum of the squares of
\(\delta_{12}\delta_{34}\), \(\delta_{13}\delta_{24}\), and
\(\delta_{14}\delta_{23}\), minus twice the sum of their pairwise products. These identities follow by the 24 permutations of the determinant.
Substituting \(\delta=z/2\) in \eqref{eq:integrated-7} gives
\begin{equation}\label{eq:integrated-8}
abst\,\mathcal C(z/2)
 =\frac14\sum_i\frac{a_i}{\prod_{j\ne i}a_j}
                  -\frac14\sum_{i<j}\frac1{a_i a_j}.
\end{equation}
Indeed, each triangle gives a term \(a_i/\prod_{j\ne i}a_j\)
with coefficient \(1/4\). For each term \(1/(a_i a_j)\), two paths
and one squared-edge term contribute.

For \(0\le\rho\le1/5\), put \(k(\rho)=1/2+\rho/7\).
Rationalization and the identity
\(\delta-z/2=-\delta^2/2\) imply
\begin{equation}\label{eq:integrated-9}
|\delta|\le k(\rho)|z|,\qquad
 |\delta-z/2|\le(\rho/7)|z| \quad (|z|\le\rho).
\end{equation}
To check both bounds explicitly, use
\[
 (1-\rho)(7+2\rho)^2-(7-2\rho)^2
   =\rho(7-28\rho-4\rho^2)\ge0,\qquad
 k(1/5)^2<2/7.
\]
In every case considered here, \(AB>5\). Set
\[
 k_1=\frac12+\frac1{7AB},\quad k_2=\frac12+\frac1{7B^2},
 \quad K_1=k_1^2k_2,\quad L_1=6K_1-\frac12,\quad
 Q_1=9k_1^2k_2^2.
\]
Every cubic monomial has at most two occurrences of edges incident to
the smallest entry, counting repeated edges. Expansion about the three
leading factors in \eqref{eq:integrated-9} therefore bounds its remainder by
\((K_1-1/8)\prod|z_{ij}|\).
The sums of the absolute cubic coefficients are 2 for terms contributing
to \(a_i/\prod_{j\ne i}a_j\), and 6 for those contributing to
\(1/(a_i a_j)\). Combining these with \eqref{eq:integrated-8}, and observing
that each quartic monomial has exactly two incident edge occurrences,
gives
\begin{equation}\label{eq:integrated-10}
1\le|\det M|\le
 2K_1\left(\frac{T}{ABS}+\frac{S}{ABT}
                    +\frac{B}{AST}+\frac{A}{BST}\right)
 +L_1\sum_{i<j}\frac1{A_iA_j}+\frac{Q_1}{ABST}.
\end{equation}

Put \(\tau=T/(ABS)\). For fixed \(\tau>0\), substituting
\(T=\tau ABS\) and \(S\ge B\) gives
\[
 \frac{S}{ABT}+\frac{B}{AST}+\frac{A}{BST}
 \le\frac1\tau\left(\frac2{A^2B^2}+\frac1{B^4}\right),
\]
\[
 \sum_{i<j}\frac1{A_iA_j}
 \le\frac2{AB}+\frac1{B^2}
       +\frac1\tau\left(\frac1{A^2B^2}+\frac2{AB^3}\right),
 \qquad
 \frac1{ABST}\le\frac1{\tau A^2B^4}.
\]
In these positive expressions, replace \(A,B\) by rational lower
bounds \(\alpha,B_0\). Evaluate \(k_1,k_2\), and hence
\(K_1,L_1,Q_1\), at \(A=\alpha,B=B_0\). Let
\[
 P_0=\frac2{\alpha B_0}+\frac1{B_0^2},\qquad
 J_0=2K_1\left(\frac2{\alpha^2B_0^2}+\frac1{B_0^4}\right)
 +L_1\left(\frac1{\alpha^2B_0^2}+\frac2{\alpha B_0^3}\right)
 +\frac{Q_1}{\alpha^2B_0^4}.
\]
Equation \eqref{eq:integrated-10} now bounds the determinant by the convex function
\[
 F(\tau)=2K_1\tau+L_1P_0+\frac{J_0}{\tau}.
\]
If \(A\ge3\), take \(\alpha=B_0=3\).
Lemma~\ref{lem:product-four} gives \(\tau>2\), while direct rational
substitution gives \(F(2)<7/10\) and \(F(3)<19/20\).
Convexity excludes \(2<\tau\le3\).

If \(A<3\), use the lower bounds in the table.
Lemma~\ref{lem:tail-irregular} gives \(\tau>2/A>2/3\).
For all five explicitly listed triples \((\alpha,B_0,t)\), direct rational
substitution gives
\[
 F(2/3)<\frac12,\qquad F(t)<\frac{499}{500}<1.
\]
The largest of the five right endpoint values is
\(964921210337222547/967173115740160000<499/500\),
in the norm-two row. The rational lower bounds for \(A\) follow by squaring.
All satisfy \(\alpha B_0>5\).
Convexity excludes \(2/3<\tau\le t\), proving all rows at once.
\end{proof}

\section{Interpolation of square roots}

We construct integral combinations of square roots by interpolation. The interpolation weights cancel the low-order terms involving the fixed entries, leaving an explicit remainder. When these combinations are nonzero, bounds for them give gaps between the last two entries. We also bound the ranges in which these combinations can vanish.

Let \(a_1,\ldots,a_m,p,q\) be distinct tuple entries, \(m=2,3,4\), with
\(1\le A_1\le\cdots\le A_m\le P=|p|\le Q=|q|\) and \(A_1P>1\).
Choose \(\eta^2=pq\), put \(x=1/p,y=1/q\), and use
\(f(w)=\sqrt{1+w}\) on its analytic branch at zero. Define
\[
 z=\eta f(xy),\qquad
 R_i=a_i\eta f(x/a_i)f(y/a_i),\qquad
 \Delta_m=\prod_{i<j}(a_i-a_j).
\]
Their squares are \(pq+1\) and \((a_ip+1)(a_iq+1)\).
Thus \(z\in\mathcal O_K\). Each \(R_i\) equals, up to sign, the
product of integral square roots of \(a_ip+1\) and \(a_iq+1\), so
\(R_i\in\mathcal O_K\) as well. We call the fixed entries $a_1,\ldots,a_m$ the anchors.
With the common branch choices above, define
\begin{equation*}
 L_m=-\Delta_m z+
 \sum_{i=1}^m
 \frac{\Delta_m a_i^{m-2}}{\prod_{j\ne i}(a_i-a_j)}R_i
 \ \in\mathcal O_K.
\end{equation*}
The coefficients are integer polynomials, since each denominator divides
the Vandermonde product up to sign. A nonzero \(L_m\) has modulus at least one.

Let \(h_j\) be the complete homogeneous symmetric polynomial of degree
\(j\) in \(1/a_1,\ldots,1/a_m\), and write
\(f(xW)f(yW)=\sum_{n\ge0}f_n(x,y)W^n\).
Partial fractions of \(\prod_i(1-W/a_i)^{-1}\) give
\[
 \sum_i\frac{a_i^{m-1-n}}{\prod_{j\ne i}(a_i-a_j)}
 =\begin{cases}
 1&n=0,\\
 0&1\le n<m,\\
 (-1)^{m-1}h_{n-m}/\prod_i a_i&n\ge m.
 \end{cases}
\]
Hence the exact identity is
\begin{equation}\label{eq:coherent-exact}
 \frac{L_m}{\Delta_m\eta}
 =1-f(xy)+\frac{(-1)^{m-1}}{\prod_i a_i}F_m,\qquad
 F_m=\sum_{n\ge m}f_n(x,y)h_{n-m}.
\end{equation}
All series converge absolutely. Indeed, for \(b_n=\binom{1/2}{n}\),
\[
 \left(\sum_{n\ge0}|b_n|W^n\right)^2
 =(2-\sqrt{1-W})^2=5-W-4\sqrt{1-W},
\]
so \(|f_n(x,y)|\le4|b_n|P^{-n}\) for \(n\ge2\).
As \(h_j\) has \(\binom{j+m-1}{m-1}\) monomials, and \(|b_n|\)
decreases for \(n\ge1\), we obtain the complete remainder estimate
\begin{equation}\label{eq:coherent-tail}
 |F_m-f_m|\le
 4\left|\binom{1/2}{m+1}\right|P^{-m}
 \left[\left(1-\frac1{A_1P}\right)^{-m}-1\right].
\end{equation}
Direct multiplication gives
\begin{equation}\label{eq:integrated-5}
f_2=-\frac{(x-y)^2}{8},\qquad
 f_3=\frac{(x^2-y^2)(x-y)}{16},\qquad
 f_4=-\frac{(x-y)^2(5x^2+6xy+5y^2)}{128}.
\end{equation}
Whenever \(PQ>10000\), rationalization and
\(\operatorname{Re}\sqrt{1+w}\ge\sqrt{1-|w|}\) give
\begin{equation}\label{eq:coherent-near-one}
 \frac{99}{100}<|f(xy)|<\frac{101}{100},\qquad
 \frac{100}{201PQ}<|1-f(xy)|<\frac{100}{199PQ}.
\end{equation}

\begin{lemma}\label{lem:interpolation-parameters}
Let \((H,\lambda,U)\) be either row of the table
\[
 \begin{array}{c|c|c}
 H&\lambda&U\\\hline
 32&1/10&179/10\\
 48&19/200&15
 \end{array}.
\]
Suppose \(A_1P\ge H,\ Q\ge4P,\ PQ>10000\). For two anchors
of moduli \(X\le Y\), and \(Z=|z|\),
\[
 |L_2|<\frac{21Y}{20Z}+\frac{Z}{2XP^2}.
\]
For three anchors of moduli \(X,Y,V\),
\[
 |L_3|<
 \frac{51|\Delta_3|}{100Z}
 +\lambda\frac{|\Delta_3|Z}{XYV P^3},\qquad
 L_3=0\ \Longrightarrow\
 5XYVP^2<Q<UXYVP^2.
\]
\end{lemma}
\begin{proof}
For two anchors, \eqref{eq:integrated-5} and the tail estimate bound \(P^2|F_2|\) by
\[
 \frac{25}{128}+\frac14\left[\left(\frac{H}{H-1}\right)^2-1\right].
\]
Equation~\eqref{eq:coherent-exact}, \(|\Delta_2|\le2Y\), and
\eqref{eq:coherent-near-one} give the required estimate because
\[
 \frac{202}{199}<\frac{21}{20},\qquad
 \frac{200}{99}\left\{\frac{25}{128}
 +\frac14\left[\left(\frac{H}{H-1}\right)^2-1\right]\right\}<\frac12.
\]
For three anchors, put $w=y/x$ and $\rho=|w|\le1/4$.
The factorization $(1-w)^2(1+w)=(1-w)(1-w^2)$ gives
\[
 (1-\rho)(1-\rho^2)\le16P^3|f_3|
 \le(1+\rho)(1+\rho^2).
\]
Thus $45/1024\le P^3|f_3|\le85/1024$. Put
\[
 r_H=\frac5{32}\left[\left(\frac{H}{H-1}\right)^3-1\right],
 \quad l_H=\frac{45}{1024}-r_H,\quad
 u_H=\frac{85}{1024}+r_H.
\]
Then \(l_H/P^3\le|F_3|\le u_H/P^3\), with \(l_H>0\).
The required rational comparisons are
\[
 \frac{101}{199}<\frac{51}{100},\quad
 \frac{100u_H}{99}<\lambda,\quad
 \frac{100}{201u_H}>5,\quad
 \frac{100}{199l_H}<U.
\]
For clarity, the last two exact fractions are
\[
 \begin{array}{c|cc}
 H&100/(201u_H)&100/(199l_H)\\\hline
 32&610119680/120943911&122023936/6879629\\
 48&2126295040/398301399&2126295040/142842001
 \end{array}.
\]
Substitution in the exact identity proves both the direct estimate and,
when \(L_3=0\), the two strict bounds for \(Q\).
\end{proof}

\begin{lemma}\label{lem:strong-two-anchor}
In the two-anchor situation of Lemma~\ref{lem:interpolation-parameters},
assume also \(Y\ge3,\ P\ge5Y\). If \(Q\ge5XYP\), then
\[
 Q>\frac{X^2P^3}{h_0},\qquad h_0=\frac{10001}{10000}.
\]
\end{lemma}
\begin{proof}
Here \(L_2=R_a-R_b-(a-b)z\). If it vanishes, subtracting the identities
for \(R_a^2,R_b^2\) and then squaring gives
\[
 (a+b-p-q)^2=4(ab+1)(pq+1).
\]
Solving this quadratic in \(q\) gives
\[
 q=a+b+p+2abp\pm2\sqrt{(ab+1)(ap+1)(bp+1)}.
\]
Since \(XY\ge3,XP\ge15,YP\ge45\),
\[
 \frac{Q}{XYP}\le\frac{19}{45}+2+
             2\sqrt{\frac43\frac{16}{15}\frac{46}{45}}<5.
\]
The last comparison is \(11776<13456\) after squaring.
Thus \(Q\ge5XYP\) makes \(L_2\ne0\). Moreover
\[
 \frac{Z^2}{Y^2}\ge\frac{5XP^2}{Y}-\frac1{Y^2}
 \ge5XP-1\ge159.
\]
The first term in the bound of Lemma~\ref{lem:interpolation-parameters}
is less than \(1/2\). Since \(1\le|L_2|\), the second term exceeds
\(1/2\), giving \(Z>XP^2\).
Finally \(Z^2\le PQ+1<h_0PQ\), proving the assertion.
\end{proof}

\begin{lemma}\label{lem:general-three}
For three anchors \(A\le B\le S\), suppose \(P\ge4S,\ AP\ge9,\ Q\ge P\).
Then
\[
 |L_3|<\frac{5BS^2}{Z}+\frac{3SZ}{AP^3}.
\]
\end{lemma}
\begin{proof}
The absolute coefficient series gives \(|f_n|\le P^{-n}/4\) for \(n\ge3\),
and hence
\[
 |F_3|\le\frac1{4P^3(1-1/(AP))^3}.
\]
Also \(P^2\ge4AP\ge36\), \(|\Delta_3|\le8BS^2\),
\(35/36<|f(xy)|<37/36\), and
\(|1-f(xy)|<36/(71PQ)\). The two contributions in
\eqref{eq:coherent-exact} are bounded by
\((296/71)BS^2/Z\) and \((6561/2240)SZ/(AP^3)\), respectively.
\end{proof}

\begin{lemma}\label{lem:separated-three}
For three anchors \(A\le B\le S\le P\), if \(P\ge1000,\ Q\ge4P\), then
\[
 |L_3|<\frac{5BS^2}{Z}+\frac{3SZ}{4AP^3},\qquad
 L_3=0\ \Longrightarrow\ 5ABSP^2<Q<15ABSP^2.
\]
\end{lemma}
\begin{proof}
Use the calculation in the proof of Lemma~\ref{lem:interpolation-parameters}
with \(r_{1000},l_{1000},u_{1000}\).
Since \(|\Delta_3|\le8BS^2\), the needed direct comparisons are
\(808/199<5\) and \(800u_{1000}/99<3/4\).
For \(L_3=0\), the required rational comparisons are
\(100/(201u_{1000})>5\) and \(100/(199l_{1000})<15\).
\end{proof}

\begin{lemma}\label{lem:four-anchor-nonvanishing}
For ordered anchors \(1\le A\le B\le S\le T\le P\le Q\), assume
\(AP>1\) and \(P^2\ge32BST/A^3\). The coherent forms \(L_3\) on
\((a,b,s)\) and \((a,b,t)\), using the same \(z,R_a,R_b\), cannot both vanish.
\end{lemma}
\begin{proof}
If both vanished, the quadratic polynomials interpolating \(uR_u\)
at the two triples would have leading coefficient \(z\) and agree at
\(a,b\), so they would coincide. Write the common polynomial
\(J(X)=zX^2+vX+w\). Since \(z^2-pq=1\),
\[
 J(X)^2-X^2(pX+1)(qX+1)=(X-a)(X-b)(X-s)(X-t).
\]
For the elementary symmetric functions \(e_j\) of the anchors,
\(w^2=e_4,\ 2vw=-e_3,\ v^2+2zw-1=e_2\). Thus
\[
 pq+1=\frac{(e_2+1-e_3^2/(4e_4))^2}{4e_4}.
\]
Using \(|e_2|\le6ST,\ |e_3|\le4BST,\ |e_4|=ABST\), this gives
\[
 |pq+1|\le\frac{(6ST+1+4BST/A)^2}{4ABST}
 \le\frac{121BST}{4A^3}.
\]
But with \(r=BST/A^3\ge1\),
\(|pq+1|\ge PQ-1\ge P^2-1\ge32r-1>121r/4\), a contradiction.
\end{proof}

\begin{lemma}\label{lem:direct-four}
For four anchors \(A\le B\le S\le T\le P\), if
\(P\ge1000,\ Q\ge1000P\), then
\[
 |L_4|<\frac{13BS^2T^3}{Z}+\frac{49ST^2Z}{50AP^4},\qquad
 L_4=0\ \Longrightarrow\ 10ABSTP^3<Q<17ABSTP^3.
\]
\end{lemma}
\begin{proof}
Expanding the Vandermonde determinant gives \(24\) monomials.
Since \(A\le B\le S\le T\), each has modulus at most
\(BS^2T^3\), by successively interchanging any inverted pair of
exponents. Thus \(|\Delta_4|\le24BS^2T^3\).

Put \(k=1/1000\) and
\[
 r=\frac7{64}\left[\left(\frac{1000}{999}\right)^4-1\right],\quad
 l=\frac{(1-k)^2(5-6k-5k^2)}{128}-r,\quad
 u=\frac{(1+k)^2(5+6k+5k^2)}{128}+r.
\]
Formula \eqref{eq:integrated-5} and \eqref{eq:coherent-tail} give
\(l/P^4\le|F_4|\le u/P^4\), with \(l>0\).
The exact identity proves the assertions using
\[
 \frac{24\cdot101}{199}<13,\quad
 \frac{2400u}{99}<\frac{49}{50},\quad
 \frac{100}{201u}>10,\quad \frac{100}{199l}<17.
\]
All four are strict rational comparisons.
\end{proof}

\subsection{A gap estimate for the last two entries}
\begin{lemma}\label{lem:retained-gap}
Let six ordered entries have \(1\le A\le B\le S\le T\le C\le D\),
\(B\ge3,\ T>2ABS\). Choose either row \((H,\lambda,U)\) of
Lemma~\ref{lem:interpolation-parameters}, and suppose \(AT\ge H\).
Put
\[
 \mathcal Q=\frac{|(a-b)(a-s)(b-s)|}{BS^2},\quad
 \mu=\frac{51051}{400000},\quad c_\lambda=\frac{9999}{10000\lambda^2}.
\]
Then
\begin{align}
 C&>c_\lambda\left(\frac1{\mathcal Q}-\frac{\mu}{A^2B^2}\right)^2
                     \frac{A^2T^5}{S^2}
   >\frac{A^2T^5}{S^2},                                      \label{eq:retained-main}\\
 D&>\frac{A^2C^5}{S^2},\qquad
 D>\frac{A^2C^7}{S^2T^4}.                                    \label{eq:retained-start}
\end{align}
Moreover \(C>1000,\ C>200S,\ C>4T,\ AC>1000,\ D>1000C\).
The conclusions for \(C\) use only the first five entries.
\end{lemma}
\begin{proof}
Lemma~\ref{lem:product-five} gives \(C>3ABST\). Since \(T>18\) and \(AT\ge H\ge32\),
\[
 C>27H\ge864,\quad TC>486H>10000,\quad
 C/T>27,\quad C/S>9H>200.
\]
Also \(T>5B,T>5S\), and \(C>5ABT\). Lemma~\ref{lem:strong-two-anchor},
applied to \(a,b,t,c\), gives \(C>A^2T^3/h_0>18^3/h_0>1000\).
As
\[
 \frac{A^2T^2}{BS}>4A^4BS\ge36>5h_0,
\]
we have \(C>5BST\), so a second application on \(b,s,t,c\) gives
\(C>B^2T^3/h_0\).
The form \(L_3\) on \(a,b,s\), evaluated at \(t,c\), cannot vanish, since
\[
 \frac{BT}{AS}>2B^2\ge18>Uh_0
\]
for both parameter rows, contradicting the upper bound for \(C\) when \(L_3=0\).

Let \(W=|\sqrt{tc+1}|\) and \(k=3999/4000\). From the second application of Lemma~\ref{lem:strong-two-anchor}
and \(W^2>kTC\),
\[
 X_0:=\frac{5BS^2}{W}
 <5\sqrt{h_0/k}\frac{S^2}{T^2}
 <\frac5{4A^2B^2}\sqrt{h_0/k}
 <\frac{5005}{4000A^2B^2}.
\]
The first term in the \(L_3\) estimate of
Lemma~\ref{lem:interpolation-parameters} is
\((51/500)\mathcal QX_0<\mu\mathcal Q/(A^2B^2)\).
Since \(L_3\) is a nonzero algebraic integer, the preceding estimate gives
\[
 W>\alpha\frac{AT^3}{S},\qquad
 \alpha=\lambda^{-1}\left(\frac1{\mathcal Q}
                              -\frac{\mu}{A^2B^2}\right).
\]
Since \(0<\mathcal Q\le8\), \(AB\ge3\), and \(\lambda\le1/10\),
\(\alpha>1\) and \((9999/10000)\alpha^2>1\).
Also \(\alpha^2A^2T^6/S^2>T^4>18^4>10000\).
After squaring, subtracting one causes a relative loss of less than \(1/10000\),
proving \eqref{eq:retained-main}.

Repeat this five-entry argument on \(a,b,s,c,d\). Here
\(C>2ABS\) and \(AC\ge H\). This gives the first inequality in
\eqref{eq:retained-start} and \(D/C>A^2C^4/S^2\ge C^2>1000\).
The form $L_4$ on $(a,b,s,t)$, evaluated at $(c,d)$, satisfies the hypotheses of Lemma~\ref{lem:direct-four}. If it vanished, that lemma would give $D<17ABSTC^3$. Together with $D>A^2C^5/S^2$, this would imply $C^2<17BS^3T/A$, contrary to
\[
 \frac{C^2}{17BS^3T/A}>
 \frac9{17}A^3\frac{BT}{S}>
 \frac{18}{17}A^4B^2>1.
\]
For its first summand, use the bounds \(D>A^2C^5/S^2\) and
\(C>A^2T^5/S^2>16BST\), the latter following from \(T>2ABS\).
They give
\[
 \left(\frac{13BS^2T^3}{|\sqrt{cd+1}|}\right)^2
 <\frac{169B^2S^6T^6}{kA^2C^6}
 <\frac{169}{16^6kA^2B^4}<\frac1{10000}.
\]
Thus \(|\sqrt{cd+1}|>(99/98)AC^4/(ST^2)\).
Squaring and subtracting one proves the last inequality in
\eqref{eq:retained-start}, since
\[
 \left(\left(\frac{99}{98}\right)^2-1\right)
       \frac{A^2C^8}{S^2T^4}>1
\]
follows from \(S,T\le C\) and \(C>1000\).
The preceding estimates also give all the stated size inequalities.
\end{proof}

\section{Simultaneous approximation}

We construct simultaneous approximations using two or three fixed tuple entries.  All polynomial divisions below are performed over $\mathbb Z$ before specializing the fixed entries. No unique factorization assumption on $\mathcal O_K$ is needed.

\subsection{A common residue construction}

\begin{lemma}\label{lem:ordinary-residue-construction}
Let $a_1,\ldots,a_m,c,d$, where $m\ge2$, be distinct nonzero members of
a $D(1)$-tuple in $\mathcal O_K$.  Put
\[
 A_j=|a_j|,\quad A=\min_j A_j,\quad S=\max_j A_j,
 \quad C=|c|>1,\quad AC>1,
\]
\[
 H=\prod_{j=1}^m a_j,\qquad
 \Delta=\prod_{1\le i<j\le m}(a_i-a_j),\qquad
 \eta=1-\frac1{AC},\qquad
 e(X)=\sqrt{1+\frac1{XC}}\left(1+\frac XC\right).
\]
For every integer $k\ge1$ there is an explicit matrix
$Q_k=(Q_{ijk})_{0\le i,j\le m}$ of integral polynomials in $a_1,\ldots,a_m$ and $c$.  Its $j$th column, for $j\ge1$, is divisible by
$a_j^{(m-1)k}$, and
\begin{equation}\label{eq:ordinary-residue-det}
 \det Q_k=
 \left((-1)^{m(m+1)/2}4^{m+1}H^{m-1}\Delta^{2m}\right)^k
 \prod_{v=-1}^{m-1}\binom{1/2-vk}{k}\ne0.
\end{equation}
With $\theta_0=1$, $\theta_j=\sqrt{1+1/(a_jc)}$ defined by the binomial series,
\begin{equation}\label{eq:ordinary-residue-error}
 \left|\sum_{j=0}^mQ_{ijk}\theta_j\right|
 \le V_{m,k}:=
 \left(\frac{4|\Delta|^2}{|H|C^m}\right)^k
 \eta^{-mk+1/2}
 \frac{\Gamma(k+3/2)\Gamma(mk-1/2)}
 {\pi\Gamma((m+1)k+1)}.
\end{equation}
If $z\in\mathcal O_K$, $z^2=cd+1$, and $Z=|z|$, then
\begin{equation}\label{eq:ordinary-root-transfer}
 1\le\frac{Ce(S)}Z
       \max_i\sum_{j=1}^m\frac{|Q_{ijk}|}{A_j}
       +V_{m,k}Z.
\end{equation}
At $k=1$ the left side can be replaced by $2$.
\end{lemma}

\begin{proof}
Set $\alpha_0=0$, $\alpha_j=H/a_j$, and $N=Hc$.  With $\delta_{ij}$
the Kronecker delta, define
\begin{equation}\label{eq:ordinary-explicit-residue}
 f_{ij}=
 \sum_{\substack{h_0,\ldots,h_m\ge0\\
                  h_0+\cdots+h_m=k+\delta_{ij}-1}}
 \binom{k+1/2}{h_j}N^{-k}(N+\alpha_j)^{k-h_j}
 \prod_{\ell\ne j}
 \binom{-k-\delta_{i\ell}}{h_\ell}
 (\alpha_j-\alpha_\ell)^{-k-\delta_{i\ell}-h_\ell},
\end{equation}
and
\[
 Q_{ijk}=(4H^mc\Delta^2)^k f_{ij}.
\]
The residue theorem, applied at the interpolation points, gives
\begin{equation}\label{eq:ordinary-contour}
 \frac1{2\pi\mathrm i}\int_\Gamma
 \frac{(1+w/N)^{k+1/2}}
 {(w-\alpha_i)\prod_{\ell=0}^m(w-\alpha_\ell)^k}\,dw
 =\sum_{j=0}^mf_{ij}\theta_j.
\end{equation}
Here $\Gamma$ surrounds the interpolation points and avoids the ray from
$-N$ away from $0$.  Such a contour exists because
$|\alpha_j/N|\le1/(AC)<1$.  Formula
\eqref{eq:ordinary-explicit-residue} follows by differentiating the
numerator and the other pole factors at each pole.  This residue construction appears in \cite[\S7]{JZ}.

We prove integrality term by term.  Apart from
a sign and the scalar
\[
 4^k\binom{k+1/2}{h_j}
       \prod_{\ell\ne j}\binom{-k-\delta_{i\ell}}{h_\ell},
\]
a summand in column zero is
\[
 c^{k-h_0}\Delta^{2k}
       \prod_{r=1}^m a_r^{\delta_{ir}+h_r},
\]
and a summand in column $j\ge1$ is
\begin{equation}\label{eq:ordinary-polynomial-summand}
 a_j^{(m-1)k}(a_jc+1)^{k-h_j}
 \prod_{\substack{\ell\ne j\\\ell\ge1}}
       a_\ell^{\delta_{i\ell}+h_\ell}
 \frac{\Delta^{2k}}
 {\displaystyle\prod_{\substack{\ell\ne j\\\ell\ge1}}
 (a_\ell-a_j)^{k+\delta_{i\ell}+h_\ell}}.
\end{equation}
For $\ell\ne j$ one has $\delta_{i\ell}+h_\ell\le k$, and
$h_j\le k$.  Thus all factors $a_\ell-a_j$ in the denominators cancel, and every remaining exponent is nonnegative.  Also
$4^{h_j}\binom{k+1/2}{h_j}\in\mathbb Z$. To see this, expand
$(1+x)^k(1+x)^{1/2}$ and use
$\binom{1/2}{r}=(-1)^{r-1}2C_{r-1}/4^r$ for $r\ge1$, where
$C_r$ is a Catalan number.  This proves polynomial integrality and the
claimed column divisibility.  At $k=1$ the scalar factors are $4$, $6$,
or multiples of $4$, so every entry is even.

The determinant identity for the residue coefficients is
\[
 \det(f_{ij})=
 N^{-(m+1)k}
 \left(\prod_{v=-1}^{m-1}\binom{1/2-vk}{k}\right)
 \prod_{\ell=0}^m\prod_{\substack{r=0\\r\ne\ell}}^m
             (\alpha_r-\alpha_\ell)^{-k}.
\]
This is the unnumbered determinant identity following (2.4) in
\cite[\S2, p.~1721]{Bennett96}. With his interpolation points replaced by
$\alpha_j$ and $\nu=1/2$, his coefficients satisfy $P_{ij}(1/N)=f_{ij}$,
since $(1+\alpha_j/N)^{k-h_j}N^{-h_j}=N^{-k}(N+\alpha_j)^{k-h_j}$.
His integral (2.1) already includes the extra pole factor
$(w-\alpha_i)^{-1}$ in \eqref{eq:ordinary-contour}.
Although stated there for distinct integer interpolation points and real $x$,
the identity extends to complex interpolation points with $\alpha_0=0$. Substitute the rational expressions
\eqref{eq:ordinary-explicit-residue} and clear all denominators.  The
result is a polynomial identity, valid for all distinct integer points
and an interval of real $x$, and hence identically valid over
$\mathbb C$.  The product of the pairwise differences of the interpolation points is, up to sign, $H^{m(m-1)/2}\Delta$.  Multiplication by the clearing factor gives
\eqref{eq:ordinary-residue-det}.  All its factors are nonzero. Each
numerator factor in the displayed binomial coefficients is a nonzero
half-integer.

Deform the contour in \eqref{eq:ordinary-contour} onto
$w=-N(1+x)$, $x\ge0$.  The numerator's branch jump has modulus
$2x^{k+1/2}$, while every pole factor after division by $N$ has modulus at least
$x+\eta$.  The large-circle contribution vanishes because the
integrand is $O(|w|^{-mk-1/2})$, and the small-circle contribution at
the branch point also tends to zero.  Multiplying by the clearing
factor therefore bounds the residual by
\[
 \left(\frac{4|\Delta|^2}{|H|C^m}\right)^k\frac1\pi
 \int_0^\infty\frac{x^{k+1/2}}{(x+\eta)^{(m+1)k+1}}\,dx.
\]
The substitution $x=\eta y$ and the beta integral give
\eqref{eq:ordinary-residue-error}.

Choose integral roots $r_j^2=a_jc+1$, $v_j^2=a_jd+1$ and put
$u_j=\pm r_jv_j/(a_jz)$. Here $z\ne0$, since otherwise
$|d|=1/C<1$, and $\theta_j\ne0$, since $A_jC>1$. We have
\[
 u_j^2=\theta_j^2\left(1+\frac{c-a_j}{a_jz^2}\right).
\]
Write $u_j=\theta_jw_j$ and choose the sign so that
$\operatorname{Re}w_j\ge0$. Then $|w_j+1|\ge1$, and hence
\[
 |\theta_j-u_j|
 =|\theta_j|\,|w_j-1|
 \le\frac{|\theta_j|\,|c-a_j|}{A_jZ^2}
 \le\frac{Ce(A_j)}{A_jZ^2}
 \le\frac{Ce(S)}{A_jZ^2}.
\]
The last inequality follows from
$(e(X)^2)'=[2+2X/C-1/X^2+1/C^2]/C>0$ for $X\ge1$.
The vector with coordinates
\[
 q=Hz,\qquad n_j=(H/a_j)(\pm r_jv_j)
\]
is nonzero.  Column divisibility makes every row form
$Q_{i0k}q+\sum_{j\ge1}Q_{ijk}n_j$ a multiple of $H$ in
$\mathcal O_K$.  Nonsingularity makes some row form nonzero, so its
modulus is at least $|H|$.  Division by $|H|$ and the triangle
inequality prove \eqref{eq:ordinary-root-transfer}.  At $k=1$ all
these forms are multiples of $2H$, proving the stronger assertion.
\end{proof}

\subsection{Two fixed entries at every order}

\begin{proposition}\label{lem:ordinary-all-orders}
Let $a,b,c,d$ be distinct nonzero members of a $D(1)$-tuple, with
$1\le A=|a|\le B=|b|$, $C=|c|>4B$ and $AC>1000$.
Put $\delta=|a-b|$ and
\[
 H_2=42B^3C,\qquad
 R_2=\frac{27B(AC-1)^2}{16A\delta^2}
      >\frac{2AC^2}{5B}.
\]
For $z^2=cd+1$, $Z=|z|$, and every integer $k\ge1$,
\begin{equation}\label{eq:ordinary-integer-k}
 1<\frac{3CH_2^k}{BZ}+\frac{Z}{2R_2^k}.
\end{equation}
If $R_2>1$ and $\lambda=1+\log H_2/\log R_2<2$, then
\begin{equation}\label{eq:ordinary-continuous}
 Z^{2-\lambda}<253B^2C^2.
\end{equation}
\end{proposition}

\begin{proof}
Use Lemma~\ref{lem:ordinary-residue-construction} with $m=2$.
Its interpolation points are $0,b,a$ and its integral matrix has
\[
 \det Q_k=(-64ab(a-b)^4)^k
           \prod_{v=-1}^{1}\binom{1/2-vk}{k}\ne0.
\]
We bound columns $1$ and $2$ using local residue integrals.
At either nonzero interpolation point let $u\le v$ be its two
distances to the other points, and integrate on the circle of radius
$2u/5$.  The product of the three pole distances is at least
\[
 \frac{6u^2(v-2u/5)}{25}.
\]
The additional factor $|w-\alpha_i|$ is at least $2u/5$, which cancels the factor given by the circle length divided by $2\pi$.  Also $|w|\le7B/5$.  Put
\[
 F_2=1+\frac7{5AC},\qquad
 p_2=\sqrt{\frac{F_2}{1-1/(AC)}}.
\]
The modulus of the numerator, divided by $|\theta_j|$, is at most
$p_2F_2^k$.  The circles lie in $|w|<|abc|$, so they support
the specified analytic branch.

At the point $a$ the two distances are $A,\delta$, and at $b$ they
are $B,\delta$.  In both cases
\begin{equation}\label{eq:ordinary-two-distance}
 \frac{v^2}{v-2u/5}\le\frac{5B}{2}.
\end{equation}
Indeed, if $\delta\le A$ at $a$, or $\delta\le B$ at $b$, the
ratio is at most $5B/3$.  Otherwise set $x=A/B\le1$ and use
$\delta\le A+B$.  The two ratios divided by $B$ are bounded by
$(1+x)^2/(1+3x/5)$ and $(1+x)^2/(x+3/5)$, respectively.
Both functions increase on $[0,1]$ and equal $5/2$ at $1$.

For column $2$, whose pole is $a$, the local estimate before
\eqref{eq:ordinary-two-distance} is
\[
 |Q_{i2k}|\le p_2
 \left(\frac{50}{3}B^2CF_2
                  \frac{v^2}{v-2u/5}\right)^k.
\]
Column $1$ has the same expression with $A^2$ in place of $B^2$
and the distances taken at $b$.  These formulas follow by multiplying
the integral estimate by $(4a^2b^2c(a-b)^2)^k$ and using $uv=A\delta$ at $a$ and $uv=B\delta$ at $b$.
Since $(125/3)(1+7/5000)<42$, they imply
\[
 |Q_{i2k}|<p_2H_2^k,\qquad
 |Q_{i1k}|<p_2\bigl((A/B)^2H_2\bigr)^k.
\]
Consequently
\[
 \frac{Ce(B)}Z\sum_{j=1}^2\frac{|Q_{ijk}|}{A_j}
 <\frac{2p_2e(B)CH_2^k}{BZ}
 <\frac{3CH_2^k}{BZ},
\]
where $A_1=A,A_2=B$.  The last strict comparison follows, for
example, by squaring and using
\[
 p_2^2\le\frac{5007}{4995},\qquad
 e(B)^2<\frac{25\cdot1001}{16000},\qquad
 4\frac{5007}{4995}\frac{25\cdot1001}{16000}<9.
\]

We next bound the residual uniformly for all $k\ge1$.  Write $\eta=1-1/(AC)$.  The exact beta bound in
\eqref{eq:ordinary-residue-error} gives
\[
 V_{2,k}R_2^k=\eta^{1/2}b_k,\qquad
 b_k=\left(\frac{27}{4}\right)^k
       \frac{\Gamma(k+3/2)\Gamma(2k-1/2)}
            {\pi\Gamma(3k+1)}.
\]
Here $b_1=27/64$, and
\[
 \frac{b_{k+1}}{b_k}=
 \frac{(27/4)(k+3/2)(2k-1/2)(2k+1/2)}
      {(3k+1)(3k+2)(3k+3)}<1,
\]
because denominator minus numerator is
$(432k^2+1110k+273)/32>0$.  Thus
$V_{2,k}<(27/64)R_2^{-k}<(1/2)R_2^{-k}$ for every $k\ge1$.
The root-transfer inequality \eqref{eq:ordinary-root-transfer} now gives \eqref{eq:ordinary-integer-k}.  Finally $\delta\le2B$ and
$(27/64)(999/1000)^2>2/5$ give the stated lower bound for $R_2$.

For \eqref{eq:ordinary-continuous}, choose
$k=1+\lfloor\log Z/\log R_2\rfloor$.  Then
$Z/(2R_2^k)<1/2$ and $H_2^k\le H_2Z^{\lambda-1}$, whence
\[
 Z^{2-\lambda}<6(C/B)H_2=252B^2C^2<253B^2C^2.
\]
\end{proof}

\begin{corollary}\label{cor:ordinary-height-envelope}
Under the preceding hypotheses, suppose $D=|d|\ge C$ and put
$x=\log_{10}C$.  Let rational $h,r,u$ satisfy
\[
 \log_{10}H_2<x+h,\qquad
 \log_{10}R_2>2x+r,\qquad
 \log_{10}(253B^2)<u.
\]
Set $t=h-r$, $\beta=4u+8h-4r$ and $\gamma=2ur+\beta t$.
If $x>t$ and $2x+r>0$, then
\[
 \log_{10}D<7x+\beta+\frac\gamma{x-t}+\frac25.
\]
\end{corollary}

\begin{proof}
The assumptions give $R_2>1$ and
$2-\lambda>(x-t)/(2x+r)>0$.  Since $CD>2$,
$CD<2Z^2$.  Apply \eqref{eq:ordinary-continuous} and
$\log_{10}2<2/5$ to obtain
\[
 \log_{10}D<\frac25-x+
       \frac{2(2x+u)(2x+r)}{x-t}
       =7x+\beta+\frac\gamma{x-t}+\frac25.
\]
\end{proof}

\subsection{Three fixed entries at order two}

\begin{lemma}\label{lem:ordinary-three-normalized}
Let $a,b,s,c,d$ be distinct nonzero members of a $D(1)$-tuple, with
\[
 1\le A=|a|\le B=|b|\le S=|s|,\qquad
 C=|c|\ge4S,\qquad AC\ge1000.
\]
For $z^2=cd+1$ and $Z=|z|$ one has
\begin{equation}\label{eq:ordinary-three-M}
 1<\frac89\,608^2\frac{B^3S^8C^3}{Z}
       +\frac{3920}{81}\frac{B^2S^5Z}{A^3C^6}.
\end{equation}
If $D=|d|\ge C$, the closed interval
\begin{equation}\label{eq:ordinary-three-coarse-window}
 \frac{512}{81}608^4B^6S^{16}C^5\le D\le
 \frac{A^6C^{11}}{2\cdot28^4(10/81)^2B^4S^{10}}
\end{equation}
is impossible.
\end{lemma}

\begin{proof}
Use the common construction with $m=3,k=2$, write
$H=abs$, $\Delta=(a-b)(a-s)(b-s)$, and divide rows by
\[
 d_0=\Delta,\quad d_1=a(b-s),\quad
 d_2=b(a-s),\quad d_3=s(a-b).
\]
Let $R_{ij}=Q_{ij2}/d_i$.  These are integral polynomials.
To check this directly in \eqref{eq:ordinary-polynomial-summand},
in the row indexed by $0$ every incident difference retains an exponent at least one.
In row $i\ge1$ and column $j\ne i,0$, the factor $a_i$ is
explicit, and the opposite difference retains an exponent at least
one because $\delta_{i\ell}=0$ and $h_\ell\le1$ at its other
endpoint.  In column $i$, the factor $a_i^4$ supplies $a_i$, and
the opposite difference has exponent $4$.  The entries in the column indexed by $0$ also contain all the required row factors.  For $j\ge1$, every entry in column $j$ of $R$ remains divisible by
$a_j$, and
\[
 \det R=-226800H^3\Delta^{10}\ne0,
\]
by \eqref{eq:ordinary-residue-det} and
$\prod_{i=0}^3d_i=H\Delta^2$.

We now bound the coefficients after these row divisions.  Put
\[
 \rho_j=\min\left(1,\min_{\substack{\ell\ne j\\\ell\ge1}}
               \frac{|a_\ell-a_j|}{A_\ell}\right),\qquad
 \sigma_j=\frac{|H|\rho_j}{A_j},
\]
\[
 F=1+\frac5{4AC},\quad
 p=\sqrt{\frac{F}{1-1/(AC)}},\quad
 K_0=\frac{1024}{27},\quad G=16K_0F.
\]
At the pole $\alpha_j$ integrate on the circle of radius
$\sigma_j/4$.  It lies in the disk $|w|<|N|$ and misses every other
pole.  The product of its four pole distances is at least
\[
 \frac14\left(\frac34\right)^3
 \sigma_j\prod_{\ell\ne j}|\alpha_j-\alpha_\ell|.
\]
Retaining the factor $|w-\alpha_i|$ in the denominator gives
\[
 |Q_{ij2}|\le p\tau_{ij}M_j^2,\qquad
 M_j=K_0FC|\Delta|^2
       \frac{A_j^3}{\rho_j\prod_{\ell\ne j,\,\ell\ge1}|a_\ell-a_j|},
\]
where
\[
 \tau_{jj}=1,\qquad \tau_{0j}\le\rho_j/3,
 \qquad \tau_{ij}\le
       \frac{A_i\rho_j}{3|a_i-a_j|}\quad(i\ne j,\ i,j\ge1).
\]
For $i\ne j$ and $w$ on this circle, the ratio
$(\sigma_j/4)/|w-\alpha_i|$ is at most $\sigma_j/(3|\alpha_i-\alpha_j|)$.  The numerator
divided by $|\theta_j|$ is at most $pF^2$, which accounts for all
remaining factors.

The ordered bounds $|a-b|\le2B$ and
$|a-s|,|b-s|\le2S$ imply
\begin{equation}\label{eq:ordinary-three-base-bound}
 M_j\le16K_0FCB^2S^5=GCB^2S^5.
\end{equation}
For the term $1$ in $1/\rho_j$, the remaining product is
$A_j^3|a_j-a_\ell||a_j-a_r||a_\ell-a_r|^2$, at most
$16B^2S^5$.  For either other term in $1/\rho_j$, one incident
difference cancels and the bound is $8B^2S^5$.

Use $M_j^2\le(GCB^2S^5)M_j$ in the weighted row sum.  For row $i\ge1$ its
diagonal contribution is
\[
 \frac{M_i}{A_i|d_i|}
 =K_0FC\frac{A_i|\Delta|}{\rho_i}\le8K_0FCBS^3.
\]
The three alternatives in $1/\rho_i$ give bounds $8BS^3,4BS^3,4BS^3$
for the remaining product.  With $r$ the remaining index, each
off-diagonal contribution satisfies
\[
 \frac{\tau_{ij}M_j}{A_j|d_i|}
 \le\frac{K_0FC}{3}A_j^2|a_i-a_r|^2
 \le\frac43K_0FCBS^3.
\]
For the row indexed by $0$ the sum is at most
\[
 \frac{K_0FC}{3}\sum_j A_j^2|a_\ell-a_r|
 \le2K_0FCBS^2.
\]
Therefore every row has weighted sum at most
$(32/3)K_0FCBS^3$.  Multiplying by the common bound $GCB^2S^5$ from \eqref{eq:ordinary-three-base-bound} gives
\begin{equation}\label{eq:ordinary-three-column-bound}
 \sum_{j=1}^3\frac{|R_{ij}|}{A_j}
 \le\frac23pG^2C^2B^3S^8.
\end{equation}

For the residual, the common beta integral at $m=3,k=2$ can be written
\[
 \left|\sum_{j=0}^3Q_{ij2}\theta_j\right|
 \le\frac{10}{81}\eta^{1/2}
       \left(\frac{27|\Delta|^2}{64|H|C^3\eta^3}\right)^2,
 \qquad\eta=1-1/(AC).
\]
The half-integer gamma identities give the factor $10/81$ directly.
For $i=0,1,2,3$ the respective bounds for
$|\Delta|^2/(|H d_i|)$ are
\[
 \frac{8S}{A},\quad\frac{32BS^2}{A^2},\quad
 \frac{32S^2}{A},\quad\frac{32S^2}{A}.
\]
All are at most $32BS^2/A^2$, and
$|\Delta|^2/|H|\le64BS^3/A$.  After row division the residual is
therefore at most
\begin{equation}\label{eq:ordinary-three-residual}
 \frac{10}{81}\frac{\eta^{1/2}}{2AS}
       \left(\frac{27BS^3}{AC^3\eta^3}\right)^2.
\end{equation}
Since $a_j$ divides $R_{ij}$ for $j\ge1$ and $\det R\ne0$, the preceding integral-form argument applies, now using
\eqref{eq:ordinary-three-column-bound} and
\eqref{eq:ordinary-three-residual}.  Finally
\[
 G\le\frac{16384}{27}\frac{4005}{4000}<608,
 \qquad
 p^2e(S)^2\le
 \frac{25\cdot1001\cdot4005}{64000\cdot999}<\frac{16}{9},
\]
\[
 27\eta^{-3}\le27(1000/999)^3<28,
 \qquad\eta^{1/2}<1.
\]
These strict bounds yield \eqref{eq:ordinary-three-M}.
If \eqref{eq:ordinary-three-coarse-window} held, then
$CD/2<Z^2<2CD$ would make each term in
\eqref{eq:ordinary-three-M} less than $1/2$, a contradiction.
\end{proof}

\section{Finite integral approximation forms}

The interval arguments below begin with a structural lower bound and end with the two-entry estimates at every order. The finite constructions in this section supply the intervening forbidden intervals. We give finite formulas for every row used here and prove the required divisibility in the polynomial ring over $\mathbb Z$ before evaluating the tuple entries.

\subsection{From polynomial rows to inequalities}

Let $x_1,\ldots,x_h,c,d$ be distinct nonzero entries of a $D(1)$-tuple
in the full ring of integers $\mathcal O_K$ of an imaginary quadratic
field. Put $X_j=|x_j|$, $C=|c|$, $D=|d|$,
$X=\min_jX_j$, $Y=\max_jX_j$, and choose $z\in\mathcal O_K$ with
$z^2=cd+1$. Suppose
\begin{equation}\label{eq:ff-basicguards}
 X\ge1,\qquad XC\ge1000,\qquad C\ge4Y,\qquad D\ge C.
\end{equation}
Then $C^2\ge4000$, so $z\ne0$. Write $Z=|z|$ and let
$\theta_x=(1+1/(xc))^{1/2}$ be the binomial branch.
If $r_x^2=xc+1$ and $v_x^2=xd+1$, signs can be chosen, independently
for the fixed entries but consistently across the rows, so that
\begin{equation}\label{eq:ff-error}
 \left|\theta_x-\frac{\pm r_xv_x}{xz}\right|
 \le \frac{C}{|x|Z^2}
       \sqrt{1+\frac1{|x|C}}\left(1+\frac{|x|}{C}\right)
 <\frac{eC}{|x|Z^2}.
\end{equation}
The first inequality is the root estimate in the proof of
Lemma~\ref{lem:ordinary-residue-construction}. Here $e=3/2$ is valid under
\eqref{eq:ff-basicguards}, and $e=101/100$ is valid if $C\ge200Y$, since
\[
 \frac{1001}{1000}\left(\frac54\right)^2<\left(\frac32\right)^2,
 \qquad
 \frac{1001}{1000}\left(\frac{201}{200}\right)^2
 <\left(\frac{101}{100}\right)^2.
\]
Neither $r_x$ nor $v_x$ is divided out.

Suppose a polynomial row has entries $(Q_0,Q_{x_1},\ldots,Q_{x_h})$
with integral coefficients and $Q_x=xP_x$, where $P_x$ is also integral.
Its arithmetic form
\[
 \Lambda=Q_0z+\sum_xP_x(\pm r_xv_x)\in\mathcal O_K
\]
satisfies
\begin{equation}\label{eq:ff-master}
 |\Lambda|\le Z\left|Q_0+\sum_xQ_x\theta_x\right|
                 +\frac{eC}{Z}\sum_x|P_x|.
\end{equation}
A matrix of such rows with rank $h+1$ supplies a nonzero $\Lambda$,
since its common evaluation vector has first coordinate $z\ne0$.
For this row $|\Lambda|\ge1$, because the norm of a nonzero algebraic
integer in $K$ is a positive rational integer. All the estimates below
enlarge a coefficient strictly and therefore give a strict inequality
$1<U/Z+VZ$.

We will use the following conversion repeatedly. Set $\varepsilon=1/4000$.
Since $D\ge C$ and $C^2\ge4000$,
\begin{equation*}
 (1-\varepsilon)CD\le Z^2\le(1+\varepsilon)CD.
\end{equation*}
Consequently $1<U/Z+VZ$ excludes the closed interval
\begin{equation}\label{eq:ff-window}
 \frac{4U^2}{C(1-\varepsilon)}\le D\le
 \frac{1}{4V^2C(1+\varepsilon)}.
\end{equation}
An interval with its lower endpoint above its upper endpoint is simply
empty. Each endpoint in a nonempty interval makes the corresponding
summand at most $1/2$, which explains both closed endpoints.

\subsection{Five families with two fixed entries}

Let $x,y,c,d$ be distinct nonzero tuple entries, with
\begin{equation}\label{eq:ff-twoguards}
 1\le X=|x|\le Y=|y|,\qquad XC\ge1000,
 \qquad C\ge200Y,\qquad D\ge C.
\end{equation}
The parameters for the five families are in Table~\ref{tab:ff-parameters}.
Here $P=m+2k$. In this subsection, $c_{\mathrm{det}}$ and $J$ denote constants in the determinant formulas. The column $R_*$ gives a uniform residual bound. Smaller bounds depending on $|x-y|$ are derived below.

\begin{table}[ht]\centering\small
\begin{tabular}{c|rrr|r|r|c|c}
 family&$k$&$m$&$g$&$c_{\mathrm{det}}$&$J$&$\beta$&$R_*$\\ \hline
 $N_2$&2&2&2&$-1800$&$-15/64$&$5/1024$&$5/4$\\
 $N_3$&3&3&4&$235200$&$105/256$&$35/65536$&$35/32$\\
 $A_3$&3&2&4&$-156800$&$-35/256$&$35/32768$&$35/16$\\
 $A_4$&4&3&2&$23284800$&$3465/16384$&$63/524288$&$63/8$\\
 $A_5$&5&4&4&$-3995671680$&$-27027/65536$&$495/33554432$&$495/64$
\end{tabular}
\caption{Parameters and exact constants of the finite two-anchor rows.}
\label{tab:ff-parameters}
\end{table}

Here is an explicit finite construction of all their coefficients.
Put $p=(m,k,k)$ and $\delta_{ij}=1$ if $i=j$, and $0$ otherwise.
For $i,j\in\{0,1,2\}$, sum over nonnegative triples $h=(h_0,h_1,h_2)$
with $h_0+h_1+h_2=p_j+\delta_{ij}-1$. Set
\[
 b_{ijh}=4^k\binom{k+1/2}{h_j}
                 \prod_{\ell\ne j}\binom{-p_\ell-\delta_{i\ell}}{h_\ell}
\]
and define $F_{ij}$ as the sum of $b_{ijh}T_{ijh}$, where
\begin{align}
 T_{i0h}&=(-1)^{2k+\delta_{i1}+\delta_{i2}+h_1+h_2}
 x^{\delta_{i1}+h_1}y^{\delta_{i2}+h_2}
 c^{k-h_0}(x-y)^{2k},\notag\\
 T_{i1h}&=(-1)^{k+\delta_{i2}+h_2}
 x^m y^{\delta_{i2}+h_2}(xc+1)^{k-h_1}
 (x-y)^{k-\delta_{i2}-h_2},\label{eq:ff-generator}\\
 T_{i2h}&=y^m x^{\delta_{i1}+h_1}(yc+1)^{k-h_2}
 (x-y)^{k-\delta_{i1}-h_1}.\notag
\end{align}
Define $\widetilde Q$ by dividing row $0$ of $F$ by $x-y$, row $1$ by $x$, and row $2$ by $y$. The integral matrix used below is $Q=\widetilde Q/g$.

All exponents in \eqref{eq:ff-generator} are nonnegative. The row
divisions hold term by term. For example, in row $0$, column $1$,
$h_2\le k-1$, so its difference exponent is positive. The same argument applies to column $2$. Row $1$ has a factor $x$ in every column, and row
$2$ a factor $y$. Moreover, after these divisions, column $1$ still
has a factor $x$ and column $2$ still has a factor $y$, since $m\ge2$.
Finally the numbers $4^k\binom{k+1/2}{h}$ for $0\le h\le k$ are
\[
\begin{array}{c|l}
 k&4^k\binom{k+1/2}{h},\quad h=0,\ldots,k\\ \hline
 2&16,40,30\\
 3&64,224,280,140\\
 4&256,1152,2016,1680,630\\
 5&1024,5632,12672,14784,9240,2772.
\end{array}
\]
They are divisible by the listed $g$, whereas all the other binomial
coefficients in $b_{ijh}$ are integers. Thus $Q$ has integral coefficients, and its columns $1$ and $2$ are divisible by $x$ and $y$, respectively. These conclusions hold
before any specialization of $x,y,c$.

We next determine the order of the residual and prove that the matrices have rank three. Put
$\alpha=(0,y,x)$, $N=xyc$. The residue sum
\begin{equation}\label{eq:ff-residue}
 I_i=\frac1{2\pi i}\int_\Gamma
 \frac{(1+w/N)^{k+1/2}\,dw}
 {(w-\alpha_i)\prod_{\ell=0}^2(w-\alpha_\ell)^{p_\ell}}
\end{equation}
uses a positively oriented contour around the three poles, excluding
the cut $w=-N(1+t)$, $t\ge0$. Their distance from this cut is positive
because $|\alpha_j/N|\le1/(XC)<1$ for $j=1,2$.
Leibniz's rule for the three pole residues gives exactly
\eqref{eq:ff-generator}, and hence
\[
 4^k c^k x^{m+k}y^{m+k}(x-y)^{2k} I_i
 =F_{i0}+F_{i1}\theta_x+F_{i2}\theta_y.
\]
The estimate below shows that these residuals are $O(c^{-(P-k)})$.
Replacing the column indexed by $0$ of $\widetilde Q$ by its residual therefore shows
that its polynomial determinant has degree at most
$2k-(P-k)=k-m$ in $c$. The entry in row $i$ and column $0$ is divisible
by $c^{k-m+1-\delta_{i0}}$, directly from the finite sum, so
the determinant is divisible by $c^{k-m}$. The coefficient of $c^{k-m}$ in $\det\widetilde Q$ is therefore
$4^k\binom{k+1/2}{m}(x-y)^{2k-1}$ times the determinant of the
constant root block in rows $1,2$.

To compute this block, put
$\lambda=k+1/2-m$ and
\begin{align*}
 A_{k,\lambda}(x,y)&=\sum_{h=0}^k
   \binom{-k}{h}\binom{\lambda}{k-h}x^h(x-y)^{k-h},\\
 B_{k,\lambda}(x,y)&=\sum_{h=0}^{k-1}
   \binom{-k-1}{h}\binom{\lambda}{k-1-h}x^h(x-y)^{k-1-h}.
\end{align*}
Vandermonde's convolution over the other two indices of
\eqref{eq:ff-generator} gives that block as
\[
 4^k\begin{pmatrix}
 x^{m-1}A_{k,\lambda}(y,x)&y^mB_{k,\lambda}(x,y)\\
 x^mB_{k,\lambda}(y,x)&y^{m-1}A_{k,\lambda}(x,y)
 \end{pmatrix}.
\]
For the five parameter pairs in Table~\ref{tab:ff-parameters},
\begin{equation}\label{eq:ff-binomialdet}
 A_{k,\lambda}(x,y)A_{k,\lambda}(y,x)
 -xyB_{k,\lambda}(x,y)B_{k,\lambda}(y,x)
 =J(x-y)^{2k},
\end{equation}
where $J=(-1)^k\binom\lambda k\binom{\lambda-k}k$ is printed in the
table. To compare the coefficients in this identity, put
\[
 a_j=(-1)^j\binom\lambda j\binom{\lambda-k-j}{k-j},\quad
 b_j=(-1)^j\binom\lambda j\binom{\lambda-k-1-j}{k-1-j},
\]
with $a_j=0$ outside $0\le j\le k$ and $b_j=0$ outside
$0\le j<k$. Expanding the preceding sums by Vandermonde gives,
for $r=0,\ldots,2k$,
\[
 \sum_{j=0}^k(a_ja_{k-r+j}-b_jb_{k-r+j})
 =J(-1)^r\binom{2k}{r}.
\]
Substitution of $(k,\lambda)=(2,1/2),(3,1/2),(3,3/2),
(4,3/2),(5,3/2)$ verifies every coefficient of
\eqref{eq:ff-binomialdet}. It follows that
\begin{equation*}
 \det \widetilde Q=c_{\mathrm{det}}(xy)^{m-1}c^{k-m}(x-y)^{4k-1},\qquad
 c_{\mathrm{det}}=4^k\binom{k+1/2}{m}16^kJ.
\end{equation*}
In particular $Q$ has rank three at every admissible evaluation.

To bound the residual, deform \eqref{eq:ff-residue} to its cut.
The large-circle contribution vanishes, since $P>k+1/2$, and the
difference of the two boundary values of the numerator has modulus
$2t^{k+1/2}$.
With $\eta=1-1/(XC)$ one obtains
\[
 |I_i|\le \frac{|N|^{-P}}\pi
       \int_0^\infty\frac{t^{k+1/2}\,dt}{(t+\eta)^{P+1}}
 =\beta\eta^{-(P-k-1/2)}|N|^{-P},
\]
where substitution $t=\eta u/(1-u)$ and the elementary half-integer
beta integral give
\begin{equation*}
 \beta=\frac{(2k+2)!(2P-2k-2)!}
 {4^P(k+1)!(P-k-1)!P!}.
\end{equation*}
The five values are in the table. Their largest exponent
$P-k-1/2$ is $17/2$, and
$(1000/999)^{17}<(101/100)^2$. Hence $\eta^{-(P-k-1/2)}$ is
strictly below $101/100$.

Put $u=X/Y$ and $\delta=|x-y|/Y$. After the row divisions and division
by $g$, each residual is strictly less than
\begin{equation}\label{eq:ff-point-residual}
 R(x,y)\frac{Y^k}{X^{k+1}C^{P-k}},\qquad
 R(x,y)=\frac{101\,4^k\beta}{100g}
             \max(\delta^{2k},u\delta^{2k-1}).
\end{equation}
Indeed, the three extra row factors are respectively
$u\delta^{2k-1},\delta^{2k},u\delta^{2k}$, and $u\le1$.
Equivalently the coefficient in front of the maximum is
$R_*101/(200\,2^{2k})$, since $R_*=2\,16^k\beta/g$.
If only $Y\ge qX$ with $q\ge1$ is known, one may use
\begin{equation}\label{eq:ff-ratio-residual}
 R(q)=\frac{101\,4^k\beta}{100g}(1+1/q)^{2k}.
\end{equation}
In particular $R(1)=101R_*/200<R_*$, so every listed $R_*$ is valid.

For completeness we give all data needed for the coefficient bounds.
Write $P_{i1}=\widetilde Q_{i1}/x$ and $P_{i2}=\widetilde Q_{i2}/y$ before division by $g$.
Every monomial $x^r y^s c^\ell$ in these polynomials satisfies
\begin{equation}\label{eq:ff-weight}
 r,s\ge0,\qquad 0\le\ell\le k,\qquad r+s=\ell+P-k-2.
\end{equation}
This follows term by term from \eqref{eq:ff-generator}.
Let $L_{i\ell}$ be the sum of absolute integer coefficients of
$c^\ell$ in $P_{i1},P_{i2}$ together. The following table is obtained
by expanding the displayed finite sums. Rows $1$ and $2$ have the same sums.
\[
\begin{array}{c|c|l}
 \text{family}&i&(L_{i\ell})_{\ell=0}^k\\ \hline
N_2&0&96,352,256\\
 &1,2&100,400,480\\
N_3&0&1344,9408,17024,8960\\
 &1,2&1368,8384,23168,16384\\
A_3&0&800,4544,10368,6144\\
 &1,2&1288,4256,9856,8960\\
A_4&0&10944,86784,350208,482304,215040\\
 &1,2&18128,80512,276096,561920,327680\\
A_5&0&156000,1601280,9707520,22379520,22087680,7864320\\
 &1,2&262152,1463264,6707712,22437888,28452864,12300288
\end{array}
\]
If $M$ is any constant strictly larger than
\begin{equation}\label{eq:ff-uniformM}
 \frac{101}{100g}\max_i\sum_{\ell=0}^k
                       L_{i\ell}1000^{\ell-k},
\end{equation}
then \eqref{eq:ff-master} proves
\begin{equation*}
 1<\frac{M Y^{P-2}C^{k+1}}Z
       +\frac{R Y^kZ}{X^{k+1}C^{P-k}},
\end{equation*}
where $R$ can be chosen from \eqref{eq:ff-point-residual},
\eqref{eq:ff-ratio-residual}, or the table.
To see \eqref{eq:ff-uniformM}, divide a monomial in
\eqref{eq:ff-weight} by $Y^{P-2}C^k$. Its modulus is
$(X/Y)^r(YC)^{\ell-k}\le1000^{\ell-k}$.
Valid integer choices of $M$ in the table order are
\[
 243,\quad4143,\quad2265,\quad165763,\quad3113013.
\]
The endpoints supplied by
\eqref{eq:ff-window}, before division by $1\mp\varepsilon$, are
\begin{equation*}
 l=4M^2Y^{2P-4}C^{2k+1},\qquad
 u=\frac{X^{2k+2}C^{2P-2k-1}}{4R^2Y^{2k}}.
\end{equation*}

We shall also use two refinements of the coefficient bound. If $Y\ge qX$, replace the summand corresponding to
$n x^ry^sc^\ell$ in \eqref{eq:ff-uniformM} by
$|n|q^{-r}1000^{\ell-k}$. The same proof applies.
For fixed values of $x$ and $y$, first combine the coefficients algebraically and
write
$P_{ij}(x,y,c)=\sum_{\ell=0}^kp_{ij\ell}(x,y)c^\ell$.
For any established lower bound $C\ge C_0>0$, it suffices to take
\begin{equation}\label{eq:ff-fixedM}
 M>\frac{101}{100gY^{P-2}}
       \max_i\sum_{j=1}^2\sum_{\ell=0}^k
                    |p_{ij\ell}(x,y)|C_0^{\ell-k}.
\end{equation}
If $x$ is a unit, \eqref{eq:ff-twoguards} permits
$C_0=\max(1000,200Y)$.
Every modulus in \eqref{eq:ff-fixedM} is the square root of the exact
integer norm of an explicitly specified algebraic integer. Each coefficient modulus can therefore be bounded above by a rational number by comparing the square of that number with the corresponding integral norm. In the field basis $\omega^2-\tau_K\omega+\nu_K=0$, the required norm
is $N(r+s\omega)=r^2+\tau_Krs+\nu_Ks^2$.

\subsection[Quadratic rows and their error bounds]{Quadratic rows with error of order $c^{-5}$}

Now let $a,b,s,c,d$ be distinct nonzero tuple entries, with ordered
moduli $1\le A\le B\le S\le C\le D$, and assume
\begin{equation*}
 C\ge200S,\qquad AC\ge1000,\qquad C\ge1000.
\end{equation*}
For an ordered pair $(u,v)$ chosen from $a,b,s$, define rows $R_{uv},T_{uv}$ as follows. Their other root columns vanish, $Q_t=tP_t$, and
$P_t=P_{t,0}+P_{t,1}c+P_{t,2}c^2$. The complete coefficient data are
\begin{align*}
 R:\quad P_{u,0}&=7u^2-35uv,&P_{v,0}&=35uv-7v^2,\\
 P_{u,1}&=64u^3-140u^2v,&P_{v,1}&=140uv^2-64v^3,\\
 P_{u,2}&=72u^4-120u^3v,&P_{v,2}&=120uv^3-72v^4;\\[2mm]
 T:\quad P_{u,0}&=-15u^3+56u^2v-105uv^2,&
 P_{v,0}&=70uv^2-6v^3,\\
 P_{u,1}&=-180u^4+512u^3v-420u^2v^2,&
 P_{v,1}&=100uv^3-12v^4,\\
 P_{u,2}&=-240u^5+576u^4v-360u^3v^2,&
 P_{v,2}&=24v^5.
\end{align*}
In both rows put
\begin{equation*}
 Q_0=-\sum_{t=u,v}\left[
 tP_{t,2}c^2+(tP_{t,1}+P_{t,2}/2)c
              +tP_{t,0}+P_{t,1}/2-P_{t,2}/(8t)\right].
\end{equation*}
This is an integral polynomial because $P_{t,1}$ has even coefficients and
$P_{t,2}$ is divisible by $8t$.

With $b_n=\binom{1/2}{n}$, the coefficient of $c^{-n}$ after removing
the polynomial part is
\begin{equation*}
 \sum_{t=u,v}\left(b_nt^{1-n}P_{t,0}
       +b_{n+1}t^{-n}P_{t,1}+b_{n+2}t^{-n-1}P_{t,2}\right).
\end{equation*}
Substituting the displayed polynomials gives zero for $n=1,2,3,4$.
The coefficients at $c^{-5}$ are respectively
\begin{align}
 E_R(u,v)&=\frac5{128}(u^{-2}-v^{-2})
              +\frac5{256}(uv^{-3}-vu^{-3}),\label{eq:ff-leading}\\
 E_T(u,v)&=-\frac{75}{128u}+\frac{5v}{16u^2}
 -\frac{15v^2}{256u^3}-\frac{35u}{256v^2}+\frac{15}{32v}.\notag
\end{align}
These identities and the bound $|b_n|\le21/1024$ for $n\ge6$ control the remaining series.

Use the four rows
\begin{equation}\label{eq:ff-prows}
 R_{ab},\quad H=R_{ab}+R_{bs}-R_{as},\quad
 T_{ab},\quad V=T_{as}-T_{bs}.
\end{equation}
In columns $(Q_0,Q_a,Q_b,Q_s)$, their determinant is
\begin{equation*}
 -78750a^2b^2s^2(a-b)^5(a-s)(b-s)F_s,
\end{equation*}
where
\[
 F_s=7(2s-a-b)+6c[3s^2-(a+b)s-ab].
\]
Replacing column zero by its $O(c^{-5})$ residual shows that the
polynomial determinant has degree at most $2+2+2-5=1$ in $c$.
Expanding its constant and linear coefficients from the displayed rows
gives the identity.

Suppose that $F_s=0$. If $3s^2-(a+b)s-ab\ne0$, this algebraic integer
has modulus at least one, and
\[
 6C\le7|2s-a-b|\le28S,
\]
contrary to $C\ge200S$. Otherwise $a+b=2s$ and $ab=s^2$, so
$a=b=s$, again a contradiction. Thus the matrix has rank four.

Put $H_* =\max(BS^3,B^3S^2)$. The coefficient sums by increasing
$c$ degree in the root columns after dividing by their anchors are
\begin{equation*}
 \begin{array}{c|rrr}
  R_{ab}&84&408&384\\
  H&210&840&720\\
  T_{ab}&252&1224&1200\\
  V&492&2424&2352.
 \end{array}
\end{equation*}
Every anchor monomial in these columns is at most $H_*$.
For $H$, the three coefficients at degree $\ell$ are
$k_\ell a^{\ell+1}(s-b)$,
$k_\ell b^{\ell+1}(a-s)$,
$k_\ell s^{\ell+1}(b-a)$, where $(k_0,k_1,k_2)=(35,140,120)$.
For $V$, the coefficient obtained after dividing the last column by $s$ is
$P_s=70(a-b)s^2+100(a-b)s^3c$. The monomials in $P_a$ and $P_b$
are bounded by $B^3S^2$.
The remaining pair rows have monomials at most $B^4$ and $B^5$,
respectively. These bounds use only $1\le A\le B\le S$ and remain valid when some moduli are equal.
Thus the second term in \eqref{eq:ff-master} is bounded above by
\[
 \frac{101}{100}\max_i(L_{i2}+L_{i1}/1000+L_{i0}/1000^2)<2400
\]
times $H_*C^3/Z$.

We now retain the ratio $B/A$ in the residual bound.
Set $\tau=21/(1024\cdot999)$, and define the polynomials
\[
\begin{aligned}
 256L_{R_{ab}}(x)&=5x+10x^2+10x^4+5x^5,\\
 W_{R_{ab}}(x)&=295x+143x^2+143x^4+295x^5,\\
 256L_H(x)&=10(x+x^4+x^5),\\
 W_H(x)&=590(x+x^4+x^5),\\
 256L_{T_{ab}}(x)&=15+80x+150x^2+120x^3+35x^4,\\
 W_{T_{ab}}(x)&=885+1144x+435x^2+42x^3+170x^4,\\
 256L_V(x)&=15+80x+150x^2+280x^3+35x^4,\\
 W_V(x)&=885+1144x+435x^2+2634x^3+170x^4.
\end{aligned}
\]
For $q\ge1$ put
\begin{equation}\label{eq:ff-adaptiveres}
 \mathcal R(q)=\max_i\{L_i(1/q)+\tau W_i(1/q)\}.
\end{equation}
To verify the table, combine the leading coefficients
\eqref{eq:ff-leading} according to \eqref{eq:ff-prows} before taking
absolute values. For a Laurent monomial $a^ib^js^k$, division of its
modulus by $S^2/A^3$ gives, on setting $u=A/B$ and $v=B/S$,
\[
 A^{i+j+k+1}u^{2-j-k}v^{2-k}.
\]
Every monomial from those combined leading coefficients satisfies
$i+j+k+1\le0$, $2-j-k\ge0$, $2-k\ge0$.
Summing its absolute coefficient times $x^{2-j-k}$ gives $L_i(x)$.
For the tail, a monomial from the coefficient $P_{t,\ell}$ has the exponent of the corresponding fixed entry decreased by $4+\ell$. The same three inequalities
hold. Use the divided root-column polynomials of \eqref{eq:ff-prows}, with
the combinations defining $H$ and $V$ already performed. Sum the absolute
values of their integer coefficients with weight $x^{2-j-k}$, before
multiplication by the binomial coefficients. This gives $W_i(x)$.
Each later inverse power brings a further factor at most $1/(AC)$.
Since $|b_n|\le21/1024$ for $n\ge6$ and
$\sum_{r\ge1}(AC)^{-r}\le1/999$, the whole tail contributes at most
$\tau W_i(1/q)$. This proves \eqref{eq:ff-adaptiveres}.

For every $q\ge1$ with $B\ge qA$, we conclude that
\begin{equation}\label{eq:ff-P}
 1<\frac{2400H_*C^3}{Z}
       +\frac{\mathcal R(q)S^2Z}{A^3C^5}.
\end{equation}
For $0\le x\le1$, the displayed polynomials satisfy
$L_i(x)\le L_V(x)$ and $W_i(x)\le W_V(x)$ for every row $i$.
These comparisons follow from $x^5\le x^4\le x^3$.
Consequently $\mathcal R(q)=L_V(1/q)+\tau W_V(1/q)$ for all $q\ge1$,
and this function is nonincreasing.
In particular,
\[
 \mathcal R(1)=\frac{35}{16}+\tau\,5268<3.
\]
Equation~\eqref{eq:ff-P}, together with \eqref{eq:ff-window}, gives
its closed forbidden interval.

\subsection{Five linear rows on four fixed entries}

Let $a,b,s,t,c,d$ be six distinct nonzero tuple entries, with
$1\le A\le B\le S\le T\le C\le D$, and suppose
\begin{equation*}
 C\ge4T,\qquad AC\ge1000.
\end{equation*}
For an ordered pair $(u,v)$ chosen from $a,b,s,t$ define a row $\mathcal L_{uv}$ by
\begin{align*}
 P_u&=u^2-5uv+(4u^3-8u^2v)c,\\
 P_v&=5uv-v^2+(8uv^2-4v^3)c,\\
 Q_u&=uP_u,\qquad Q_v=vP_v,\qquad
 Q_0=-(u-v)^3[3+4(u+v)c],
\end{align*}
and put the other root columns equal to zero.
For $P_x=P_{x,0}+P_{x,1}c$, the coefficient of $c^{-n}$ in the
residual is
\[
 b_n\sum_{x=u,v}x^{1-n}P_{x,0}
 +b_{n+1}\sum_{x=u,v}x^{-n}P_{x,1}.
\]
The two sums for $n=1,2,3$ are, respectively,
\[
 \begin{array}{c|cc}
 1&u^2-v^2&4(u^2-v^2)\\
 2&6(u-v)&12(u-v)\\
 3&5(u/v-v/u)&8(u/v-v/u).
 \end{array}
\]
After multiplication by $b_n$ and $b_{n+1}$, the two entries in each row sum to zero. The coefficient $Q_0$ cancels the two nonnegative powers.
The five rows
$\mathcal L_{ab},\mathcal L_{as},\mathcal L_{bs},
\mathcal L_{at},\mathcal L_{bt}$, in columns $(Q_0,Q_a,Q_b,Q_s,Q_t)$,
have determinant
\begin{equation}\label{eq:ff-qdet}
 180a^2b^2s^2t^2(a-b)^2(a-s)(a-t)(b-s)(b-t)(s-t).
\end{equation}
Indeed, the column indexed by $0$, after replacement by the residual, is $O(c^{-4})$ and the remaining four columns have degree at most one, so their polynomial determinant is
constant in $c$. Expanding its constant term from the displayed rows
gives \eqref{eq:ff-qdet}. Every factor is nonzero, giving rank five
without an additional assumption on the moduli.

The sum of absolute integer coefficients of the $P$ columns in each
row is $36$. Each coefficient monomial, including its power of $c$,
is bounded by $T^3C$, so $\sum_x|P_x|\le36T^3C$.
After the cancellations, the first residual monomials are
$u^{-1},v/u^2,u/v^2,v^{-1}$, each bounded by $T/A^2$.
Since $|b_n|\le5/128$ for $n\ge4$, the entire convergent residual is
at most
\[
 \frac{36\cdot5}{128(1-1/(AC))}\frac{T}{A^2C^4}
 \le\frac{625}{444}\frac{T}{A^2C^4}
 <\frac{3T}{2A^2C^4}.
\]
Apply \eqref{eq:ff-master} with $e=3/2$ to a nonzero form and enlarge
$54$ to $55$. This gives
\begin{equation*}
 1<\frac{55T^3C^2}{Z}+\frac{3TZ}{2A^2C^4}.
\end{equation*}
Its endpoints before applying the factors in \eqref{eq:ff-window} are
$l=4\cdot55^2T^6C^3$ and $u=A^4C^7/(9T^2)$, to be multiplied by $1/(1-\varepsilon)$ and $1/(1+\varepsilon)$, respectively.

\subsection{The interval calculation}
\label{sec:common-intervals}

Let $a,b,s,t,c,d$ be six distinct nonzero tuple entries, ordered by
$1\le A\le B\le S\le T\le C\le D$. Assume
\[
 C>\max(1000,200S,4T).
\]
These conditions imply $AC>1000$ and $C^2>4000$, and supply the
hypotheses of all the approximation constructions used below.
As in \eqref{eq:ff-window}, put
\[
 \varepsilon=\frac1{4000},\qquad h=\frac{4001}{3999},\qquad
 M_m=\frac89\,608^2,\qquad R_m=\frac{3920}{81}.
\]
Write $M_2,M_n,M_3,M_4,M_5$ for valid coefficient bounds for
$N_2,N_3,A_3,A_4,A_5$ on $a,b$, and write
$R_2,R_n,R_3,R_4,R_5$ for the corresponding residual bounds.
For $N_3$ on $b,s$, use valid bounds $M_{bs},R_{bs}$.
The uniform choice $M_{bs}=4143$ is sufficient.
Write $R_p$ for an upper bound for $\mathcal R(q)$, where
$q\ge1$ and $B\ge qA$.
The rows denoted by $P$ and $Q$ are those of
\eqref{eq:ff-prows} and \eqref{eq:ff-qdet}, respectively,
and $M$ denotes \eqref{eq:ordinary-three-M}.

Equation~\eqref{eq:ff-window} shows that $D$ avoids every closed
interval $[l/(1-\varepsilon),u/(1+\varepsilon)]$ in
Table~\ref{tab:pr-windows}. Both terms in the maximum for $P$ must
be checked. The ordinary family $O_j$ is available for every positive
integer $j$ by \eqref{eq:ordinary-integer-k}.
\begin{table}[htbp]
\centering\small
\caption{Forbidden-interval endpoints.}\label{tab:pr-windows}
\setlength{\tabcolsep}{4pt}
\begin{tabular}{c|c|c}
Family&\(l\)&\(u\)\\\hline
\(Q\)&\(4\cdot55^2T^6C^3\)&\(A^4C^7/(9T^2)\)\\
\(N_2\)&\(4M_2^2B^8C^5\)&\(A^6C^7/(4R_2^2B^4)\)\\
\(P\)&\(4\cdot2400^2\max(BS^3,B^3S^2)^2C^5\)&\(A^6C^9/(4R_p^2S^4)\)\\
\(A_3\)&\(4M_3^2B^{12}C^7\)&\(A^8C^9/(4R_3^2B^6)\)\\
\(M\)&\(4M_m^2B^6S^{16}C^5\)&\(A^6C^{11}/(4R_m^2B^4S^{10})\)\\
\(N_3\)&\(4M_n^2B^{14}C^7\)&\(A^8C^{11}/(4R_n^2B^6)\)\\
\(N_{3,bs}\)&\(4M_{bs}^2S^{14}C^7\)&\(B^8C^{11}/(4R_{bs}^2S^6)\)\\
\(A_4\)&\(4M_4^2B^{18}C^9\)&\(A^{10}C^{13}/(4R_4^2B^8)\)\\
\(A_5\)&\(4M_5^2B^{24}C^{11}\)&\(A^{12}C^{17}/(4R_5^2B^{10})\)\\
\(O_j\)&\(36(42B^3C)^{2j}C/B^2\)&\((2AC^2/(5B))^{2j}/C\)
\end{tabular}
\end{table}

We use the following elementary observation. Suppose that $D$ is
above the lower endpoint of the first forbidden interval, and that
each upper endpoint is at least the next lower endpoint. Avoidance
of the first interval puts $D$ above its upper endpoint, hence above
the next lower endpoint. Induction applies even when an interval is
empty. A chain whose lower endpoints tend to infinity is therefore
impossible.
For the ordinary intervals, put
\[
 H_0=42B^3C,\qquad R_0=2AC^2/(5B).
\]
The quotient of the upper endpoint of $O_j$ by the lower endpoint
of $O_{j+1}$ is multiplied by $(R_0/H_0)^2$ when $j$ increases by one.
Thus, if $R_0/H_0>1$, one checked overlap starts the entire ordinary
tail. Since $H_0>1$, its lower endpoints tend to infinity.
The same argument applies after replacing $A$ by one in the upper
endpoints, using $R_0=2C^2/(5B)$.

For the applications in Sections~\ref{sec:all-large} and~\ref{sec:pair-reduction}, suppose in addition that
\[
 T>2BS,\qquad C>\frac{A^2T^5}{S^2},\qquad
 D>\frac{A^2C^7}{S^2T^4},
\]
and that
\begin{equation}\label{eq:pr-coarse-gap}
 C>GA^7B^5S^3,
 \qquad G_* /R_{bs}\ge\rho,
\end{equation}
where $G,G_*,\rho>0$, both $G$ and $G_*$ are valid in the first inequality,
and $G_*$ may depend on the actual triple $a,b,s$.
The constants $G$ and $\rho$ may be estimated separately.
In Table~\ref{tab:pr-overlaps}, each displayed quotient is a sufficient
lower bound for the preceding upper endpoint divided by the next
lower endpoint, including the factor $h$. The first line instead
compares the starting lower bound for $D$ with the lower endpoint of $Q$.
\begin{table}[htbp]
\centering\small
\caption{Sufficient overlap quotients, each required to exceed one.}
\label{tab:pr-overlaps}
\begin{tabular}{c|c|c}
Junction&Quotient&Additional substitution\\\hline
START--\(Q\)&\(1024A^{10}B^{10}(1-\varepsilon)/12100\)&none\\
\(Q\)--\(N_2\)&\(256A^8S^4/(36M_2^2h)\)&none\\
\(N_2\)--\(P\), \(BS^3\)&\(G^2A^{20}B^4/(16R_2^2\,2400^2h)\)&none\\
\(N_2\)--\(P\), \(B^3S^2\)&\(G^2A^{20}S^2/(16R_2^2\,2400^2h)\)&none\\
\(P\)--\(N_3\)&\(G^2A^{20}/(16R_p^2M_n^2h)\)&\(S\ge B^2\)\\
\(P\)--\(A_3\)&\(G^2A^{20}/(16R_p^2M_3^2h)\)&\(S\ge B\)\\
\(A_3\)--\(M\)&\(G^4A^{36}/(16R_3^2M_m^2h)\)&\(S\le B^2\)\\
\(M\)--\(N_3\)&\(G^4A^{34}B^2S^2/(16R_m^2M_n^2h)\)&none\\
\(N_3\)--\(A_4\)&\(G^2A^{22}B^4/(16R_n^2M_4^2h)\)&\(S\ge B^3\)\\
\(N_3\)--\(N_{3,bs}\)&\(G^4A^{36}B^8/(16R_n^2M_{bs}^2h)\)&\(S\le B^3\)\\
\(N_{3,bs}\)--\(A_4\)&\(\rho^2A^{14}/(16M_4^2h)\)&pointwise bound\\
\(A_4\)--\(O_4\)&\(G^4A^{38}B^2/(144R_4^2\,42^8h)\)&\(S\ge B\)\\
\(O_4\)--\(A_5\)&\(G^4A^{36}/(4M_5^2(5/2)^8h)\)&\(S\ge B\)\\
\(A_5\)--\(O_5\)&\(G^6A^{54}B^{10}/(144R_5^2\,42^{10}h)\)&\(S\ge B\)\\
\(O_5\)--\(O_6\)&\(G^6A^{52}B^4/(36\cdot42^{12}(5/2)^{10}h)\)&\(S\ge B\)
\end{tabular}
\end{table}

The first two comparisons use $C>A^2T^5/S^2$ and $T>2BS$.
The others use \eqref{eq:pr-coarse-gap}.
For example, before substitution, the $N_3$--$A_4$ quotient is
\[
 \frac{A^8C^2}{16R_n^2M_4^2B^{24}h}.
\]
Substituting the gap and then $S\ge B^3$ gives its table entry.
At the relabeled junction the raw quotient is
$C^2/(16R_{bs}^2M_4^2B^{10}S^6h)$.
Use the coefficient $G_*$ here to obtain the displayed
$\rho^2A^{14}$ bound.
The remaining entries follow by the same divisions. All substitutions
use a positive exponent for a lower bound or a negative exponent for an
upper bound.

The chain begins START, $Q,N_2,P$. If $S\ge B^2$, continue directly
to $N_3$. Otherwise use $A_3,M,N_3$. If $S\ge B^3$, continue
to $A_4$. Otherwise use $N_{3,bs},A_4$. Then use
\[
 A_4,O_4,A_5,O_5,O_6,O_7,\ldots .
\]
Both equality cases belong to a checked branch. Moreover,
\[
 R_0/H_0=AC/(105B^4)>(G/105)A^8BS^3.
\]
Consequently, if all the quotients in Table~\ref{tab:pr-overlaps}
and the last quantity exceed one, the ordinary-tail argument gives
a contradiction.

\section{Entries of modulus at least three}
\label{sec:all-large}


\begin{proposition}\label{prop:no-all-large-six}
A \(D(1)\)-tuple contains at most five entries of modulus at least three.
\end{proposition}
\begin{proof}
Suppose six such entries are ordered as \(a,b,s,t,c,d\), with
\(3\le A\le B\le S\le T\le C\le D\).
Lemma~\ref{lem:five-frontier-gaps} gives \(T>3ABS\) and hence \(AT>243\).
We may therefore apply the row \(H=48,\lambda=19/200\) of
Lemma~\ref{lem:retained-gap}. It gives all the required size conditions and
\begin{equation}\label{eq:integrated-6}
C>\frac{A^2T^5}{S^2},\qquad
 D>\mathcal G:=\frac{A^2C^7}{S^2T^4}.
\end{equation}
Put
\[
 \mathcal Q=\frac{|(a-b)(a-s)(b-s)|}{BS^2},\qquad
 c_*=\frac{9999}{10000}\left(\frac{200}{19}\right)^2.
\]
The retained inequality and \(T>3ABS\) give
\[
 C>3^5c_*\left(\frac1{\mathcal Q}-\frac{\mu}{A^2B^2}\right)^2
                         A^7B^5S^3.
\]
Since \(A^2B^2\ge81\) and \(0<\mathcal Q\le8\),
\[
 3^5c_*\left(\frac18-\frac{\mu}{81}\right)^2
 =\frac{16405041495117}{40000000000}>400.
\]
Consequently, with \(K=400\),
\begin{equation}\label{eq:large-gap}
 C>KA^7B^5S^3.
\end{equation}

Apply the interval calculation of Section~\ref{sec:common-intervals}.
Lemma~\ref{lem:retained-gap} gives $C>\max(1000,200S,4T)$
and $D>1000C$, so all the approximation hypotheses hold.
Use the uniform coefficient and residual bounds
\[
 M_2=243,\quad M_n=M_{bs}=4143,\quad M_3=2265,\quad
 M_4=165763,\quad M_5=3113013,
\]
\[
 R_2=\frac54,\quad R_n=R_{bs}=\frac{35}{32},\quad
 R_3=\frac{35}{16},\quad R_4=\frac{63}{8},\quad
 R_5=\frac{495}{64},\quad R_p=3.
\]
The choice of $R_p$ follows from $\mathcal R(1)<3$.
Take $G=G_*=K=400$ and $\rho=400/R_{bs}$ in
\eqref{eq:pr-coarse-gap}, using \eqref{eq:large-gap}.
Every remaining modulus power in Table~\ref{tab:pr-overlaps}
is nonnegative, so all its comparisons follow by substituting
$A=B=S=3$ and clearing positive denominators.
The smallest resulting quotient occurs at $N_{3,bs}$--$A_4$ and equals
\[
 \frac{400^2 3^{14}}{16R_{bs}^2M_4^2h}
 =\frac{7834457305497600}{5386911336360281}>\frac75.
\]
All other quotients exceed $13/10$. Finally,
\[
 (G/105)A^8BS^3\ge\frac{400\cdot3^{12}}{105}>1.
\]
The common interval calculation therefore gives a contradiction.
\end{proof}

\section{Tuples with at least two small entries}
\label{sec:two-small}
Call an entry small when its modulus is less than $3$.
The purpose of this section is to prove the following result.
\begin{proposition}\label{sm:main}
A $D(1)$-tuple in the full ring of integers of an imaginary quadratic
field with at least two small entries has at most five elements.
\end{proposition}
The finite calculations below use exact integer algorithms and explicit
search domains. We tabulate their counts and the required extensions.
Completeness follows from a reversible descent and the search bounds proved below.

\subsection{Auxiliary bounds}

\begin{lemma}\label{sm:crossblock}
Let $a_1,a_2,a_3,b_1,b_2,b_3$ be six distinct members of a
$D(1)$-tuple, separately ordered by modulus. Put
$A_i=|a_i|$, $B_i=|b_i|$, and $W_i=A_iB_i$.
If $W_1\ge9$, then
\[
 W_3>\frac{W_1^3W_2}{7}.
\]
\end{lemma}
\begin{proof}
Choose complex square roots $u_i^2=a_i$, $v_j^2=b_j$ and signs
of the integral roots such that
\[
 r_{ij}=u_iv_j f(1/(a_i b_j)),\qquad f(z)=\sqrt{1+z},\quad f(0)=1.
\]
We first show $D_3=\det(r_{ij})\ne0$. No two rows are proportional. Indeed,
squaring such a relation gives a linear polynomial vanishing at the
three distinct $b_j$, forcing the corresponding $a_i$ to coincide.
If the determinant vanished, its third row would therefore be
$\lambda$ times the first plus $\mu$ times the second, with
$\lambda\mu\ne0$. Squaring this relation gives
$2\lambda\mu r_{1j}r_{2j}=Pb_j+Q$. Squaring again at the three
values $b_j$ proves the polynomial identity
\[
 (PX+Q)^2=4\lambda^2\mu^2(a_1X+1)(a_2X+1).
\]
The right side has two distinct simple roots, a contradiction.
Thus $D_3$ is a nonzero algebraic integer and $|D_3|\ge1$.

Write $c_k=\binom{1/2}{k}$. Cauchy--Binet for the power series gives
\[
 \det(f(1/(a_i b_j)))=
 \sum_{0\le k_0<k_1<k_2}
  c_{k_0}c_{k_1}c_{k_2}
  \det(a_i^{-k_j})\det(b_i^{-k_j}).
\]
For each term, rearrangement of the exponents gives
\[
 |\det(a_i^{-k_j})|\le6\prod_{i=1}^3A_i^{-k_{3-i}},
 \qquad |c_{k_0}c_{k_1}c_{k_2}|\le|c_0c_1c_2|=1/16.
\]
The analogous estimate holds for $\det(b_i^{-k_j})$. Summing the successive
nonnegative gaps between the $k_j$ gives
\[
 \sum_{0\le k_0<k_1<k_2}\prod_{i=1}^3W_i^{-k_{3-i}}
 =\frac{W_1^{-2}W_2^{-1}}
 {(1-W_1^{-1})(1-(W_1W_2)^{-1})(1-(W_1W_2W_3)^{-1})}.
\]
This convergent geometric sum also justifies passage from truncated
series to Cauchy--Binet. Since $W_i\ge9$, the denominator is greater
than $1-(9^{-1}+9^{-2}+9^{-3})>7/8$. Hence
\[
 1\le|D_3|<\frac{18}{7}\frac{W_3^{1/2}}{W_1^{3/2}W_2^{1/2}}.
\]
Squaring proves the assertion because $49/324>1/7$.
\end{proof}
\begin{lemma}[Four roots at order one]\label{sm:lem-four-root-k1}
Let \(a_1,a_2,a_3,a_4,c,d\) be distinct nonzero entries of a
\(D(1)\)-tuple, with
\[
 1\le A=|a_1|\le B=|a_2|\le S=|a_3|\le T=|a_4|,
 \quad C=|c|>1,\quad AC>1.
\]
Put \(Z=|\sqrt{cd+1}|\), \(\eta=1-1/(AC)\), and use the function
\(e\) from Lemma~\ref{lem:ordinary-residue-construction}. Then
\begin{equation}
 2\le8192e(T)\left(1+\frac{11}{8AC}\right)
       \frac{B^2S^4T^6C^2}{Z}
       +\frac{192BS^3T^5Z}{AC^4\eta^{7/2}}.
 \label{sm:eq-four-root-exact}
\end{equation}
If \(C\ge4T\) and \(AC\ge1000\), then
\begin{equation}
 2<10300\frac{B^2S^4T^6C^2}{Z}
       +193\frac{BS^3T^5Z}{AC^4}.
 \label{sm:eq-four-root-safe}
\end{equation}
If in addition \(D=|d|\ge C\), the closed interval
\begin{equation}\label{sm:eq-four-root-window}
 2\cdot10300^2B^4S^8T^{12}C^3\le D
 \le\frac{A^2C^7}{2\cdot193^2B^2S^6T^{10}}
\end{equation}
is impossible.
\end{lemma}

\begin{proof}
Use the order-one matrix in Lemma~\ref{lem:ordinary-residue-construction}
with \(m=4\), and put
\[
 H=\prod_{i=1}^4a_i,\qquad
 \Delta=\prod_{1\le i<j\le4}(a_i-a_j).
\]
For \(j\ge1\), write
\[
 E_j=\frac{\Delta^2}{\prod_{\ell\ne j}(a_\ell-a_j)},\qquad
 E_{ji}=\frac{E_j}{a_i-a_j}\quad(i\ne j).
\]
These are polynomials. The simple and double residue formulas give
\begin{align*}
 Q_{0j1}&=4a_j^3(a_jc+1)E_j,\\
 Q_{ij1}&=4a_j^3(a_jc+1)a_iE_{ji}\quad(i\ge1,\ i\ne j),\\
 Q_{jj1}&=6a_j^3E_j-Q_{0j1}-\sum_{\substack{i\ge1\\i\ne j}}Q_{ij1}.
\end{align*}
In this case \eqref{eq:ordinary-residue-det} reads
\(\det Q_1=-1440H^3\Delta^8\ne0\), and every column with index $j\ge1$
is divisible by \(a_j^3\).

Set \(K=B^2S^4T^6\) and \(L=C+1/A\). Using $|a_i-a_j|\le2\max(A_i,A_j)$ for each anchor difference gives
\[
 A_j^3|E_j|\le2^9K,\qquad
 A_iA_j^3|E_{ji}|\le2^8K.
\]
Indeed, the product of \(\max(A_i,A_j)^2\) over all six pairs is
\(K\). In \(E_j\), one factor is removed for each pair containing
\(j\), and their product is at least \(A_j^3\). Forming \(E_{ji}\)
removes one further factor \(\max(A_i,A_j)\ge A_i\). Nine and eight
difference factors remain, respectively, giving the stated powers of two.
It follows that
\[
 \frac{|Q_{0j1}|}{A_j}\le2048LK,\qquad
 \frac{|Q_{ij1}|}{A_j}\le1024LK\quad(i\ge1,\ i\ne j),
 \qquad 6A_j^2|E_j|\le\frac{3072K}{A}.
\]
Row zero has four terms. In a row with index $i\ge1$, substituting the formula for $Q_{ii1}$
gives one $Q_{0i1}$ term and six off-diagonal terms, together with
$6a_i^3E_i$. Thus every row satisfies
\[
 \sum_{j=1}^4\frac{|Q_{ij1}|}{A_j}
 \le(8192C+11264/A)B^2S^4T^6.
\]

For \(m=4,k=1\), the integral in the proof of
\eqref{eq:ordinary-residue-error} equals
\[
 \int_0^\infty\frac{x^{3/2}}{(x+\eta)^6}\,dx
 =\eta^{-7/2}B(5/2,7/2)
 =\frac{3\pi}{256}\eta^{-7/2}.
\]
Hence the residual is at most
\(3|\Delta|^2/(64|H|C^4\eta^{7/2})\).
Since \(|\Delta|^2\le2^{12}B^2S^4T^6\) and \(|H|=ABST\),
this is at most \(192BS^3T^5/(AC^4\eta^{7/2})\).
Apply the final, even-coefficient assertion of
Lemma~\ref{lem:ordinary-residue-construction} to obtain
\eqref{sm:eq-four-root-exact}.

Under the additional size assumptions,
\(e(T)^2\le25\cdot1001/16000\) and \(\eta\ge999/1000\).
The exact rational comparisons
\[
 \left[8192\left(1+\frac{11}{8000}\right)\right]^2
 \frac{25\cdot1001}{16000}<10300^2,
 \qquad
 192^2\left(\frac{1000}{999}\right)^7<193^2
\]
prove \eqref{sm:eq-four-root-safe}.
Finally \(CD/2<Z^2<2CD\). Membership in
\eqref{sm:eq-four-root-window} would therefore imply
\[
 Z>10300B^2S^4T^6C^2,\qquad
 Z<\frac{AC^4}{193BS^3T^5}.
\]
Both terms on the right of \eqref{sm:eq-four-root-safe} would be less
than one, a contradiction.
\end{proof}


\subsection{Two small entries}

We first bound the third large entry. The descent in Lemma~\ref{sm:bs7-descent} generates all common extensions within the resulting disk from a finite initial set.

Let $a,b,s,t,p,q$ be six distinct tuple entries, with moduli
$A\le B\le S\le T\le P\le Q$. Assume
\begin{equation*}
 P\ge\max\{4T,1000\},\qquad Q\ge4P,
 \qquad P^2\ge\frac{32BST}{A^3}.
\end{equation*}
Lemma~\ref{lem:four-anchor-nonvanishing} gives a nonzero three-anchor form among the two triples sharing the first two anchors. Indeed, simultaneous vanishing would give
\begin{equation*}
 |pq+1|\le\frac{121BST}{4A^3}
 <P^2-1\le PQ-1,
\end{equation*}
where the strict inequality follows from $BST/A^3\ge1$.
Lemma~\ref{lem:separated-three} therefore gives
\begin{equation*}
 1<\frac{5BT^2}{Z}+\frac{3TZ}{4AP^3},
 \qquad Z=|\sqrt{pq+1}|.
\end{equation*}
Put $\varepsilon=1/4000$. Since $P^2\ge4AP\ge4000$, we have
$(1-\varepsilon)PQ\le Z^2\le(1+\varepsilon)PQ$.
Consequently the closed interval
\begin{equation}\label{sm:six-V-interval}
 \frac{100B^2T^4}{(1-\varepsilon)P}\le Q
 \le\frac{4A^2P^5}{9(1+\varepsilon)T^2}
\end{equation}
is impossible. We shall use the exact comparisons
\begin{equation*}
 \frac{100}{1-\varepsilon}<101,
 \qquad \frac4{9(1+\varepsilon)}>\frac25.
\end{equation*}

\begin{lemma}\label{sm:bs7-cutoff}
Suppose a $D(1)$-tuple has six entries with ordered moduli
\[
 1\le A\le B<3\le S\le T\le C\le D.
\]
Then $C<10^{550}$.
\end{lemma}
\begin{proof}
Lemmas~\ref{lem:product-four} and~\ref{sm:crossblock} give
\begin{equation}\label{sm:five-two-small-products}
 C>\frac38AST,\qquad D>2STC,\qquad
 D>\frac{A^3BT^3C}{7S},
\end{equation}
where the last inequality is used only when $AT\ge9$.
For the first inequality apply Lemma~\ref{lem:product-four} to
$(a,s,t,c)$. In the last one the row entries are $(a,b,s)$ and the
column entries are $(t,c,d)$.

Taking $a,b$ as anchors and any of $s,t,c$ with modulus $P\ge10^8$ as the interpolation entry, Corollary~\ref{cor:ordinary-height-envelope} gives
\begin{equation}\label{sm:five-uniform-continuous}
 \log_{10}D<g(\log_{10}P),\qquad
 g(x)=7x+\frac{212}{5}+\frac{4047/25}{x-4}.
\end{equation}
Indeed, $H_2<1134P$, $R_2>(2/15)P^2$ and $253B^2<2277$.
The parameters $(h,r,u)=(31/10,-9/10,17/5)$ follow from
$1134^{10}<10^{31}$, $(15/2)^{10}<10^9$ and $2277^5<10^{17}$.
The function $g$ is convex on $(4,\infty)$.

First assume $S\ge10^8$. The first inequality in
\eqref{sm:five-two-small-products} implies
$C\ge\max\{4T,1000\}$, $AC\ge1000$ and $C^2\ge32BST/A^3$.
Moreover, $D>18C$ and
\[
 DC>\frac9{448}A^5BST^5>101B^2T^4,
\]
since $9A^5ST>45248B$. Thus $D$ lies above the lower endpoint of
\eqref{sm:six-V-interval}, and it must exceed the upper endpoint. Hence
\begin{equation*}
 D>\frac25\frac{A^2C^5}{T^2}
  >\frac{243}{81920}A^7S^5T^3
  >\frac{S^5T^3}{400}\ge\frac{S^8}{400}.
\end{equation*}
If $x=\log_{10}S\ge49$, then $400<10^3$ gives $\log_{10}D>8x-3$. This contradicts~\eqref{sm:five-uniform-continuous}, since $(8x-3)-g(x)=x-227/5-(4047/25)/(x-4)$ is increasing and equals $1/375$ at $49$.
Therefore $S<10^{49}$. In the whole case $S\ge10^8$, convexity gives
$\log_{10}D<389$, since $g(8)<139$ and $g(49)=389-1/375$.

Now assume $S<10^8$ and $T\ge10^{72}$. Use $L_3$ with anchors
$(a,b,s)$ and terminal pair $(t,c)$. Lemma~\ref{lem:general-three}
applies. If $L_3\ne0$, its forbidden interval contains
\[
 \frac{101B^2S^4}{T}\le C\le\frac{A^2T^5}{37S^2}.
\]
Here we used $100/(1-1/4000)<101$ and
$36(1+1/4000)<37$. The lower endpoint is below one, so
$C>T^5/(37S^2)$. Equation~\eqref{sm:five-two-small-products} gives
$D>T^8/(259S^3)>T^8/10^{27}$, contrary to
\eqref{sm:five-uniform-continuous}. Indeed,
the function $x-347/5-(4047/25)/(x-4)$ is increasing for $x\ge72$ and positive at $x=72$.

If $L_3=0$, we first prove $C>4T$ in order to apply
Lemma~\ref{lem:separated-three}. Since $AT\ge T\ge10^{72}>11$,
the bound on $F_3$ in the proof of Lemma~\ref{lem:general-three}
gives $|F_3|\le1331/(4000T^3)<1/(3T^3)$.
The two terms in~\eqref{eq:coherent-exact} satisfy
\[
 |1-f(1/(tc))|>\frac1{3TC},\qquad
 \frac{|F_3|}{ABS}<\frac1{3ABS T^3}.
\]
Their equality implies $C>ABS T^2>4T$.
Lemma~\ref{lem:separated-three} now gives
$C>5ABS T^2\ge5ST^2$. The four-anchor nonvanishing hypotheses
at $C$ follow, and
$DC>2STC^2>50S^3T^5>101B^2T^4$.
Thus~\eqref{sm:six-V-interval} gives
$D>(2/5)C^5/T^2>1250S^5T^8>T^8$, again contradicting
\eqref{sm:five-uniform-continuous}. Consequently $T<10^{72}$.
If $T\ge10^8$, convexity and $g(72)<549$ give $D<10^{549}$.

It remains to treat $S,T<10^8$. Put $M=10^8$ and suppose
$C\ge10^{100}$. Four-anchor nonvanishing and
Lemma~\ref{lem:general-three} give
\[
 1<\frac{15M^2}{Z}+\frac{3MZ}{C^3},\qquad
 Z=|\sqrt{cd+1}|>C/2.
\]
The first term is less than $1/2$, so
$D>C^5/(100M^2)>\max\{1000C,17ABSTC^3\}$.
The upper bound for $D$ when $L_4=0$ therefore implies $L_4\ne0$.
Lemma~\ref{lem:direct-four} gives
\[
 1<\frac{39M^5}{Z}+\frac{49M^3Z}{50C^4}.
\]
Again the first term is less than $1/2$. Hence
\begin{equation*}
 D>\frac{C^7}{100M^6}=10^{-50}C^7.
\end{equation*}
Lemma~\ref{lem:ordinary-three-normalized} at order two, with anchors
$(a,b,s)$, forbids
\[
 \frac{512}{81}608^4B^6S^{16}C^5\le D
 \le\frac{A^6C^{11}}{2\cdot28^4(10/81)^2B^4S^{10}}.
\]
Since $(512/81)608^4 3^6<10^{15}$ and
$2\cdot28^4(10/81)^2 3^4<10^7$, the smaller interval
$10^{143}C^5\le D\le10^{-87}C^{11}$ is also forbidden.
Its lower endpoint is less than $10^{-50}C^7$.
Its upper endpoint exceeds the bound $D<10^{60}C^7$ obtained from
\eqref{sm:five-uniform-continuous}. This contradiction completes all cases.
\end{proof}

\begin{lemma}\label{sm:bs7-descent}
Let $a,b\in\OK$ be distinct nonzero elements with $|a|,|b|<3$, and suppose $r_0\in\OK$ satisfies $r_0^2=ab+1\ne0$. Every integral common extension~$z$ of~$a,b$ with $|z|\ge33$ has an integral common extension~$w$ satisfying
\begin{equation*}
 |w|<|z|/5+26<|z|\,.
\end{equation*}
Here a common extension may be zero or equal to an anchor. The transformations used to obtain~$w$ are reversible within this enlarged set.
\end{lemma}
\begin{proof}
Choose roots $r^2=az+1$ and $s^2=bz+1$. The two transformations are
\begin{equation}\label{sm:bs7-mutation}
 z_\pm=a+b+(2ab+1)z\pm2r_0rs\,.
\end{equation}
For the plus sign, the integral roots $r_+=r_0r+as$ and $s_+=br+r_0s$ satisfy $r_+^2=az_++1$ and $s_+^2=bz_++1$. Changing the sign of~$s$ proves the other case. Expansion also gives
\begin{equation*}
 a+b+(2ab+1)z_+-2r_0r_+s_+=z\,.
\end{equation*}
These identities involve no division by roots or extensions, so they remain valid when a root vanishes or a transformed value is zero or equals an anchor.

Put $p=ab$ and choose $\eta^2=p$. The numbers
\begin{equation*}
 \lambda_\pm=2p+1\pm2r_0\eta
\end{equation*}
have product one. We show that the smaller modulus is less than~$1/5$. If $|r_0|^2\ge2$, then $|2r_0^2-1|\ge3$. Otherwise $r_0$ is a unit of modulus one, and its integral trace lies in $\{-2,-1,0,1,2\}$. Trace~$\pm2$ would give $r_0=\pm1$ and $ab=0$. Trace zero gives $r_0^2=-1$. If the trace is $\pm1$, then $r_0^2$ is a primitive cube root of unity. Thus in every permitted case $|2r_0^2-1|^2\ge7$.

Let $q\ge1$ be the larger modulus of~$\lambda_+,\lambda_-$. If $q\le5$, then
\begin{equation*}
 2\sqrt7\le|\lambda_++\lambda_-|\le q+q^{-1}\le5+1/5\,,
\end{equation*}
contrary to $28>676/25$. Choose the sign of~$\eta$ so that $|2p+1-2r_0\eta|<1/5$.

One of $\tau=rs,-rs$ satisfies $|\tau+\eta z|\ge|\eta z|$. Since $\tau^2-(\eta z)^2=(a+b)z+1$, it follows that
\begin{equation*}
 |\tau-\eta z|\le\frac{|a|+|b|+1/|z|}{\sqrt{|ab|}}\,.
\end{equation*}
The corresponding value $w=a+b+(2p+1)z-2r_0\tau$ in~\eqref{sm:bs7-mutation} therefore satisfies
\begin{equation*}
 |w|<\frac{|z|}{5}+6+14\sqrt2<\frac{|z|}{5}+26\,,
\end{equation*}
using $|a|+|b|<6$, $|ab|\ge1$, $|z|\ge1$, and $|r_0|/\sqrt{|ab|}\le\sqrt2$. For $|z|\ge33$ this is strictly smaller than~$|z|$.
\end{proof}

Repeated application of Lemma~\ref{sm:bs7-descent} terminates, since squared moduli of algebraic integers are nonnegative integers. Consequently all common extensions in any closed disk of radius~$R\ge33$ can be generated from those in the closed disk of radius~$33$. Repeatedly apply both transformations~\eqref{sm:bs7-mutation}, retaining new values only when their moduli are at most~$R$. Reversing a decreasing chain proves completeness of this procedure. Zero and anchor coincidences must be retained until the generation is finished.

\begin{lemma}\label{sm:bs7-finite-pairs}
For every squarefree $d<36$ and every compatible small pair $a,b\in\mathcal O_{\mathbb Q(\sqrt{-d})}$ with $ab+1\ne0$, at most two members of a $D(1)$-tuple containing~$a,b$ have moduli in $[3,10^{550}]$.
\end{lemma}
\begin{proof}
We describe the complete finite calculation, including its arithmetic tests and search domains. Write $t=\sqrt{-d}$ and represent $(u+vt)/2$ by~$(u,v)$. Valid coordinates have matching parity when $d\equiv3\pmod4$, and are both even otherwise. Their norm is $(u^2+dv^2)/4$.

For an element $(u+vt)/2$, the following is a necessary and sufficient square test. Require
\begin{equation*}
 \frac{u^2+dv^2}{4}=L^2,\qquad
 U^2=2L+u,\qquad dV^2=2L-u,\qquad UV=|v|\,,
\end{equation*}
with nonnegative integers~$L,U,V$, and require that $(U,V)$ satisfy the coordinate parity rule. A square root is $(U+\epsilon Vt)/2$, where $\epsilon=1$ if $v\ge0$ and $\epsilon=-1$ otherwise. Squaring this expression proves sufficiency. Conversely, the norm of a square gives~$L$, and adding and subtracting its real-coordinate equations gives the displayed conditions. Thus the test and a choice of square root use only integer square roots, divisibility, and parity.

Enumerate the squarefree integers $1\le d<36$. For each~$d$, enumerate all nonzero valid coordinates with $u^2+dv^2<36$, form all distinct compatible pairs, and select the lexicographically least representative under simultaneous negation, conjugation, and interchange of the two entries. This gives 72 classes. Twenty-three are the pair $\{-1,1\}$, one in each field. One further class has product~$-1$, namely
\begin{equation}\label{sm:bs7-exceptional-pair}
 d=3,\qquad a=(-1-t)/2,\qquad b=(1-t)/2\,.
\end{equation}
The remaining 48 representatives are listed in Table~\ref{sm:bs7-pair-table}.

For each of these 48 pairs, enumerate every valid coordinate with $u^2+dv^2\le4\cdot33^2$ and retain precisely its common extensions, including zero and any anchor coincidences. Close this finite initial set under both transformations~\eqref{sm:bs7-mutation}, retaining a new value exactly when $u^2+dv^2\le4\cdot10^{1100}$. Any root returned by the square test may be used, since both signs of the transformation are included. Continue until no new value is produced. The descent lemma proves that the resulting set is the complete closed radius-$10^{550}$ set of common extensions.

Now discard zero, the two anchors, and entries of modulus below three. Join every pair of remaining vertices whose product plus one passes the square test. A triangle is detected by taking, for every edge $\{x,y\}$, the intersection of the two neighbor sets and testing whether it contains a third vertex. Table~\ref{sm:bs7-pair-table} gives the exact vertex and edge counts. The 48 graphs contain 36931 vertices and 36854 edges in total, and every triangle test is negative. Since any three compatible extensions in the disk would form a triangle, the assertion follows.
\end{proof}

\begin{longtable}{rccrr}
\caption{Nondegenerate small pairs and their complete extension graphs through modulus $10^{550}$. A coordinate pair $(u,v)$ represents $(u+v\sqrt{-d})/2$. Every graph has zero triangles.}\label{sm:bs7-pair-table}\\
\toprule
$d$ & First anchor & Second anchor & Vertices & Edges\\
\midrule
\endfirsthead
\toprule
$d$ & First anchor & Second anchor & Vertices & Edges\\
\midrule
\endhead
\midrule
\multicolumn{5}{r}{Continued on the next page}\\
\endfoot
\bottomrule
\endlastfoot
1 & $(-4,-4)$ & $(-4,4)$ & 718 & 716 \\
1 & $(-4,-4)$ & $(0,-4)$ & 414 & 413 \\
1 & $(-4,-4)$ & $(4,0)$ & 827 & 825 \\
1 & $(-4,-2)$ & $(0,-2)$ & 1192 & 1190 \\
1 & $(-4,-2)$ & $(0,4)$ & 861 & 859 \\
1 & $(-4,-2)$ & $(4,-2)$ & 1315 & 1312 \\
1 & $(-4,0)$ & $(-2,-4)$ & 862 & 860 \\
1 & $(-4,0)$ & $(2,0)$ & 1434 & 1432 \\
1 & $(-2,-4)$ & $(2,-4)$ & 438 & 437 \\
1 & $(-2,-4)$ & $(2,0)$ & 1192 & 1190 \\
1 & $(-2,-2)$ & $(2,-2)$ & 718 & 717 \\
1 & $(0,-4)$ & $(0,-2)$ & 718 & 717 \\
2 & $(-4,-2)$ & $(0,-2)$ & 496 & 495 \\
2 & $(-4,0)$ & $(2,-2)$ & 991 & 989 \\
2 & $(-2,-2)$ & $(-2,2)$ & 961 & 959 \\
2 & $(-2,-2)$ & $(2,-2)$ & 552 & 551 \\
2 & $(0,-4)$ & $(0,4)$ & 718 & 716 \\
3 & $(-5,-1)$ & $(-1,-1)$ & 524 & 523 \\
3 & $(-5,-1)$ & $(3,1)$ & 892 & 890 \\
3 & $(-4,-2)$ & $(0,-2)$ & 446 & 445 \\
3 & $(-4,-2)$ & $(1,1)$ & 1048 & 1046 \\
3 & $(-4,0)$ & $(4,0)$ & 962 & 960 \\
3 & $(-3,-1)$ & $(-3,1)$ & 962 & 960 \\
3 & $(-3,-1)$ & $(1,-3)$ & 892 & 890 \\
3 & $(-3,-1)$ & $(1,-1)$ & 761 & 760 \\
3 & $(-3,-1)$ & $(2,0)$ & 1522 & 1520 \\
3 & $(-2,-2)$ & $(2,-2)$ & 481 & 480 \\
3 & $(-2,0)$ & $(-1,-3)$ & 1048 & 1046 \\
3 & $(-1,-1)$ & $(0,-2)$ & 761 & 760 \\
3 & $(0,-2)$ & $(0,2)$ & 962 & 960 \\
5 & $(-2,-2)$ & $(2,-2)$ & 410 & 409 \\
6 & $(-2,-2)$ & $(2,-2)$ & 389 & 388 \\
7 & $(-5,-1)$ & $(-5,1)$ & 718 & 716 \\
7 & $(-5,-1)$ & $(-1,-1)$ & 457 & 456 \\
7 & $(-5,-1)$ & $(0,2)$ & 737 & 735 \\
7 & $(-5,-1)$ & $(2,0)$ & 1121 & 1119 \\
7 & $(-5,-1)$ & $(5,-1)$ & 744 & 742 \\
7 & $(-3,-1)$ & $(1,-1)$ & 561 & 560 \\
7 & $(-3,-1)$ & $(3,1)$ & 914 & 912 \\
7 & $(-2,-2)$ & $(-2,2)$ & 718 & 716 \\
7 & $(-2,-2)$ & $(2,-2)$ & 372 & 371 \\
11 & $(-3,-1)$ & $(1,-1)$ & 484 & 483 \\
11 & $(-1,-1)$ & $(-1,1)$ & 962 & 960 \\
15 & $(-3,-1)$ & $(1,-1)$ & 440 & 439 \\
19 & $(-3,-1)$ & $(1,-1)$ & 411 & 410 \\
23 & $(-3,-1)$ & $(-3,1)$ & 718 & 716 \\
23 & $(-3,-1)$ & $(1,-1)$ & 389 & 388 \\
31 & $(-1,-1)$ & $(-1,1)$ & 718 & 716 \\
\midrule
\multicolumn{3}{r}{Total} & 36931 & 36854\\
\end{longtable}

\begin{samepage}
\begin{lemma}\label{sm:bs7-unit-axes}
Let $d$ be positive and squarefree, and put $t=\sqrt{-d}$. Every nonzero common extension of $\{-1,1\}$ has one of the forms
\begin{align*}
 z&=\pm(dv^2+1),& u^2-dv^2&=2,\\
 z&=\pm2uvt,& u^2-dv^2&=1,
\end{align*}
where $u,v\in\mathbb Z$ and $v\ne0$. In particular $|z|\ge d+1$.
\end{lemma}
\end{samepage}
\begin{proof}
Choose integral roots $r^2=z+1$ and $s^2=1-z$, and write $2r=U+Vt$, $2s=X+Yt$. The equation $r^2+s^2=2$ gives
\begin{equation*}
 U^2+X^2-d(V^2+Y^2)=8,\qquad UV+XY=0\,.
\end{equation*}
The real coordinate vector is nonzero, since otherwise the left side of the first equation would be nonpositive. If the imaginary coordinate vector vanished, both roots would be rational integers, forcing $r^2=s^2=1$ and $z=0$. Write $(U,X)=g(p,q)$ with $g>0$ and $\gcd(p,q)=1$. Orthogonality gives $(V,Y)=h(q,-p)$ for a nonzero integer~$h$, so
\begin{equation*}
 (p^2+q^2)(g^2-dh^2)=8\,.
\end{equation*}
Primitivity excludes $p^2+q^2=4,8$. Thus this sum is one or two. If it is one, one root is rational and the other purely imaginary. Integrality forces $g,h$ even, including in the half-integral basis. Setting $u=g/2$, $v=h/2$ gives the first alternative.

If $p^2+q^2=2$, then $p,q\in\{-1,1\}$ and $g^2-dh^2=4$. Outside the half-integral basis, $g,h$ are even. In that basis they have matching parity. If both were odd, then $g^2-dh^2\equiv2\pmod4$, contrary to its value four. Thus $g,h$ are again even, and expansion gives the second alternative. The first has modulus at least~$d+1$, while in the second
\begin{equation*}
 |z|^2=4du^2v^2\ge4d(d+1)>(d+1)^2\,.
\end{equation*}
\end{proof}

\begin{lemma}\label{sm:five-unit-cubic}
Suppose $\{-1,1,c,d\}$ is a $D(1)$-quadruple, with
$9\le C=|c|\le D=|d|$ and $d\ne-c$. Then $D>C^3/2$.
\end{lemma}
\begin{proof}
Choose $z\in\OK$ with $z^2=cd+1$, and put $w=c+d$, $h=w/z^2$.
Then $|h|\le2C/(C^2-1)\le9/40<1/4$.
With the analytic square-root branches at zero, define
$u=z\sqrt{1+h}$, $v=z\sqrt{1-h}$ and $L=u+v-2z$.
We have $u^2=(c+1)(d+1)$ and $v^2=(1-c)(1-d)$.
Both products are squares in $\OK$, so $u,v$, and hence $L$, lie in $\OK$.
If $L=0$, squaring $u+v=2z$ gives $uv=z^2$, and squaring again gives
$w^2=0$. This contradicts $d\ne-c$, so $|L|\ge1$.
The even binomial coefficients have modulus at most $1/8$ after the
constant term. Therefore
\[
 |L|\le\frac{|z||h|^2}{4(1-|h|^2)}
 <\frac4{15}\frac{|w|^2}{|z|^3}
 \le\frac{16}{15}\left(\frac{81}{80}\right)^{3/2}
       \frac{\sqrt D}{C^{3/2}}
 <\sqrt{\frac{2D}{C^3}}.
\]
Here $|w|\le2D$, $|z|^2\ge(80/81)CD$, and
$(16/15)^2(81/80)^3<2$. The conclusion follows.
\end{proof}

\begin{lemma}\label{sm:bs7-unit-pair}
No $D(1)$-tuple in an imaginary quadratic integer ring contains $\{-1,1\}$ and four further entries.
\end{lemma}
\begin{proof}
We first restrict the axes and field parameters. A negative rational integer which is an integral square in $\mathbb Q(\sqrt{-d})$ is $-dm^2$ for an integer~$m$. Indeed, the square of a rational linear combination of $1,t$ is rational only when one coordinate vanishes. Integrality of a purely imaginary root forces its coefficient of~$t$ to be an integer, also in the half-integral basis. A positive rational integer that is a square in $\OK$ is therefore an ordinary integer square.

There are no nonzero common extensions when $d=1$, since neither $u^2-v^2=2$ nor $u^2-v^2=1$ has a solution with $v\ne0$. For $d>1$, at most one extension can lie on each imaginary half-axis. Two with coefficients of the same sign would have product plus one equal to $1-dkl<0$, which is one modulo~$d$ and cannot be $-dm^2$.

Real and imaginary extensions cannot coexist in one tuple.

Suppose first that $\{-1,1,a,-a,b\}$ is a $D(1)$-tuple with $|a|,|b|\ge3$. The element $-b$ may also be adjoined. Indeed, $1-b^2=(1-b)(1+b)$ is a square, and its other compatibility conditions follow by changing signs. We may therefore interchange~$a,b$ and assume $A=|a|\le B=|b|$. Choose roots $x^2=1-b$, $y^2=1+b$, $v^2=1-ab$, $w^2=1+ab$, and put
\begin{equation*}
 G=2a^2b^2-a^2-1,\qquad J=2axyvw\,.
\end{equation*}
Expansion gives $(G-J)(G+J)=(a^2-1)^2$. Both factors are nonzero algebraic integers. Since $|G-J|,|G+J|\ge1$, we have
\[
 |G-J|+|G+J|\le|G-J||G+J|+1.
\] Consequently
\begin{equation*}
 2|G|\le(A^2+1)^2+1,\qquad |G|\ge2A^4-A^2-1\,.
\end{equation*}
These imply $3A^4-4A^2-4\le0$, impossible for $A\ge3$. Now if a real extension~$a$ and an imaginary extension~$b$ coexisted, conjugation would make $1-ab$ a square, while $(1-a)(1+a)$ would make $1-a^2$ a square. Adjoining~$-a$ would give the forbidden five entries. Lemma~\ref{sm:bs7-unit-axes} ensures that all these extensions have modulus at least three.

If $d$ is even, a real extension has positive magnitude $u^2-1$ with $u$ even. In $u^2-dv^2=2$, an even~$v$ is impossible modulo four, and an odd~$v$ forces~$u$ even. Thus every positive magnitude is three modulo four. Two real extensions of the same sign would have a positive product plus one congruent to two modulo four. There are therefore at most two real extensions. If $d$ is odd, an even~$v$ in the same Pell equation gives $u^2\equiv2$ or~$6\pmod8$, which is impossible. An odd~$v$ gives $u^2\equiv d+2\pmod8$, forcing $d\equiv7\pmod8$. We need only consider real extensions for odd squarefree~$d\equiv7\pmod8$.

The argument excluding $\{-1,1,a,-a,b\}$ shows that the four real extensions have distinct absolute values. Write these as $S<T<C<D$. If $S\ge10^4$, Lemma~\ref{sm:five-unit-cubic}, applied three times, gives $D>S^{27}/2^{13}>S^{27}/10^4$. On the other hand, Corollary~\ref{cor:ordinary-height-envelope} with parameters
\[
 (h,r,u)=\left(\frac{163}{100},-\frac25,\frac{241}{100}\right)
\]
gives, for $s=\log_{10}S$,
\[
 \log_{10}D<7s+\frac{617}{25}
       +\frac{118401/2500}{s-203/100}.
\]
The parameters follow from $42^{100}<10^{163}$, $(5/2)^5<10^2$ and $253^{100}<10^{241}$. The difference from the lower bound $27s-4$ is increasing for $s\ge4$ and positive at $4$, since
$80-717/25-118401/4925>0$. Hence $S<10^4$.
Lemma~\ref{sm:bs7-unit-axes} gives $d<S$, so it remains to consider squarefree $7\le d<10^4$ with $d\equiv7\pmod8$.

Put $R=10^{550}$. We enumerate all real common extensions in the closed radius-$R$ disk by ordinary continued fractions. By Lemma~\ref{sm:bs7-unit-axes}, their positive magnitudes are $u^2-1$, where
\begin{equation*}
 u^2-dv^2=2,\qquad
 1\le u\le\lfloor\sqrt{R+1}\rfloor,\qquad v\ge1.
\end{equation*}
Every such pair $u,v$ is coprime, and
$u/v+\sqrt d>2\sqrt d>4$, since $d\ge7$. Hence
\begin{equation*}
 \left|\frac uv-\sqrt d\right|
 =\frac{2}{v^2(u/v+\sqrt d)}<\frac1{2v^2}.
\end{equation*}
We recall a short proof of the criterion used here. If $p/q$ is reduced
and $|\alpha-p/q|<1/(2q^2)$, choose consecutive convergents
$p_n/q_n,p_{n+1}/q_{n+1}$ with $q_n\le q<q_{n+1}$.
Their determinant is $\pm1$, and the errors
$e_j=q_j\alpha-p_j$ alternate in sign. Indeed, the recurrence gives
$p_nq_{n-1}-p_{n-1}q_n=(-1)^{n-1}$ and, if
$\alpha_{n+1}>0$ is the next complete quotient,
\[
 \alpha=\frac{p_n\alpha_{n+1}+p_{n-1}}
 {q_n\alpha_{n+1}+q_{n-1}}.
\]
These identities give the two assertions directly. If $p/q$ is not a convergent, express
$(p,q)=h(p_n,q_n)+k(p_{n+1},q_{n+1})$ with integers $h,k$.
The denominator bounds and reducedness force $h,k$ to be nonzero
and of opposite signs. Hence
$|q\alpha-p|=|h e_n+k e_{n+1}|\ge|e_n|$. But then
\[
 1\le|p q_n-q p_n|
 \le(q_n+q)|q\alpha-p|\le2q|q\alpha-p|<1,
\]
a contradiction. Applying this criterion with $\alpha=\sqrt d$
shows that every fraction $u/v$ above is a convergent.

For completeness, the entire computation uses the following integer recurrences. For each admissible $d$, set $m_0=0$, $q_0=1$, and $a_0=\lfloor\sqrt d\rfloor$, and iterate
\begin{equation*}
 m_{n+1}=q_na_n-m_n,\qquad
 q_{n+1}=\frac{d-m_{n+1}^2}{q_n},\qquad
 a_{n+1}=\left\lfloor\frac{a_0+m_{n+1}}{q_{n+1}}\right\rfloor.
\end{equation*}
The convergent numerators and denominators satisfy
\begin{equation*}
 u_n=a_nu_{n-1}+u_{n-2},\qquad
 v_n=a_nv_{n-1}+v_{n-2},
\end{equation*}
with $(u_{-2},u_{-1})=(0,1)$ and $(v_{-2},v_{-1})=(1,0)$.
Stop at the first numerator exceeding $\lfloor\sqrt{R+1}\rfloor$, and retain precisely the convergents for which $v_n>0$ and $u_n^2-dv_n^2=2$. Add both vertices $\pm(u_n^2-1)$. The increasing numerators and Legendre's criterion prove that no extension in the disk is omitted.

Test every pair of distinct vertices by the rational square criterion at the start of this proof. As before, detect triangles by intersecting the neighbor sets at each edge. The complete computation gives the following counts.
\begin{center}
\begin{tabular}{rrrrr}
\toprule
Fields & Nonempty graphs & Vertices & Edges & Triangles\\
\midrule
1016 &409 &19612 &9806 &0\\
\bottomrule
\end{tabular}
\end{center}
Each vertex $x$ is adjacent to $-x$, since
$1-x^2=(1-x)(1+x)$ is a square. The total number of edges is half
the total number of vertices, so these pairs are all the edges of
the computed graphs.
Thus four compatible extensions would have their third modulus greater than $10^{550}$, contrary to Lemma~\ref{sm:bs7-cutoff}.
\end{proof}
\begin{lemma}\label{sm:bs7-eisenstein}
A $D(1)$-tuple containing the pair in~\eqref{sm:bs7-exceptional-pair} has at most four entries.
\end{lemma}
\begin{proof}
Put $t=\sqrt{-3}$, $a=(-1-t)/2$, and $b=(1-t)/2$. Then $a^3=1$, $b=-a^2=-\overline a$, and $|a|=1$. For a common extension~$z$, choose roots $r^2=az+1$ and $s^2=bz+1$, and set $q=as$. Their equations give
\begin{equation*}
 r^2+q^2=1+a^2=-a\,.
\end{equation*}
The integral transformation
\begin{equation*}
 T(r,q)=(2r+tq,-tr+2q)
\end{equation*}
preserves this sum of squares and has integral inverse $T^{-1}(r,q)=(2r-tq,tr+2q)$, since $4+t^2=1$. Under the embedding $t=i\sqrt3$, the coordinates $u=r+iq$, $v=r-iq$ are multiplied by $\alpha=2+\sqrt3$ and~$\alpha^{-1}$, respectively. Their product is~$-a$, so $|u||v|=1$.

Choose an integer power of~$T$ for which $\alpha^{-1/2}\le|u|\le\alpha^{1/2}$. The transformed integral roots satisfy
\begin{equation*}
 |r|,|q|\le\frac{|u|+|v|}{2}
 \le\frac{\sqrt\alpha+1/\sqrt\alpha}{2}=\frac{\sqrt6}{2}\,.
\end{equation*}
Their squared moduli are nonnegative integers, hence at most one. The six units in this ring are $\pm1,\pm a,\pm a^2$, as follows by solving $U^2+3V^2=4$ with matching coordinate parity. Their squares are $1,a,a^2$. Since $-a$ is not one of those squares, neither root vanishes. The only two unit squares adding to~$-a$ are $1$ and~$a^2$. Thus the transformed pair is $(\pm1,\pm a)$ or $(\pm a,\pm1)$, and satisfies
\begin{equation*}
 q=\epsilon a\overline r,\qquad\epsilon\in\{-1,1\}\,.
\end{equation*}
This relation is preserved by~$T$ and its inverse. Indeed, if it holds, then $\epsilon a\overline{2r+tq}=2q-tr$. Replacing $t$ by $-t$ gives the identity for the inverse. It therefore holds for the original roots. As $q=as$, we have $s=\epsilon\overline r$, and the original equations give $bz=\overline a\,\overline z$. Since $b=-\overline a$, every common extension is purely imaginary.

A purely imaginary Eisenstein integer is~$kt$ with $k\in\mathbb Z$, by the coordinate parity rule. Two nonzero extensions on the same half-axis have product plus one equal to $1-3kl<0$, with $k,l$ positive integers. This is one modulo three and cannot be an integral square, since such a negative rational square would be~$-3m^2$. At most one extension lies on each half-axis.
\end{proof}

\begin{proposition}\label{sm:five-two-small-four-large}
A $D(1)$-tuple in any imaginary quadratic integer ring cannot contain two entries of modulus below three and four further entries of modulus at least three.
\end{proposition}
\begin{proof}
Write $K=\mathbb Q(\sqrt{-d})$ with~$d$ positive and squarefree. If $d\ge36$, every nonreal algebraic integer has modulus at least~$\sqrt d/2\ge3$. The small entries are among $-2,-1,1,2$. Testing their six products plus one gives only the compatible pair $\{-1,1\}$. The other values are $3,-1,-3$, and a negative rational square in this field is~$-dm^2$. Lemma~\ref{sm:bs7-unit-pair} excludes this case.

If $d<36$, the complete enumeration in Lemma~\ref{sm:bs7-finite-pairs} gives 72 classes. The unit pairs are excluded by Lemma~\ref{sm:bs7-unit-pair}, and the exceptional product-$-1$ pair by Lemma~\ref{sm:bs7-eisenstein}. For every other pair, Lemma~\ref{sm:bs7-finite-pairs} forces the third large modulus to exceed~$10^{550}$, contradicting Lemma~\ref{sm:bs7-cutoff}.
\end{proof}




\subsection{Three small entries}

\begin{lemma}\label{sm:six-small-triple-classification}
A $D(1)$-tuple in an imaginary quadratic integer ring has at most three entries of modulus less than $3$. Up to simultaneous negation and conjugation, its small triples are precisely those in Table~\ref{sm:six-small-triples-table}. The table lists all further nonzero common extensions of each triple with modulus at most $R=10^{40}$.
\end{lemma}
\begin{proof}
Write $K=\mathbb Q(t)$ with $t=\sqrt{-d}$, where $d$ is positive and squarefree. Represent an algebraic integer by $(u+vt)/2$. The integers $u,v$ are both even if $d\not\equiv3\pmod4$, and have the same parity otherwise. Small nonzero elements satisfy $0<u^2+dv^2<36$.

If $d\ge36$, all such elements are rational and belong to $\{-2,-1,1,2\}$. A square root of a product of two of them plus one has modulus at most $\sqrt5<3$, so it too is rational. The only compatible pair is $\{-1,1\}$. Thus a small triple requires $d<36$.

For each squarefree $d<36$, form the graph on the admissible pairs with $0<u^2+dv^2<36$, joining distinct entries when their product plus one is a square. The required roots have norm at most $9$, so the square test uses the squares of all admissible pairs $(r,s)$ with $r^2+ds^2\le36$, including $(0,0)$. Testing all triangles gives $24$ triples in the ten orbits displayed in the table. Their orbit sizes are listed in its last column. No fourth small entry is adjacent to all three vertices of any triangle.

For the extension column we use the reversible descent of Lemma~\ref{sm:bs7-descent}. Number the three entries in each displayed triple as $1,2,3$. In the successive ten rows choose the pairs
\[
 (1,3),(1,2),(1,2),(1,3),(2,3),
 (1,3),(1,3),(1,2),(1,3),(1,2).
\]
Each chosen pair $a,b$ satisfies $ab+1\ne0$. Enumerate all its common extensions in the closed disk of radius $33$, retaining zero and any anchor coincidences. This initial list is obtained by enumerating roots $q$ with
\[
 N(q)\le\left\lfloor\sqrt{N(a)33^2}\right\rfloor+1,
 \qquad z=(q^2-1)/a,
\]
then checking integrality, $N(z)\le33^2$, and the other square condition. Starting from this list, repeatedly append both values
\[
 a+b+(2ab+1)z\pm2\sqrt{ab+1}\sqrt{az+1}\sqrt{bz+1}
\]
when their moduli are at most $R$, until no new value is obtained. The descent lemma proves completeness, since reversing a strictly decreasing descent from any extension in the disk stays in that disk.

Every step uses exact integer arithmetic. To specify the root test throughout this procedure, an integral $z=(u+vt)/2$ is a square precisely when $N(z)=h^2$ for a nonnegative integer $h$, and there are nonnegative integers $r,s$ satisfying
\[
 r^2=2h+u,\qquad ds^2=2h-u,\qquad rs=|v|,
\]
with $(r,s)$ admissible. The signs then construct an integral root $(r\pm st)/2$. This criterion includes zero.

The numbers of distinct pair extensions obtained, including zero and small entries, are respectively
\[
 61,89,64,105,71,112,112,57,83,55.
\]
Testing the third anchor condition and removing zero and the original three entries gives the extension column of the table. Thus the finite procedure proves all its assertions.
\end{proof}

\par\medskip
\noindent\begin{minipage}{\textwidth}
\makeatletter\def\@captype{table}\makeatother
\centering
\begin{tabular}{ccclc}
\toprule
$d$&$j$&Small triple&Extension&Orbit size\\
\midrule
1&1&$\{-2-2t,-2+2t,2\}$&$60$&2\\
1&2&$\{-2-t,-t,2-t\}$&$15t$&2\\
1&3&$\{-2-t,2t,2-t\}$&$-40t$&2\\
1&4&$\{-2,-1-2t,1\}$&$4+12t$&4\\
2&1&$\{-2,1-t,1+t\}$&$-24$&2\\
3&1&$\{(-3-t)/2,(-3+t)/2,1\}$&$8$&2\\
3&2&$\{(-3-t)/2,(1-3t)/2,1\}$&$-12+4t$&4\\
3&3&$\{(-1-t)/2,-t,(1-t)/2\}$&none&2\\
7&1&$\{(-5-t)/2,(-5+t)/2,1\}$&$24$&2\\
7&2&$\{(-5-t)/2,t,(5-t)/2\}$&$-32t$&2\\
\bottomrule
\end{tabular}
\caption{Small triples and all their further nonzero extensions of modulus at most $10^{40}$. Here $t=\sqrt{-d}$.}
\label{sm:six-small-triples-table}
\end{minipage}\par\medskip




\begin{lemma}\label{sm:five-small-anchor-alternative}
Fix three entries $a,b,s$ of a $D(1)$-tuple with
$1\le A\le B\le S<3$. Let $v,e$ be two further entries with
$1000\le V=|v|\le E=|e|$. Their coherent form satisfies
\[
 \begin{cases}
  5V^2<E<405V^2,&L_3(v,e)=0,\\
  E>V^5/100,&L_3(v,e)\ne0.
 \end{cases}
\]
In particular, $E>5V^2$ in either case.
\end{lemma}
\begin{proof}
First suppose $E\le4V$. Put $x=1/v$ and $y=1/e$.
The exact identity is
\[
 \frac{L_3}{\Delta_3\sqrt{ve}}
   =1-\sqrt{1+xy}+\frac{F_3}{abs}.
\]
Here $|F_3|<V^{-3}$, since
$(1000/999)^3<4$. Rationalization gives
\[
 |1-\sqrt{1+xy}|>\frac1{3VE}
       \ge\frac1{12V^2}>\frac1{V^3}\ge\frac{|F_3|}{ABS}.
\]
Thus $L_3\ne0$. But $V/2<Z<3V$, so
Lemma~\ref{lem:general-three} implies
\[
 |L_3|<\frac{270}{V}+\frac{27}{V^2}<1,
\]
contrary to integrality. Therefore $E>4V$.

If $L_3=0$, Lemma~\ref{lem:separated-three} gives
$5ABS V^2<E<15ABS V^2$. Since $1\le ABS<27$, this is the
first alternative. Otherwise the same lemma gives
\[
 1\le|L_3|<\frac{270}{V}+\frac{9Z}{4V^3}.
\]
The first term is less than $1/2$. Hence $Z>2V^3/9$ and
\[
 E>\frac{4V^5}{81}-\frac1V>\frac{V^5}{100}.
\]
Finally $V^5/100>5V^2$ for $V\ge1000$.
\end{proof}

\begin{proposition}\label{sm:five-three-small}
A $D(1)$-tuple in an imaginary quadratic integer ring containing
three entries of modulus less than $3$ has at most five entries.
\end{proposition}
\begin{proof}
Put $R=10^{40}$ and fix the three small entries, with ordered
moduli $1\le A\le B\le S<3$.
Lemma~\ref{sm:six-small-triple-classification} and
Table~\ref{sm:six-small-triples-table} give at most one further entry
of modulus at most $R$. If such an entry exists, its modulus is
less than $85$. These are complete closed-disk extension lists,
obtained by the reversible descent of Lemma~\ref{sm:bs7-descent}.

We record the approximation bounds needed below. For a further
pair of ordered moduli $R\le V\le E$, the two small anchors give
\begin{equation}\label{sm:five-small-upper}
 E<10^{60}V^7.
\end{equation}
Indeed, in Corollary~\ref{cor:ordinary-height-envelope} we may take
$h=3$, $r=-2$, and $u=4$. The squared small moduli are integers at most $8$. The necessary comparisons are
\[
 42^2\cdot8^3<10^6,\qquad 253\cdot8<10^4,\qquad
 \frac2{15}>10^{-2}.
\]
With $v=\log_{10}V\ge40$, that corollary consequently gives
\[
 \log_{10}E<7v+48+\frac{224}{v-5}+\frac25<7v+60.
\]

The three small anchors also forbid the closed interval
\begin{equation}\label{sm:five-three-small-window}
 10^{40}V^5\le E\le10^{-22}V^{11}.
\end{equation}
Apply the normalized three-root estimate
in Lemma~\ref{lem:ordinary-three-normalized}. Its order-two forbidden
interval contains~\eqref{sm:five-three-small-window}. This follows from
$A\ge1$, $B,S<3$, and the exact comparisons
\[
 \frac{512}{81}\,608^4\,3^{22}<10^{40},\qquad
 2\cdot28^4\left(\frac{10}{81}\right)^2 3^{14}<10^{22}.
\]
The estimate applies since $V\ge10^{40}>4S$ and $AV>1000$.

If the extension of modulus at most $R$ exists, use it as a fourth
anchor of modulus $T<85$. Lemma~\ref{sm:lem-four-root-k1} additionally forbids
\begin{equation}\label{sm:five-four-fixed-window}
 10^{38}V^3\le E\le10^{-29}V^7.
\end{equation}
Its hypotheses hold because $V\ge R>4T$ and $AV>1000$.
The displayed endpoints follow from
\[
 2\cdot10300^2\cdot3^{12}\cdot85^{12}<10^{38},\qquad
 2\cdot193^2\cdot3^8\cdot85^{10}<10^{29}.
\]

Suppose now that three further entries exist. First consider the
case where all three moduli satisfy $R<U\le V\le W$.
Apply Lemma~\ref{sm:five-small-anchor-alternative} to their three pairs.
If $L_3(u,v)\ne0$, then
\[
 V>U^5/100,\qquad W>5V^2>U^{10}/2000.
\]
If $L_3(u,v)=0$ and $L_3(v,w)\ne0$, then
$V>5U^2$ and $W>V^5/100>U^{10}/2000$ as well.
Either conclusion contradicts~\eqref{sm:five-small-upper}, since
$U^3>10^{120}>2000\cdot10^{60}$.

If both consecutive forms vanish, their quadratic bounds give
\[
 5U^2<V<405U^2,\qquad
 125U^4<W<405^3U^4.
\]
As $125U^4>405U^2$, the form $L_3(u,w)$ cannot vanish.
Its nonzero alternative gives $W>U^5/100$. Hence
$U<100\cdot405^3$, contrary to $U>R$.

It remains to consider the unique extension of modulus at most $R$,
whose modulus is $T<85$, together with two other entries of moduli
$R<V\le W$.
Lemma~\ref{lem:four-anchor-nonvanishing} applies because
\[
 V^2>R^2>32\cdot3\cdot3\cdot85\ge\frac{32BST}{A^3}.
\]
At least one coherent form on $(a,b,s)$ or $(a,b,t)$ is therefore
nonzero. Lemma~\ref{lem:general-three} gives
\[
 1<\frac{108375}{Z}+\frac{255Z}{V^3},\qquad Z^2=|vw+1|.
\]
Since $Z>V/2$ and $216750/V<1/2$, we obtain $Z>V^3/510$.
Consequently
\[
 W>\frac{V^5}{260100}-\frac1V
   >\frac{V^5}{10^6}>10^{38}V^3.
\]
The intervals~\eqref{sm:five-four-fixed-window} and
\eqref{sm:five-three-small-window} overlap because $V^2>10^{69}$.
Their union forbids every value from $10^{38}V^3$ to
$10^{-22}V^{11}$. But~\eqref{sm:five-small-upper} gives
$W<10^{60}V^7<10^{-22}V^{11}$, since $V^4>10^{82}$.
This is the final contradiction.
\end{proof}


\begin{proof}[Proof of Proposition~\ref{sm:main}]
Choose six distinct entries, including two of the small entries.
Lemma~\ref{sm:six-small-triple-classification} shows that at most three
are small. If exactly two are small, Proposition~\ref{sm:five-two-small-four-large}
excludes the four large entries. If three are small,
Proposition~\ref{sm:five-three-small} excludes the three remaining entries.
\end{proof}

\section{Reduction of the first three entries}
\label{sec:pair-reduction}

We now treat an ordered sextuple with
\[
 1\le A<3\le B\le S\le T\le C\le D.
\]
The preceding small-entry and all-large arguments reduce the proof to this
case. All norms below are the field norms, so \(N(z)=|z|^2\) is a
nonnegative rational integer.

We use the interval calculation of Section~\ref{sec:common-intervals}.
For the gap and the residual of $N_3$ on $b,s$, put
\[
 \mu=\frac{51051}{400000},\qquad \kappa=\frac{3535}{409600},\qquad
 v=B/S,\qquad d_0=|b-s|/S,
\]
\[
 \Phi=\max(d_0^6,vd_0^5),\qquad R_{bs}=\kappa\Phi.
\]
The coefficient $M_{bs}=4143$ is sufficient in both applications below.

\subsection{The five regions with an unbounded second entry}

The following elementary phase estimate keeps the two distances in the
gap and in the relabeled residual coupled. Put
\[
 \mathcal Q=\frac{|(a-b)(a-s)(b-s)|}{BS^2}.
\]
If \(B\ge qA\), \(q\ge1\), then
\begin{equation}\label{eq:pr-phase}
 \mathcal Q^2\Phi\le
 F(q):=\max_{0\le x\le2}(4-x^2)^4(1+q^{-2}+x/q)^2.
\end{equation}
Indeed, writing \(u=A/B\), \(v=B/S\), \(d_0=|b-s|/S\), the identity
\((a-b)(a-s)=a^2-a(b+s)+bs\) and the parallelogram identity give
\[
 \mathcal Q\le d_0\{1+u^2v+u\sqrt{2+2v^2-d_0^2}\}
 \le d_0\{1+q^{-2}+q^{-1}\sqrt{4-d_0^2}\}.
\]
For \(d_0\ge v\), substitute \(x=\sqrt{4-d_0^2}\).
For \(d_0<v\le1\), the same expression gives
\(\mathcal Q^2\Phi\le(1+1/q)^4\le16<F(q)\).
The unique maximizing interior point in \eqref{eq:pr-phase} is the
positive zero of
\begin{equation}\label{eq:pr-phase-critical}
 5x^2/q+4(1+q^{-2})x-4/q=0.
\end{equation}
The derivative changes sign there from positive to negative.

\begin{proposition}\label{prop:pr-H2}
In the one-small-entry branch, every hypothetical sextuple satisfies
\[
 N(a)=1,\ B<8,\qquad\hbox{or}\qquad N(a)=2,\ B<4.
\]
\end{proposition}
\begin{proof}
Since \(N(a)\) is an integer between one and eight, it is enough to
exclude the five rows of Table~\ref{tab:pr-five-regions}. In each row
\(B\ge m\), \(A^2\ge n\), \(A\le U\), and \(q=m/U\).
Lemma~\ref{lem:five-frontier-gaps} gives \(T>tABS\).
Consequently
\[
 AT>tnm^2\ge54,
\]
so Lemma~\ref{lem:retained-gap} applies with \(\lambda=19/200\).
In particular it supplies the size assumptions for all approximation
estimates and the starting lower bound. Put \(c_*=(9999/10000)(200/19)^2\),
\[
 Q_0=2(1+1/q)^2,\quad
 H=c_*t^5(1-\mu Q_0/m^2)^2.
\]
The triangle inequalities give \(\mathcal Q\le Q_0\).
Weakening \(\mu/(A^2B^2)\) to \(\mu/m^2\) therefore gives
\[
 C>\frac{H}{\mathcal Q^2}A^7B^5S^3.
\]
For every row of Table~\ref{tab:pr-five-regions}, exact substitution
shows \(H/Q_0^2>G\) and \(H/(\kappa F_p)>\rho\), where
\(F(q)<F_p\). This establishes both bounds in
\eqref{eq:pr-coarse-gap}.

\begin{table}[htbp]
\centering\small
\caption{Five excluded regions and their gap bounds.}
\label{tab:pr-five-regions}
\begin{tabular}{c|rrrrrrrrr}
\(N(a)\)&\(m\)&\(n\)&\(U\)&\(t\)&\(q\)&\(F_p\)&\(G\)&\(\rho\)&\(R_p\)\\\hline
1&8&1&1&27/8&8&269&7495&20687&7/50\\
2&4&2&3/2&3&8/3&368&1771&7973&7/20\\
3&3&3&7/4&11/4&12/7&540&598&3226&18/25\\
4&3&4&2&29/10&3/2&637&624&3507&47/50\\
5--8&3&5&17/6&3&18/17&1121&375&2217&201/100
\end{tabular}
\end{table}
The finite-family residuals are
\begin{equation*}
 R_k=R_*\frac{101}{200}
             \left(\frac{1+1/q}{2}\right)^{2k}.
\end{equation*}
The coefficient bounds and overlap margins are given in
Table~\ref{tab:pr-five-coefficients}. If a root-column monomial after division by its
anchor is \(n a^r b^s c^\ell\), then
\[
 r,s\ge0,\quad0\le\ell\le k,\quad r+s=\ell+P-k-2,
 \qquad
 \frac{|n a^r b^s c^\ell|}{B^{P-2}C^k}
 \le |n|q^{-r}1000^{\ell-k}.
\]
Thus a strict upper bound for
\begin{equation}\label{eq:pr-Mq}
 \frac{101}{100g}\max_{\text{rows}}
       \sum_{\text{two root columns}}\sum_{r,s,\ell}
                    |n|q^{-r}1000^{\ell-k}
\end{equation}
is a valid \(M_k\). The explicit polynomial rows constructed earlier,
applied to \eqref{eq:pr-Mq}, give the displayed integers.
The adaptive residual is the maximum of its four explicit positive
polynomials at $x=1/q$. Their maximum is strictly below the displayed
\(R_p\).
\begin{table}[htbp]
\centering\small
\caption{Coefficient bounds for the five regions. The last column is a
strict lower bound for every overlap quotient.}
\label{tab:pr-five-coefficients}
\begin{tabular}{c|rrrrrr}
\(N(a)\)&\(M_2\)&\(M_n\)&\(M_3\)&\(M_4\)&\(M_5\)&margin\\\hline
1&36&282&145&4605&38390&1.260\\
2&69&697&326&14771&171856&2.329\\
3&109&1306&663&34919&482131&1.166\\
4&130&1652&863&48019&701908&5.459\\
5--8&220&3553&1944&134712&2400382&1.321
\end{tabular}
\end{table}

Here and below terminating decimals denote the corresponding exact
rational numbers. All the numerical verifications in these two tables
can be performed as follows. For \(F_p\), bisect \([0,2]\) thirty times
using the sign of \eqref{eq:pr-phase-critical}, obtaining \([l,r]\).
Then evaluate
\((4-l^2)^4(1+q^{-2}+r/q)^2<F_p\).
Evaluate \eqref{eq:pr-Mq}, the four residual polynomials, and the two
bounds for \(G,\rho\) by rational arithmetic. In each row of
Table~\ref{tab:pr-overlaps}, substitute \(A^2=n\) and \(B=S=m\).
All remaining exponents are nonnegative and every exponent of \(A\)
is even, so this gives a lower bound throughout the region. The
resulting fifteen quotients are all larger than the last column of
Table~\ref{tab:pr-five-coefficients}. Finally
\((G/105)n^4m^4>1\), proving the infinite-tail condition. The common
interval calculation gives a contradiction in every region.
\end{proof}

\subsection{The field and the finite first-pair list}

\begin{lemma}\label{lem:pr-one-three}
Every common extension of $\{1,3\}$ in the full integer ring of an
imaginary quadratic field is rational. A $D(1)$-tuple containing
$1,3$ and a negative entry has at most four elements.
\end{lemma}
\begin{proof}
First every common extension is rational. Otherwise choose roots
\(r=(e+f\sqrt{-\Delta_K})/2\) and
\(u=(g+h\sqrt{-\Delta_K})/2\) of \(c+1\) and \(3c+1\), respectively.
Nonreality makes \(e,f,g,h\) nonzero. From \(3r^2-u^2=2\), put
\[
 X=3e^2-g^2,\quad V=h^2-3f^2,\quad k=(eh-fg)/2.
\]
If $\Delta_K$ is odd, then $e\equiv f$ and $g\equiv h\pmod2$.
If $\Delta_K$ is even, then $e$ and $g$ are even.
In either case $k$ is an integer. The imaginary and real parts
give $3ef=gh$ and $X+\Delta_K V=8$. Multiplication gives
\(XV=12k^2\). Moreover $V=3f^2X/g^2$, so $X,V>0$. Also $k\ne0$, since $k=0$ would give $g^2=3e^2$, which is impossible for nonzero integers. Hence
\[
 64=(X+\Delta_K V)^2\ge4\Delta_K XV=48\Delta_K k^2\ge144,
\]
a contradiction.

Write a negative rational extension as \(-x\), where \(x\) is a
positive integer. Purely imaginary integral roots in
\(\mathbb Q(\sqrt{-d})\) are integer multiples of \(\sqrt{-d}\),
including in the half-integral full rings. Thus
\(x-1=du^2\) and \(3x-1=dv^2\), where \(u=0\) is allowed,
and \(d\mid2\).
The case \(d=1\) is impossible modulo three, so \(d=2\) and \(x\)
is odd. A positive extra entry \(p\) satisfies
\(p+1=r^2\), \(3p+1=s^2\) with integer roots. Reduction modulo four
forces \(r\) odd, hence \(8\mid p\). Such a positive entry cannot coexist with $-x$, because the equation $px-1=2w^2$ has an odd left side. If negative entries occur, all extra entries are therefore
negative odd integers. At most two can coexist, since among three
odd absolute values two have the same residue modulo four and their
product plus one is congruent to two.
\end{proof}

\begin{lemma}\label{lem:pr-field}
Under Proposition~\ref{prop:pr-H2}, either the first pair is
$(1,3)$ up to simultaneous negation, or the field has absolute
discriminant $\Delta_K\le24$.
\end{lemma}
\begin{proof}
For a nonreal algebraic integer, \(N(z)\ge\Delta_K/4\).
If \(\Delta_K\ge25\), the condition \(N(a)\le2\) therefore makes
\(a\) rational. Simultaneous negation allows \(a=1\), and then
\(B<8\). For \(r^2=b+1\),
\(N(r)=|b+1|<9\). If \(\Delta_K\ge36\), this forces \(r\) rational.
The only absolute negative fundamental discriminants between 25 and 35
are 31 and 35. At 35 a nonreal integer has integral norm at least nine.
At 31 the only nonreal roots of norm below nine are
$(\pm1\pm\sqrt{-31})/2$. They give \(N(r^2-1)=80\), contrary to
\(N(b)<64\). Thus \(r\) is rational in every case. Now
$b=r^2-1$, $r\in\mathbb Z$, and $3\le|b|<8$ force $b=3$,
which is the stated alternative.
\end{proof}

We enumerate the remaining pairs in the fields allowed by
Lemma~\ref{lem:pr-field}. For squarefree
\[
 d\in\{1,2,3,5,6,7,11,15,19,23\},
\]
write \(\omega=(1+\sqrt{-d})/2\) when \(d\equiv3\pmod4\), and
\(\omega=\sqrt{-d}\) otherwise. Put
\(\omega^2=\tau_K\omega-\nu_K\), \(\Delta_K=4\nu_K-\tau_K^2\).
The exact arithmetic is
\begin{equation}\label{eq:pr-ring-arithmetic}
\begin{aligned}
 (x,y)(v,w)&=(xv-\nu_Kyw,xw+yv+\tau_Kyw),\\
 N(x,y)&=x^2+\tau_Kxy+\nu_Ky^2,\qquad
 \overline{(x,y)}=(x+\tau_Ky,-y).
\end{aligned}
\end{equation}
All elements of norm at most \(L\) are given by
\begin{equation}\label{eq:pr-disk}
\begin{aligned}
 |y|&\le\left\lfloor\sqrt{4L/\Delta_K}\right\rfloor,&
 |e|&\le\left\lfloor\sqrt{4L-\Delta_Ky^2}\right\rfloor,\\
 e&\equiv\tau_Ky\pmod2,&x&=(e-\tau_Ky)/2.
\end{aligned}
\end{equation}
This follows from \(4N(x,y)=(2x+\tau_Ky)^2+\Delta_Ky^2\), and is both
necessary and sufficient.

Enumerate every \(a\) of norm one or two and every \(r\) of norm at
most eight. Form
\[
 b=\overline a(r^2-1)/N(a),
\]
retaining it only when both coordinates are integers and
\[
 9\le N(b)\le63\quad(N(a)=1),\qquad
 9\le N(b)\le15\quad(N(a)=2).
\]
The root range is exhaustive because $N(r)=|ab+1|\le AB+1<9$ in the first case and $AB+1\le\sqrt{30}+1<7$ in the second.
Set aside $(a,b)=(1,3),(-1,-3)$, which will be treated separately. Quotient only by simultaneous
negation and conjugation of both coordinates. For an explicit
representative rule, minimize lexicographically
\[
 \bigl(2a_x+\tau_Ka_y<0,\ a_y<0,\ 2b_x+\tau_Kb_y<0,\ b_y<0,\ a_x,a_y,b_x,b_y\bigr)
\]
in each orbit, with false preceding true. The result is exactly the
forty pairs in Table~\ref{tab:pr-forty-pairs}. The number of orbits
and the number before these symmetries, respectively, are
\[
\begin{array}{c|rrrrrrrrrr}
 d&1&2&3&5&6&7&11&15&19&23\\\hline
 \text{orbits}&9&3&15&2&2&2&2&2&2&1\\
 \text{pairs}&30&10&58&6&6&8&8&8&8&4.
\end{array}
\]
Their totals are forty and 146. The symmetries preserve every square
condition. No arbitrary unit multiplication is used.

\begingroup\small
\setlength{\tabcolsep}{5pt}
\begin{longtable}{r|r|c|c|r|r}
\caption{The forty first-pair representatives. Coordinates mean $x+y\omega$.}\label{tab:pr-forty-pairs}\\
$i$&$d$&$a$&$b$&$N(a)$&$N(b)$\\\hline\endfirsthead
$i$&$d$&$a$&$b$&$N(a)$&$N(b)$\\\hline\endhead
0&1&$(0,1)$&$(0,-3)$&1&9\\
1&1&$(0,1)$&$(0,5)$&1&25\\
2&1&$(0,1)$&$(4,-2)$&1&20\\
3&1&$(0,1)$&$(4,4)$&1&32\\
4&1&$(1,0)$&$(-5,0)$&1&25\\
5&1&$(1,0)$&$(-4,4)$&1&32\\
6&1&$(1,0)$&$(2,4)$&1&20\\
7&2&$(1,0)$&$(-3,0)$&1&9\\
8&2&$(1,0)$&$(-2,2)$&1&12\\
9&2&$(1,0)$&$(1,4)$&1&33\\
10&3&$(0,1)$&$(-9,6)$&1&63\\
11&3&$(0,1)$&$(-6,9)$&1&63\\
12&3&$(0,1)$&$(-5,1)$&1&21\\
13&3&$(0,1)$&$(-4,-4)$&1&48\\
14&3&$(0,1)$&$(-4,4)$&1&16\\
15&3&$(0,1)$&$(-1,5)$&1&21\\
16&3&$(0,1)$&$(2,-7)$&1&39\\
17&3&$(0,1)$&$(3,-3)$&1&9\\
18&3&$(0,1)$&$(4,4)$&1&48\\
19&3&$(0,1)$&$(7,-2)$&1&39\\
20&3&$(1,0)$&$(-9,3)$&1&63\\
21&3&$(1,0)$&$(-5,4)$&1&21\\
22&3&$(1,0)$&$(-4,0)$&1&16\\
23&3&$(1,0)$&$(-4,8)$&1&48\\
24&3&$(1,0)$&$(2,5)$&1&39\\
25&5&$(1,0)$&$(-6,0)$&1&36\\
26&5&$(1,0)$&$(-5,2)$&1&45\\
27&6&$(1,0)$&$(-7,0)$&1&49\\
28&6&$(1,0)$&$(-6,2)$&1&60\\
29&7&$(1,0)$&$(-2,3)$&1&16\\
30&7&$(1,0)$&$(1,5)$&1&56\\
31&11&$(1,0)$&$(-4,1)$&1&15\\
32&11&$(1,0)$&$(-3,3)$&1&27\\
33&15&$(1,0)$&$(-5,1)$&1&24\\
34&15&$(1,0)$&$(-4,3)$&1&40\\
35&19&$(1,0)$&$(-6,1)$&1&35\\
36&19&$(1,0)$&$(-5,3)$&1&55\\
37&23&$(1,0)$&$(-7,1)$&1&48\\
38&1&$(1,1)$&$(-1,-3)$&2&10\\
39&1&$(1,1)$&$(3,1)$&2&10\\
\end{longtable}
\endgroup
\begingroup\small
\setlength{\tabcolsep}{4pt}
\begin{longtable}{r|rrrrr}
\caption{Coefficient bounds for the forty fixed pairs.}\label{tab:pr-forty-coefficients}\\
$i$&$M_2$&$M_n$&$M_3$&$M_4$&$M_5$\\\hline\endfirsthead
$i$&$M_2$&$M_n$&$M_3$&$M_4$&$M_5$\\\hline\endhead
0&62&608&280&12222&136530\\
1&17&77&44&808&4011\\
2&37&281&147&4610&38256\\
3&21&110&64&1402&8160\\
4&44&379&180&6340&61370\\
5&37&293&150&4791&40361\\
6&23&129&77&1798&11204\\
7&62&608&280&12222&136530\\
8&45&377&183&6284&60438\\
9&25&158&95&2471&16976\\
10&24&156&94&2424&16601\\
11&21&119&70&1599&9691\\
12&34&252&136&4117&32986\\
13&37&286&147&4661&38983\\
14&22&121&72&1618&9751\\
15&17&73&41&741&3580\\
16&39&318&159&5209&45613\\
17&48&401&193&6835&66703\\
18&20&106&62&1336&7678\\
19&24&152&91&2326&15710\\
20&36&276&143&4502&37275\\
21&40&323&162&5308&46658\\
22&51&456&205&8194&83910\\
23&26&173&103&2780&19826\\
24&21&113&66&1460&8607\\
25&41&333&164&5461&49340\\
26&35&273&142&4450&36676\\
27&38&303&153&4954&42245\\
28&34&259&137&4205&34076\\
29&32&224&127&3666&28238\\
30&23&137&81&1973&12709\\
31&50&440&201&7791&78839\\
32&32&230&127&3748&29208\\
33&44&367&176&6049&57935\\
34&32&229&126&3718&28965\\
35&40&324&161&5305&46974\\
36&31&226&124&3660&28401\\
37&37&296&151&4834&40902\\
38&77&803&387&17946&216405\\
39&19&38&41&370&985\\
\end{longtable}
\endgroup

For the rational pair $(a,b)=(1,3)$, the same coefficient calculation gives
\[
 (M_2,M_n,M_3,M_4,M_5)=(12,41,25,300,1100).
\]
These constants are independent of the imaginary quadratic field.


\subsection{A common bound \texorpdfstring{$S<15$}{S<15}}

\begin{proposition}\label{prop:pr-S15}
No ordered sextuple whose first pair is $(1,3)$ or one of the pairs in
Table~\ref{tab:pr-forty-pairs} has $S\ge15$.
\end{proposition}
\begin{proof}
Suppose that $S\ge15$. Lemma~\ref{lem:tail-irregular} gives
$T>2BS\ge90$. We first prove the common estimate
\begin{equation}\label{eq:pr-fixed-fourth-gap}
 T>\tau ABS,\qquad \tau=\frac{59}{20}.
\end{equation}
For the edges $12,13,14,23,24,34$ of $(a,b,s,t)$, use the bounds
\[
 r=\left(\frac13,\frac1{15},\frac1{90},
          \frac1{45},\frac1{270},\frac1{1350}\right),\qquad
 k=\frac1{10000}(5506,5087,5014,5029,5005,5001).
\]
For every corresponding pair $r_e,k_e$, we have
$1/2<k_e<1$ and $r_ek_e^2<2k_e-1$.
Hence, for $z_e=1/(a_i a_j)$ and
$\delta_e=\sqrt{1+z_e}-1$ on the branch with positive real part,
\[
 |\delta_e|\le k_e|z_e|,\qquad
 |\delta_e-z_e/2|\le(k_e-1/2)|z_e|.
\]
The first inequality follows by rationalization and the second from
$\delta-z/2=-\delta^2/2$.

Let $M$ be the integral matrix of order four defined in the proof of
Lemma~\ref{lem:five-frontier-gaps}, with $(a_1,a_2,a_3,a_4)=(a,b,s,t)$.
It is nonsingular by the even-dimensional congruence used there.
Apply the determinant expansion
\eqref{eq:integrated-7}--\eqref{eq:integrated-8}. The coefficient
bounding the term $T/(ABS)$ is
$\alpha=2k_{12}k_{13}k_{23}$. For complementary pairs $\{i,j\}$
and $\{k,l\}$, the coefficient bounding $1/(A_iA_j)$ is
\begin{equation*}
 c_{ij}=2\{k_{ij}(k_{ik}k_{jl}+k_{il}k_{jk})
                  +k_{ij}^2k_{kl}\}-\frac12.
\end{equation*}
Indeed, retain the signed leading cubic in
\eqref{eq:integrated-8} and bound each cubic remainder by the
product of its three $k$'s minus $1/8$, counting repeated edges
with their multiplicities. Put
\[
 Q_4=(k_{12}k_{34}+k_{13}k_{24}+k_{14}k_{23})^2,\qquad
 W=\frac1{4050}.
\]
Since $A<3$, $B\ge3$, $S\ge15$ and $T>2BS$, the three remaining
cubic terms $A/(BST)$, $B/(AST)$ and $S/(ABT)$ are at most
$3W$, $1/450$ and $1/18$, respectively. The quartic factor
$1/(ABST)$ is at most $W$. Consequently
\[
 1\le|\det M|\le\alpha\frac{T}{ABS}+\beta,
\]
where
\[
 \beta=6k_{23}k_{24}k_{34}W
       +\frac{2k_{13}k_{14}k_{34}}{450}
       +\frac{2k_{12}k_{14}k_{24}}{18}
       +\sum_{i<j}c_{ij}r_{ij}+Q_4W.
\]
Substitution of the six rational $k_e$ gives
$\alpha<141/500$ and $\beta<21/125$. If $T\le(59/20)ABS$,
these inequalities would give
\[
 1<\frac{141}{500}\frac{59}{20}+\frac{21}{125}
   =\frac{9999}{10000},
\]
a contradiction. This proves~\eqref{eq:pr-fixed-fourth-gap}.
Since $AT>90$, Lemma~\ref{lem:retained-gap} applies with
$\lambda=1/10$.

Write $n_a=N(a)$, $n_b=N(b)$ and $n_{ab}=N(a-b)$. For a nonnegative
rational $x$, define
\[
 \lfloor\sqrt{x}\rfloor_j=10^{-j}\lfloor10^j\sqrt{x}\rfloor,
 \qquad
 \lceil\sqrt{x}\rceil_j=10^{-j}\lceil10^j\sqrt{x}\rceil.
\]
These bounds use only integer square roots and comparisons. Put
\[
 A_1=\lceil\sqrt{n_a}\rceil_8,\quad
 B_0=\lfloor\sqrt{n_b}\rfloor_8,\quad
 B_1=\lceil\sqrt{n_b}\rceil_8,\quad c_0=\frac{9999}{100},
\]
and define
\begin{align*}
 d_1&=\lceil\sqrt{n_{ab}/n_b}\rceil_8,&
 Q_0&=d_1(1+A_1/15)(1+B_1/15),\\
 K&=c_0n_a(1/Q_0-\mu/n_b)^2,&G&=K\tau^5,\\
 H&=c_0n_a\tau^5(1-\mu Q_0/n_b)^2,&
 F_0&=(n_{ab}/n_b)(1+A_1/15)^2(1+B_1/15)^8,\\
 \rho&=H/(\kappa F_0).&&
\end{align*}
For all forty-one pairs, $1/Q_0>\mu/n_b$ and $K>7/4$.
The retained gap, after replacing $\mu/(A^2B^2)$ by $\mu/B^2$, gives
\[
 C>K T^5/S^2>\frac74T^5/S^2,\qquad C>GB^5S^3.
\]
Moreover, $\mathcal Q\le Q_0$ and $\mathcal Q^2\Phi\le F_0$.
For the second inequality, use $|a-s|/S\le1+A/S$,
$|b-s|/S\le1+B/S$ and $\Phi\le(1+B/S)^6$.
Thus the pointwise coefficient of $B^5S^3$, divided by $R_{bs}$,
is at least $\rho$. All other positive powers of $A$ in the
interval endpoints may be replaced by one.

We evaluate the finite-family coefficients at $a,b$ before taking
absolute values. If the two root columns after division by their
anchors are $\sum_{\ell=0}^k p_{j,\ell}(a,b)c^\ell$, put
$C_0=\max(1000,200B)$. Any integer strictly larger than
\begin{equation*}
 \frac{101}{100g}\max_{\mathrm{rows}}
 \sum_{j=1}^2\sum_{\ell=0}^k
 \frac{|p_{j,\ell}(a,b)|}{B^{P-2}C_0^{k-\ell}}
\end{equation*}
is a valid coefficient bound, since $C\ge C_0$.
To reproduce the displayed integers, write
$n_p=N(p_{j,\ell}(a,b))$ and $r=k-\ell$, and replace each summand
before $101/(100g)$ by
\[
 \begin{cases}
 \lceil\sqrt{n_p/n_b^{P-2}}\rceil_{10}/1000^r,&n_b\le25,\\
 \lceil\sqrt{n_p/n_b^{P-2+r}}\rceil_{10}/200^r,&n_b>25.
 \end{cases}
\]
Take the integer part of the largest rational row sum, plus one.
Every polynomial evaluation uses~\eqref{eq:pr-ring-arithmetic}.

For every pair, $|a-b|\ge B-A>A$. This also holds for the two
norm-two pairs, since $B\ge3>2\sqrt2$. Hence the residual constants are
\begin{equation*}
 R_k=R_*\frac{101}{200}
                 \left(\frac{n_{ab}}{4n_b}\right)^k.
\end{equation*}
For the adaptive residual use $q=B_0/A_1>2$ and $x=1/q$.
Its largest residual bound is
\begin{equation*}
 R_p=\frac{15+80x+150x^2+280x^3+35x^4}{256}
 +\frac{21(885+1144x+435x^2+2634x^3+170x^4)}{1024\cdot999}.
\end{equation*}
The comparison with the other three residual polynomials follows
directly from $0\le x\le1$ and $x^5\le x^4\le x^3$.

Use Table~\ref{tab:pr-overlaps} with every $A$ replaced by one.
The first two numerators improve from $1024,256$ to $9604,784$,
respectively, because $C>(7/4)T^5/S^2$.
Substitute $B^2=n_b$ and $S^2=225$ in the remaining nonnegative
powers. For each of the forty-one pairs, all fifteen rational
quotients are strictly larger than $59/50$.
These checks use only the displayed formulas and family polynomials.
The infinite-tail condition holds as well, since the same constants give
\[
 \frac{R_0}{H_0}\ge\frac{C}{105B^4}
 >\frac{GBS^3}{105}\ge\frac{GB_0\,15^3}{105}>1.
\]
The common interval argument gives the contradiction.
\end{proof}

\subsection{The remaining third entries}

The remaining third entries are found by a bounded enumeration.
For each pair of
Table~\ref{tab:pr-forty-pairs}, enumerate \(s\) by
\eqref{eq:pr-disk} with \(L=224\), retain \(N(s)\ge N(b)\) and
\(s\ne a,b\), and test \(as+1\) and \(bs+1\) for being squares.
The norm cutoff is exact because Proposition~\ref{prop:pr-S15}
gives \(S<15\) and the norm is integral.
An exact full-ring square test, including zero, is as follows.
For \(z=(z_x,z_y)\), put \(p=2z_x+\tau_Kz_y\), \(q=z_y\).
Require \(N(z)=n^2\) with \(n\ge0\) integral, then require integers
\(e\ge0,f\) satisfying
\[
 e^2=2n+p,\quad \Delta_Kf^2=2n-p,\quad ef=q,
 \quad e\equiv\tau_Kf\pmod2.
\]
If they exist, \(((e-\tau_Kf)/2,f)\) is an integral root, and conversely
every root gives these equations after changing its sign if necessary.
When \(e=0\), either sign of \(f\) gives the same square. A final
multiplication verifies every retained witness.

This calculation gives 57 triples. One is
\((1,-3,3)\), already excluded by Lemma~\ref{lem:pr-one-three}.
There remain exactly 56 triples, retaining all equal-modulus cases and
without any additional symmetry on the third entry. They are listed, together with their cutoffs $N_0$, in
Table~\ref{tab:fixed56}.
It remains to exclude these fifty-six triples and the pair $(1,3)$.

\section{Completion for one small entry}
\label{sec:one-small-completion}

Proposition~\ref{prop:pr-S15} gives $S<15$. For the forty pairs in
Table~\ref{tab:pr-forty-pairs}, it leaves fifty-six triples.
By Lemma~\ref{lem:pr-one-three}, any sextuple containing the
additional pair $(1,3)$ consists of positive rational integers.
The conditions $3<s<15$, $s+1$ square and $3s+1$ square give $s=8$.
For each of these fifty-six triples and for $(1,3,8)$, the forbidden intervals
give a bound $N(t)<N_0$. We then contradict this bound using the growth of an
irregular triple, with a separate square calculation for the last two cases.

\subsection{A common fourth-entry cutoff}

Fix one of these triples $(a,b,s)$ and put
\[
 a_0=N(a),\quad b_0=N(b),\quad s_0=N(s),\qquad
 d_{xy}=N(x-y),\quad L=10^{10}.
\]
For a nonnegative rational number $v$ define rational enclosures
\[
 [v]_- =\frac{\lfloor L\sqrt v\rfloor}{L},\qquad
 [v]_+ =\frac{\lceil L\sqrt v\rceil}{L}.
\]
These definitions require only integer square comparisons. The lower
integer $n$ is characterized by $n^2\le L^2v<(n+1)^2$.
No floating-point comparison is used below. Define
\begin{align*}
 q_1&=\left[\frac{d_{ab}d_{as}d_{bs}}{b_0s_0^2}\right]_+,
 &K&=\frac{9999}{100}\,a_0
       \left(\frac1{q_1}-\frac{\mu}{b_0}\right)^2,\\
 \phi_1&=\left[\max\left\{\frac{d_{bs}^6}{s_0^6},
                  \frac{b_0d_{bs}^5}{s_0^6}\right\}\right]_+,
 &R_{bs}&=\frac{3535}{409600}\phi_1,
 \qquad \mu=\frac{51051}{400000}.
\end{align*}
In every row $1/q_1>\mu/b_0$ and $K>1$.
For $a_0=1$, Lemma~\ref{lem:tail-irregular} gives $T>2BS=2ABS$.
For $a_0=2$ we will impose $N(t)\ge4a_0b_0s_0+1$.
Thus, whenever $N(t)\ge1024$ and this additional condition holds,
the row $H=32$, $\lambda=1/10$ of Lemma~\ref{lem:retained-gap}
implies
\begin{equation}\label{eq:fixed56-gap}
 C>K\frac{T^5}{S^2},\qquad
 D>\frac{C^7}{S^2T^4}.
\end{equation}
Indeed, $\mathcal Q\le q_1$ and $\mu/(a_0b_0)\le\mu/b_0$,
so the coefficient of $T^5/S^2$ in Lemma~\ref{lem:retained-gap}
is at least $K$. The other enclosure satisfies $\Phi\le\phi_1$, so the
displayed $R_{bs}$ is a valid residual constant for $N_{3,bs}$.
The same lemma gives $C>\max(1000,200S,4T)$ and $D>1000C$,
which supply all approximation hypotheses in this section.

Use the constants $M_2,M_n,M_3,M_4,M_5$ in
Table~\ref{tab:pr-forty-coefficients}. For $(a,b)=(1,3)$, take
$12,41,25,300,1100$, respectively. Their normalizations are those of
$N_2,N_3,A_3,A_4,A_5$.
For $N_{3,bs}$ evaluate the same $N_3$ polynomials at $b,s$ and use
\eqref{eq:ff-fixedM} with $C_0=\max(1000,200S)$.
More explicitly, in the notation of that formula, replace every
summand $|p_{ij\ell}|C_0^{\ell-3}/S^7$ by
\[
 \begin{cases}
 [N(p_{ij\ell})/s_0^7]_+/1000^{3-\ell},&s_0\le25,\\
 [N(p_{ij\ell})/s_0^{10-\ell}]_+/200^{3-\ell},&s_0>25.
 \end{cases}
\]
Multiply the largest row sum by $101/400$ and take its floor plus one.
Denote this integer by $M_{bs}$.
This is the integer printed below. Evaluation precedes taking norms,
so cancellation between monomials of the same $c$-degree is preserved.

Since $|a-b|\ge B-A>A$ for all these pairs, the distance-sensitive
residual constants in the first five families are
\begin{equation*}
 R_k=R_*\frac{101}{200}\left(\frac{d_{ab}}{4b_0}\right)^k,
 \qquad (R_*,k)=
 (5/4,2),(35/32,3),(35/16,3),(63/8,4),(495/64,5).
\end{equation*}
They will be denoted by $R_2,R_n,R_3,R_4,R_5$, respectively.
For $P$ take $R_p$ from \eqref{eq:ff-adaptiveres} with
$q=[b_0/a_0]_-$, and take its coefficient bound $2400$.
Finally put $M_m=(8/9)608^2$ and $R_m=3920/81$.

Use the endpoints in Table~\ref{tab:pr-windows}, replacing $A$ by
one only in the upper endpoints. This discards positive powers of
$A\ge1$, so each resulting closed interval remains forbidden for
both possible values of $a_0$. The constants $M_{bs}$ and $R_{bs}$
are the specialized ones just defined.

Put $h=4001/3999$ and $x=N(t)$.
The next table gives a strict lower bound $c_ex^{p_e}$ for each
ratio of endpoints, using only \eqref{eq:fixed56-gap}.
For the first row the ratio is $(1-\varepsilon)C^7/(S^2T^4l_Q)$.
For every other row it is $u_{\mathrm{left}}/(hl_{\mathrm{right}})$.
Both terms in the maximum defining $P$ are checked separately.
\begin{center}\small
\begin{tabular}{c|c}
junction&$c_ex^{p_e}$\\\hline
start--$Q$&$(1-\varepsilon)K^4x^5/(12100s_0^5)$\\
$Q$--$N_2$&$K^2x^4/(36M_2^2b_0^4s_0^2h)$\\
$N_2$--$P$, first term&$K^2x^5/(16R_2^2\cdot2400^2b_0^3s_0^5h)$\\
$N_2$--$P$, second term&$K^2x^5/(16R_2^2\cdot2400^2b_0^5s_0^4h)$\\
$P$--$A_3$&$K^2x^5/(16R_p^2M_3^2b_0^6s_0^4h)$\\
$A_3$--$M$&$K^4x^{10}/(16R_3^2M_m^2b_0^6s_0^{12}h)$\\
$M$--$N_3$&$K^4x^{10}/(16R_m^2M_n^2b_0^9s_0^9h)$\\
$N_3$--$N_{3,bs}$&$K^4x^{10}/(16R_n^2M_{bs}^2b_0^3s_0^{11}h)$\\
$N_{3,bs}$--$A_4$&$K^2x^5/(16R_{bs}^2M_4^2b_0^5s_0^5h)$\\
$A_4$--$O_4$&$K^4x^{10}/(144R_4^2\cdot42^8b_0^{15}s_0^4h)$\\
$O_4$--$A_5$&$(2/5)^8K^4x^{10}/(4M_5^2b_0^{16}s_0^4h)$\\
$A_5$--$O_5$&$K^6x^{15}/(144R_5^2\cdot42^{10}b_0^{19}s_0^6h)$\\
$O_5$--$O_6$&$(2/5)^{10}K^6x^{15}/(36\cdot42^{12}b_0^{22}s_0^6h)$
\end{tabular}
\end{center}
For each row let $n_e$ be the least positive integer with $c_en_e^{p_e}>1$.
Doubling gives an upper bound, after which integer binary search
determines $n_e$. The final checks are $c_en_e^{p_e}>1$ and
$c_e(n_e-1)^{p_e}\le1$.
Let $n_\infty$ be defined in the same way by
\[
 \frac{K^2x^5}{105^2b_0^4s_0^2}>1,
\]
and define
\begin{equation*}
 N_0=\max\bigl(1024,n_\infty,\{n_e\}_e,\eta\bigr),\qquad
 \eta=\begin{cases}1,&a_0=1,\\4a_0b_0s_0+1,&a_0=2.\end{cases}
\end{equation*}
All numbers and exponents in this prescription are rational or integral.
The values in Table~\ref{tab:fixed56} result directly from these formulas.

If $N(t)\ge N_0$, every overlap is strict. The starting lower bound
and successive forbidden intervals force $D$ beyond the upper endpoint
of each of
\[
 Q,N_2,P,A_3,M,N_3,N_{3,bs},A_4,O_4,A_5,O_5,O_6,\ldots.
\]
With $A$ replaced by one, the squared ratio of the ordinary bases
is greater than $K^2x^5/(105^2b_0^4s_0^2)>1$.
The ordinary-tail argument of Section~\ref{sec:common-intervals}
therefore gives a contradiction and proves $N(t)<N_0$.

For $(a,b,s)=(1,3,8)$ the formulas give
\[
 q_1=\frac{1822916667}{5000000000},\qquad K>744,\qquad
 M_{bs}=34,\qquad N_0=1024.
\]
For clarity, the thirteen integers $n_e$, in the order of the overlap
table, are
\[
 3,23,96,101,108,179,131,9,34,76,123,29,438,
\]
and $n_\infty=14$. Thus $T<32$, whereas
Lemma~\ref{lem:tail-irregular} gives $T>2\cdot3\cdot8=48$.
This contradiction excludes the rational pair.

\subsection{The finite complement}

The fifty-six triples in Table~\ref{tab:fixed56} were enumerated at the
end of Section~\ref{sec:pair-reduction}. We now exclude $N(t)<N_0$.
Lemma~\ref{lem:tail-irregular} gives $T>2BS$. In thirty-eight rows,
$4b_0s_0\ge N_0$, so this already contradicts $N(t)<N_0$.

For the other rows, put
\[
 q=b+s,\qquad h=1+2bs,\qquad \delta=(b-s)^2-4,
\]
and define
\[
 f(X)=X^2-2(|h|-|q|)X+|\delta|+2|q|+1.
\]
The triple $(b,s,t)$ is irregular by Lemma~\ref{lem:tail-irregular},
so \eqref{eq:irregular-quadratic} gives $f(T)\ge0$.
This quadratic is convex. For $X\ge0$, the bounds
$|h|\ge2BS-1$, $|q|\le B+S$ and
$|\delta|\le(B+S)^2+4$ give
\[
 f(X)\le(B+S+X+1)^2+4-4BSX.
\]
The endpoint calculation in the proof of
Lemma~\ref{lem:irregular-growth} therefore gives $f(2BS)<0$.

For sixteen of the remaining rows, substitution of the displayed
entries and $N_0$ gives $f(\sqrt{N_0})<-100$. These are exact checks
using the rational enclosures defined above. Explicitly,
$[N(h)]_--[N(q)]_+>0$, and
\[
 \begin{split}
 f(\sqrt{N_0})\le{}&N_0
 -2\bigl([N(h)]_--[N(q)]_+\bigr)[N_0]_-\\
 &+[N(\delta)]_++2[N(q)]_++1<-100.
 \end{split}
\]
Convexity excludes $2BS<T<\sqrt{N_0}$ in each of these rows.

The two remaining triples are
$(a,b,s)=(1,-3,-2\pm2\sqrt{-2})$, both with $N_0=1024$.
Here $N(q)=33$, $N(h)=457$ and $N(\delta)=153$, so
\[
 f(30)=901-60\sqrt{457}+62\sqrt{33}+\sqrt{153}<-5.
\]
Use $\sqrt{457}>213/10$, $\sqrt{33}<29/5$ and
$\sqrt{153}<62/5$. Convexity gives $T>30$.
Write $t+1=(u+v\sqrt{-2})^2$ with $u,v\in\mathbb Z$.
If $30<T<32$, then
\[
 29<|t+1|=u^2+2v^2<33.
\]
The values $30$ and $31$ are impossible modulo $16$.
For $u^2+2v^2=32$, reduction modulo $16$ forces $4\mid u,v$,
and then $u=0$, $v=\pm4$. This gives $t=-33$, contrary to $T<32$.
Thus $N(t)<N_0$ is impossible in the last two rows as well.

\begin{longtable}{r|r r|r|r|r}
\caption{All remaining triples and the common cutoff procedure. Pair indices refer to Table~\ref{tab:pr-forty-pairs}. Coordinates mean $u+v\omega$.}\label{tab:fixed56}\\
Pair&$u$&$v$&$N(s)$&$M_{bs}$&$N_0$\\\hline
\endfirsthead
Pair&$u$&$v$&$N(s)$&$M_{bs}$&$N_0$\\\hline
\endhead
0 & -4 & -2 & 20 & 285 & 1024 \\
0 & 4 & -2 & 20 & 285 & 1024 \\
0 & 0 & 5 & 25 & 1368 & 1934 \\
1 & 0 & 10 & 100 & 42 & 2311 \\
2 & -4 & 4 & 32 & 2238 & 4770 \\
2 & 8 & 1 & 65 & 72 & 2359 \\
3 & 6 & 9 & 117 & 44 & 3896 \\
4 & -4 & -4 & 32 & 396 & 3299 \\
4 & -4 & 4 & 32 & 396 & 3299 \\
5 & -1 & 8 & 65 & 172 & 4857 \\
6 & 7 & 6 & 85 & 29 & 2018 \\
7 & -2 & -2 & 12 & 474 & 1024 \\
7 & -2 & 2 & 12 & 474 & 1024 \\
8 & 1 & 4 & 33 & 133 & 1024 \\
9 & 6 & 6 & 108 & 56 & 4775 \\
10 & -13 & 13 & 169 & 77 & 13059 \\
11 & -8 & 16 & 192 & 60 & 12244 \\
12 & -9 & 6 & 63 & 74 & 2301 \\
13 & -10 & 1 & 91 & 180 & 9744 \\
14 & -6 & 9 & 63 & 35 & 1195 \\
15 & -1 & 10 & 91 & 37 & 1760 \\
16 & -4 & -4 & 48 & 430 & 7354 \\
16 & 8 & -8 & 64 & 271 & 7441 \\
17 & 7 & -2 & 39 & 78 & 1024 \\
18 & 6 & 9 & 171 & 48 & 7943 \\
19 & 11 & 1 & 133 & 48 & 6361 \\
20 & -6 & -3 & 63 & 473 & 14404 \\
20 & -10 & 9 & 91 & 272 & 14608 \\
21 & -4 & 8 & 48 & 155 & 2461 \\
22 & -5 & 4 & 21 & 422 & 1622 \\
22 & -1 & -4 & 21 & 422 & 1622 \\
23 & 4 & -8 & 48 & 4138 & 21270 \\
23 & -1 & 12 & 133 & 78 & 8968 \\
24 & 7 & 7 & 147 & 41 & 5928 \\
25 & -5 & -2 & 45 & 381 & 5913 \\
25 & -5 & 2 & 45 & 381 & 5913 \\
26 & -2 & 4 & 84 & 186 & 8435 \\
27 & -6 & -2 & 60 & 370 & 9706 \\
27 & -6 & 2 & 60 & 370 & 9706 \\
28 & -3 & 4 & 105 & 197 & 13453 \\
29 & 1 & 5 & 56 & 67 & 1517 \\
30 & 6 & 7 & 176 & 57 & 11143 \\
31 & -3 & -1 & 15 & 646 & 1451 \\
31 & -3 & 3 & 27 & 261 & 1460 \\
32 & 0 & 5 & 75 & 94 & 3590 \\
33 & -4 & -1 & 24 & 579 & 3059 \\
33 & -4 & 3 & 40 & 263 & 3090 \\
34 & -1 & 5 & 96 & 115 & 6843 \\
35 & -5 & -1 & 35 & 532 & 5592 \\
35 & -5 & 3 & 55 & 266 & 5661 \\
36 & -2 & 5 & 119 & 133 & 11532 \\
37 & -6 & -1 & 48 & 499 & 9289 \\
37 & -6 & 3 & 72 & 269 & 9415 \\
38 & -4 & 0 & 16 & 540 & 1281 \\
38 & 4 & -4 & 32 & 182 & 2561 \\
39 & 8 & 4 & 80 & 38 & 6401 \\
\end{longtable}
\begin{proposition}\label{prop:one-small}
A $D(1)$-tuple with exactly one entry of modulus less than three has
at most five elements.
\end{proposition}
\begin{proof}
The pair $(1,3)$ has just been excluded. Every other sextuple would
yield one of the forty pairs in Table~\ref{tab:pr-forty-pairs}, then
one of the fifty-six triples in Table~\ref{tab:fixed56}. The common
analytic argument excludes $N(t)\ge N_0$, and the preceding argument
excludes $N(t)<N_0$. This contradiction proves the assertion.
\end{proof}

\section{Proof of the uniform bound}

\begin{proof}[Proof of Theorem~\ref{thm:main}]
If a tuple had at least six elements, choose six and order them by
modulus. If none has modulus less than three, Proposition~\ref{prop:no-all-large-six}
gives a contradiction. If exactly one has modulus less than three,
Proposition~\ref{prop:one-small} gives a contradiction. If at least
two have modulus less than three, Proposition~\ref{sm:main} gives a
contradiction. These cases exhaust all sextuples, including ties in
modulus, and prove the assertion.
\end{proof}

\par\medskip
\noindent\textbf{Funding.} The authors were supported by the Croatian Science Foundation under the project no. IP-2022-10-5008 (TEBAG). The authors acknowledge support from the project ``Implementation of cutting-edge research and its application as part of the Scientific Center of Excellence for Quantum and Complex Systems, and Representations of Lie Algebras'', Grant No. PK.1.1.10.0004, co-financed by the European Union through the European Regional Development Fund -- Competitiveness and Cohesion Programme 2021--2027. This research was funded by the European Union's NextGenerationEU through the National Recovery and Resilience Plan 2021--2026 Institutional grant of University of Zagreb Faculty of Science (IK IA 1.1.3. Impact4Math).

\medskip
\noindent\textbf{Use of artificial intelligence.} The authors used ChatGPT 6 (OpenAI) and Claude Opus 5.5 (Anthropic) to assist with developing proofs, constructing examples and counterexamples, searching the literature, writing code for exact computations, and drafting and revising the manuscript. The authors take full responsibility for the content of this paper.

\begin{thebibliography}{99}
\bibitem{A} N.~Ad\v zaga, On the size of Diophantine $m$-tuples in
imaginary quadratic number rings, \emph{Bull. Math. Sci.} \textbf{11}
(2021), 1950020.
\href{https://doi.org/10.1142/S1664360719500206}{doi:10.1142/S1664360719500206}.

\bibitem{Bennett96} M.~A. Bennett, Simultaneous rational approximation
to binomial functions, \emph{Trans. Amer. Math. Soc.} \textbf{348}
(1996), no.~5, 1717--1738.
\href{https://doi.org/10.1090/S0002-9947-96-01480-8}{doi:10.1090/S0002-9947-96-01480-8}.

\bibitem{CF} M.~Cipu and Y.~Fujita, On the length of $D(\pm1)$-tuples
in imaginary quadratic rings, \emph{Bull. Lond. Math. Soc.}
\textbf{56} (2024), no.~1, 274--287.
\href{https://doi.org/10.1112/blms.12929}{doi:10.1112/blms.12929}.

\bibitem{Cipu25} M.~Cipu, \emph{$D(\pm1)$-quintuples in imaginary
quadratic fields}, lecture slides, Representation Theory XIX,
Dubrovnik, 24 June 2025. Joint work in progress with Y.~Fujita.
\href{https://web.math.pmf.unizg.hr/~congress/Dubrovnik25/Slides-NT-25/Cipu.pdf}{Lecture slides}.

\bibitem{Dujella24} A.~Dujella, \emph{Diophantine $m$-tuples and Elliptic
Curves}, Developments in Mathematics \textbf{79}, Springer, Cham, 2024.
\href{https://doi.org/10.1007/978-3-031-56724-7}{doi:10.1007/978-3-031-56724-7}.

\bibitem{DujellaOpen} A.~Dujella, \emph{Open problems on Diophantine
$m$-tuples and elliptic curves}, version of 15 October 2026.
\url{https://web.math.pmf.unizg.hr/~duje/pdf/open2.pdf}
(accessed 15 October 2026).

\bibitem{F10} Z.~Franu\v{s}i\'c, On the extension of the Diophantine
pair $\{1,3\}$ in $\mathbb Z[\sqrt d]$, \emph{J. Integer Seq.}
\textbf{13} (2010), Article 10.9.6.
\href{https://cs.uwaterloo.ca/journals/JIS/VOL13/Franusic/franusic4.html}{Journal version}.

\bibitem{FK11} Z.~Franu\v{s}i\'c and D.~Kreso, Nonextensibility
of the pair $\{1,3\}$ to a Diophantine quintuple in
$\mathbb Z[\sqrt{-2}]$, \emph{J. Comb. Number Theory}
\textbf{3} (2011), no.~3.
\href{https://web.math.pmf.unizg.hr/~fran/clanci/dz-web.pdf}{Author's version}.
\href{https://doi.org/10.48550/arXiv.1312.5241}{Preprint doi:10.48550/arXiv.1312.5241}.

\bibitem{JZ} B.~Jadrijevi\'c and V.~Ziegler, A system of relative
Pellian equations and a related family of relative Thue equations,
\emph{Int. J. Number Theory} \textbf{2} (2006), no.~4, 569--590.
\href{https://doi.org/10.1142/S1793042106000735}{doi:10.1142/S1793042106000735}.

\bibitem{W} N.~Wildner Fraga, \emph{Uniform bounds for $D(1)$-tuples in
rings of integers}, electronic preprint, 16 September 2026,
version at commit \texttt{768dda32edc7}.
\href{https://github.com/natanaelwf/Solving-Open-Problems/blob/768dda32edc7166f9dd094b419df250b1e008531/dujella-problems-1.12-1.13/README.md}{Pinned preprint}.
\end{thebibliography}
\bigskip
\noindent Zrinka Franu\v{s}i\'{c} and Tomislav Pejkovi\'{c}\\
Department of Mathematics, Faculty of Science\\
University of Zagreb\\
Bijenička cesta 30, 10000 Zagreb, Croatia. \texttt{fran@math.hr}, \texttt{pejkovic@math.hr}
\end{document}
